\documentclass[11pt,a4paper]{book}

% ==========================================
% PAQUETES MATEMÁTICOS Y DE FORMATO
% ==========================================
\usepackage[utf8]{inputenc}
\usepackage[T1]{fontenc}
\usepackage[spanish,es-nodecimaldot]{babel}
\usepackage{amsmath,amssymb,amsfonts,amsthm}
\usepackage{mathrsfs}
\usepackage{mathtools}
\usepackage{physics}
\usepackage{bm}
\usepackage{siunitx}
\usepackage{geometry}
\geometry{left=2.5cm,right=2.5cm,top=2.5cm,bottom=2.5cm}
\usepackage{fancyhdr}
\usepackage{hyperref}
\usepackage{xcolor}
\usepackage{tcolorbox}
\tcbuselibrary{theorems,breakable,skins}
\usepackage{enumitem}

% Configuración de hipervínculos
\hypersetup{
    colorlinks=true,
    linkcolor=blue!70!black,
    urlcolor=blue!70!black,
    citecolor=blue!70!black,
    pdftitle={Solucionario Detallado: Electrodinámica (Griffiths) - Capítulos 4 al 12},
    pdfauthor={Física Teórica}
}

% Configuración de cabeceras
\pagestyle{fancy}
\fancyhf{}
\fancyhead[LE,RO]{\thepage}
\fancyhead[RE]{\leftmark}
\fancyhead[LO]{\rightmark}
\setlength{\headheight}{28pt}
\renewcommand{\headrulewidth}{0.4pt}

% ==========================================
% ENTORNOS PERSONALIZADOS PARA EJERCICIOS
% ==========================================
\newtcbtheorem[number within=chapter]{ejercicio}{Problema}{
    enhanced,
    breakable,
    colback=blue!3!white,
    colframe=blue!65!black,
    fonttitle=\bfseries\raggedright,
    separator sign={\ ---},
    attach boxed title to top left={yshift=-2mm,xshift=4mm},
    minipage boxed title,
    boxed title style={
        colback=blue!65!black,
        boxsep=2mm
    }
}{prob}

\newtheorem{definicion}{Definición}[chapter]
\newtheorem{ejemplo}{Ejemplo}[chapter]
\newtheorem{propiedad}{Propiedad}[chapter]
\newtheorem{teorema}{Teorema}[chapter]

\newenvironment{solucion}{%
    \par\medskip\noindent\textbf{\color{blue!65!black}\Large $\blacktriangleright$ Solución paso a paso:}\par\medskip
}{%
    \hfill$\blacksquare$\par\bigskip
}

% Comandos personalizados para vectores unitarios al estilo Griffiths
\newcommand{\rhat}{\hat{\bm{r}}}
\newcommand{\thetahat}{\hat{\bm{\theta}}}
\newcommand{\phihat}{\hat{\bm{\phi}}}
\newcommand{\shat}{\hat{\bm{s}}}
\newcommand{\xhat}{\hat{\bm{x}}}
\newcommand{\yhat}{\hat{\bm{y}}}
\newcommand{\zhat}{\hat{\bm{z}}}
\newcommand{\nhat}{\hat{\bm{n}}}
\newcommand{\Evec}{\bm{E}}
\newcommand{\Dvec}{\bm{D}}
\newcommand{\Pvec}{\bm{P}}
\newcommand{\Bvec}{\bm{B}}
\newcommand{\Hvec}{\bm{H}}
\newcommand{\Mvec}{\bm{M}}
\newcommand{\Avec}{\bm{A}}
\newcommand{\Jvec}{\bm{J}}
\newcommand{\Svec}{\bm{S}}
\newcommand{\rsep}{\boldsymbol{\mathscr{r}}}

\title{\Huge \textbf{Solucionario Avanzado y Detallado}\\
\LARGE Basado en \textit{Introduction to Electrodynamics} (D. J. Griffiths)\\
\large Desde Electrostática en Medios Materiales hasta Electrodinámica Relativista}
\author{}
\date{}

\begin{document}

\maketitle
\begingroup
\let\cleardoublepage\clearpage
\tableofcontents
\endgroup
\newpage

% =========================================================================
% CAPÍTULO 4: CAMPOS ELÉCTRICOS EN LA MATERIA
% =========================================================================
\chapter[Campos eléctricos en la materia]{Campos Eléctricos en la Materia (Polarización y Dieléctricos)}

\section{Polarización, Cargas Ligadas y Campo Eléctrico Macroscópico}

\begin{ejercicio}{Esfera con polarización congelada uniforme}{esfera_polarizada}
Una esfera de radio $R$ posee una polarización uniforme $\Pvec(\bm{r}) = P_0 \zhat$ en todo su volumen (donde $P_0$ es una constante). 
\begin{enumerate}[label=\alph*)]
    \item Calcule las densidades de carga ligada de volumen $\rho_b$ y de superficie $\sigma_b$.
    \item Determine el potencial eléctrico $V(\bm{r})$ y el campo eléctrico $\Evec(\bm{r})$ tanto en el interior ($r < R$) como en el exterior ($r > R$) de la esfera, integrando directamente el potencial generado por las cargas ligadas.
\end{enumerate}
\end{ejercicio}

\begin{solucion}
\textbf{Parte a) Cálculo de las densidades de carga ligada.}

Por definición, una distribución de polarización $\Pvec(\bm{r})$ dentro de un material equivale electrostáticamente a una densidad volumétrica de carga ligada $\rho_b$ y una densidad superficial de carga ligada $\sigma_b$ dadas por las expresiones:
\begin{equation}
    \rho_b = -\div{\Pvec}, \qquad \sigma_b = \Pvec \cdot \nhat
\end{equation}
donde $\nhat$ es el vector unitario normal saliente a la superficie del material.

Primero calculamos la divergencia de $\Pvec$ en coordenadas cartesianas, dado que $\Pvec = P_0 \zhat$ tiene únicamente componente en el eje $z$ y $P_0$ es constante en el espacio:
\begin{equation}
    \rho_b = -\left( \pdv{P_x}{x} + \pdv{P_y}{y} + \pdv{P_z}{z} \right) = -\left( 0 + 0 + \pdv{P_0}{z} \right) = 0
\end{equation}
Por lo tanto, no hay carga ligada en el interior del volumen de la esfera ($\rho_b = 0$).

Para la superficie de la esfera ($r = R$), el vector normal unitario saliente en coordenadas esféricas es $\nhat = \rhat$. Recordando la relación geométrica entre el vector unitario cartesiano $\zhat$ y el vector radial esférico $\rhat$, donde el ángulo polar $\theta$ se mide desde el eje $z$ positivo:
\begin{equation}
    \zhat \cdot \rhat = \cos\theta
\end{equation}
Sustituimos esto en la definición de $\sigma_b$:
\begin{equation}
    \sigma_b = \Pvec \cdot \rhat = (P_0 \zhat) \cdot \rhat = P_0 \cos\theta
\end{equation}

\medskip
\textbf{Parte b) Cálculo del potencial $V(\bm{r})$ y del campo eléctrico $\Evec(\bm{r})$.}

Como $\rho_b = 0$ y no existen cargas libres ($\rho_f = 0$), el problema se reduce exactamente al de una cáscara esférica de radio $R$ con una densidad superficial de carga $\sigma(\theta) = P_0 \cos\theta$ en el vacío.

Para una cáscara esférica de radio $R$ con distribución azimutalmente simétrica $\sigma(\theta')$, el potencial electrostático general solución de la ecuación de Laplace $\nabla^2 V = 0$ viene dado por el desarrollo en polinomios de Legendre:
\begin{equation}
    V(r,\theta) = \begin{cases} 
    \displaystyle \sum_{l=0}^{\infty} A_l r^l P_l(\cos\theta), & \text{para } r \le R \\[12pt]
    \displaystyle \sum_{l=0}^{\infty} \frac{B_l}{r^{l+1}} P_l(\cos\theta), & \text{para } r \ge R
    \end{cases}
\end{equation}

Imponemos las dos condiciones de contorno electrostáticas en $r = R$:
\begin{enumerate}
    \item \textbf{Continuidad del potencial en $r = R$:}
    \begin{equation}
        V_{\text{int}}(R,\theta) = V_{\text{ext}}(R,\theta) \implies A_l R^l = \frac{B_l}{R^{l+1}} \implies B_l = A_l R^{2l+1}
    \end{equation}
    
    \item \textbf{Discontinuidad de la derivada normal del potencial en $r = R$:}
    Sabemos que el salto en la componente normal del campo eléctrico está gobernado por la carga superficial total (aquí solo carga ligada $\sigma_b$):
    \begin{equation}
        \left. \pdv{V_{\text{ext}}}{r} \right|_{r=R} - \left. \pdv{V_{\text{int}}}{r} \right|_{r=R} = -\frac{\sigma_b(\theta)}{\varepsilon_0}
    \end{equation}
\end{enumerate}

Derivamos las series respecto a $r$:
\begin{align}
    \pdv{V_{\text{ext}}}{r} &= \sum_{l=0}^{\infty} -(l+1)\frac{B_l}{r^{l+2}} P_l(\cos\theta) \\
    \pdv{V_{\text{int}}}{r} &= \sum_{l=0}^{\infty} l A_l r^{l-1} P_l(\cos\theta)
\end{align}
Evaluando en $r = R$ y usando $B_l = A_l R^{2l+1}$:
\begin{align}
    \sum_{l=0}^{\infty} \left[ -(l+1)\frac{A_l R^{2l+1}}{R^{l+2}} - l A_l R^{l-1} \right] P_l(\cos\theta) &= -\frac{P_0 \cos\theta}{\varepsilon_0} \\
    \sum_{l=0}^{\infty} \left[ -(l+1) - l \right] A_l R^{l-1} P_l(\cos\theta) &= -\frac{P_0 \cos\theta}{\varepsilon_0} \\
    \sum_{l=0}^{\infty} (2l+1) A_l R^{l-1} P_l(\cos\theta) &= \frac{P_0}{\varepsilon_0} P_1(\cos\theta)
\end{align}
donde hemos identificado que $\cos\theta = P_1(\cos\theta)$ es el polinomio de Legendre de orden $l=1$. Por la ortogonalidad de los polinomios de Legendre en el intervalo $[-1,1]$, todos los coeficientes $A_l$ son nulos excepto para $l = 1$:
\begin{equation}
    A_l = 0 \quad (\forall l \neq 1)
\end{equation}
Para $l = 1$, igualamos coeficientes:
\begin{equation}
    (2(1)+1) A_1 R^{1-1} = \frac{P_0}{\varepsilon_0} \implies 3 A_1 = \frac{P_0}{\varepsilon_0} \implies A_1 = \frac{P_0}{3\varepsilon_0}
\end{equation}
Con $A_1$ conocido, determinamos $B_1$:
\begin{equation}
    B_1 = A_1 R^{2(1)+1} = \frac{P_0 R^3}{3\varepsilon_0}
\end{equation}

Sustituyendo los coeficientes en la expresión del potencial:
\begin{itemize}
    \item \textbf{Para el interior ($r \le R$):}
    \begin{equation}
        V_{\text{int}}(r,\theta) = A_1 r P_1(\cos\theta) = \frac{P_0}{3\varepsilon_0} r \cos\theta
    \end{equation}
    Como en coordenadas cartesianas $z = r \cos\theta$, esto se escribe como:
    \begin{equation}
        V_{\text{int}}(x,y,z) = \frac{P_0}{3\varepsilon_0} z
    \end{equation}
    El campo eléctrico interno se halla aplicando el operador $-\grad V$:
    \begin{equation}
        \Evec_{\text{int}} = -\grad V_{\text{int}} = -\pdv{V_{\text{int}}}{z}\zhat = -\frac{P_0}{3\varepsilon_0}\zhat = -\frac{1}{3\varepsilon_0}\Pvec
    \end{equation}
    Observamos el resultado notable de que el campo eléctrico en el interior de una esfera uniformemente polarizada es \textbf{completamente uniforme} y opuesto al vector polarización.

    \item \textbf{Para el exterior ($r \ge R$):}
    \begin{equation}
        V_{\text{ext}}(r,\theta) = \frac{B_1}{r^2} P_1(\cos\theta) = \frac{P_0 R^3}{3\varepsilon_0} \frac{\cos\theta}{r^2}
    \end{equation}
    Si definimos el momento dipolar total de la esfera como el producto de su volumen por su polarización:
    \begin{equation}
        \bm{p} = \left(\frac{4}{3}\pi R^3\right) \Pvec = \frac{4}{3}\pi R^3 P_0 \zhat
    \end{equation}
    entonces podemos reescribir $\frac{P_0 R^3}{3\varepsilon_0} = \frac{p}{4\pi\varepsilon_0}$, quedando el potencial exterior como:
    \begin{equation}
        V_{\text{ext}}(r,\theta) = \frac{1}{4\pi\varepsilon_0}\frac{p \cos\theta}{r^2} = \frac{1}{4\pi\varepsilon_0}\frac{\bm{p}\cdot\rhat}{r^2}
    \end{equation}
    El campo eléctrico en el exterior es exactamente el de un dipolo puntual ubicado en el origen:
    \begin{align}
        \Evec_{\text{ext}}(r,\theta) &= -\grad V_{\text{ext}} = -\left( \pdv{V_{\text{ext}}}{r}\rhat + \frac{1}{r}\pdv{V_{\text{ext}}}{\theta}\thetahat \right) \nonumber \\
        &= -\left( -\frac{2p\cos\theta}{4\pi\varepsilon_0 r^3}\rhat - \frac{p\sin\theta}{4\pi\varepsilon_0 r^3}\thetahat \right) \nonumber \\
        &= \frac{1}{4\pi\varepsilon_0}\frac{p}{r^3}\left(2\cos\theta\,\rhat + \sin\theta\,\thetahat\right) = \frac{P_0 R^3}{3\varepsilon_0 r^3}\left(2\cos\theta\,\rhat + \sin\theta\,\thetahat\right)
    \end{align}
\end{itemize}
\end{solucion}

% -------------------------------------------------------------------------
\section[Desplazamiento eléctrico y ley de Gauss]{El Desplazamiento Eléctrico \texorpdfstring{$\Dvec$}{D} y la Ley de Gauss en Dieléctricos}

\begin{ejercicio}{Esfera metálica rodeada de una capa dieléctrica no homogénea}{esfera_dielectrica}
Una esfera conductora de radio $a$ porta una carga libre neta $Q$. Está rodeada por una cáscara dieléctrica lineal concéntrica de radio interno $a$ y radio externo $b$, cuya susceptibilidad eléctrica depende de la distancia radial según $\chi_e(r) = \frac{k}{r^2}$, donde $k$ es una constante positiva con dimensiones de área.
\begin{enumerate}[label=\alph*)]
    \item Encuentre los vectores $\Dvec(\bm{r})$, $\Evec(\bm{r})$ y $\Pvec(\bm{r})$ en todas las regiones del espacio.
    \item Calcule la densidad volumétrica de carga ligada $\rho_b$ y las densidades superficiales de carga ligada $\sigma_b$ en $r=a$ y $r=b$. Verifique explícitamente que la carga ligada total del sistema es exactamente cero.
    \item Determine la energía electrostática total almacenada en el sistema.
\end{enumerate}
\end{ejercicio}

\begin{solucion}
\textbf{Parte a) Cálculo de $\Dvec$, $\Evec$ y $\Pvec$ en todo el espacio.}

Dada la simetría esférica del conductor y el hecho de que $\chi_e(r)$ depende únicamente de la coordenada radial $r$, los campos $\Dvec$, $\Evec$ y $\Pvec$ apuntan puramente en la dirección radial $\rhat$ y sus magnitudes dependen solo de $r$.

Aplicamos la ley de Gauss para el desplazamiento eléctrico $\Dvec$ sobre una superficie gaussiana esférica de radio $r$:
\begin{equation}
    \oint_{\mathcal{S}} \Dvec \cdot d\bm{a} = Q_{f,\text{enc}}
\end{equation}
Analizamos por regiones:
\begin{enumerate}[label=\roman*)]
    \item \textbf{Región $r < a$ (Interior del conductor):}
    En equilibrio electrostático, el campo eléctrico dentro de un conductor perfecto es nulo y toda la carga libre reside en la superficie $r=a$. Por tanto:
    \begin{equation}
        \Evec(\bm{r}) = \bm{0}, \qquad \Pvec(\bm{r}) = \bm{0}, \qquad \Dvec(\bm{r}) = \bm{0} \qquad (r < a)
    \end{equation}

    \item \textbf{Región $a < r < b$ (Dentro del dieléctrico):}
    La carga libre encerrada por una esfera gaussiana de radio $r \in (a,b)$ es $Q_{f,\text{enc}} = Q$. Por simetría esférica:
    \begin{equation}
        D(r) (4\pi r^2) = Q \implies \Dvec(\bm{r}) = \frac{Q}{4\pi r^2}\rhat
    \end{equation}
    En un medio dieléctrico lineal, la relación constitutiva entre $\Dvec$ y $\Evec$ es $\Dvec = \varepsilon\Evec = \varepsilon_0(1 + \chi_e(r))\Evec$. Despejando $\Evec(\bm{r})$:
    \begin{equation}
        \Evec(\bm{r}) = \frac{\Dvec(\bm{r})}{\varepsilon_0\big(1 + \chi_e(r)\big)} = \frac{Q}{4\pi\varepsilon_0 r^2 \left(1 + \frac{k}{r^2}\right)}\rhat = \frac{Q}{4\pi\varepsilon_0 (r^2 + k)}\rhat
    \end{equation}
    La polarización $\Pvec$ se obtiene mediante $\Pvec = \varepsilon_0 \chi_e(r) \Evec$ (o equivalentemente $\Pvec = \Dvec - \varepsilon_0\Evec$):
    \begin{equation}
        \Pvec(\bm{r}) = \varepsilon_0 \left(\frac{k}{r^2}\right) \frac{Q}{4\pi\varepsilon_0 (r^2 + k)}\rhat = \frac{Q k}{4\pi r^2 (r^2 + k)}\rhat
    \end{equation}

    \item \textbf{Región $r > b$ (Vacío exterior):}
    Para $r > b$, la carga libre encerrada sigue siendo $Q$, por lo que:
    \begin{equation}
        \Dvec(\bm{r}) = \frac{Q}{4\pi r^2}\rhat
    \end{equation}
    Al estar en el vacío ($\chi_e = 0$):
    \begin{equation}
        \Pvec(\bm{r}) = \bm{0}, \qquad \Evec(\bm{r}) = \frac{\Dvec(\bm{r})}{\varepsilon_0} = \frac{Q}{4\pi\varepsilon_0 r^2}\rhat
    \end{equation}
\end{enumerate}

\medskip
\textbf{Parte b) Cálculo de cargas ligadas y comprobación de neutralidad.}

Calculamos primero la densidad volumétrica de carga ligada $\rho_b = -\div{\Pvec}$ en la región $a < r < b$. Usando la expresión de la divergencia en coordenadas esféricas para un campo puramente radial $\Pvec = P_r(r)\rhat$:
\begin{equation}
    \rho_b(r) = -\frac{1}{r^2}\dv{}{r}\left(r^2 P_r(r)\right)
\end{equation}
Observamos el término dentro de la derivada:
\begin{equation}
    r^2 P_r(r) = r^2 \left[\frac{Q k}{4\pi r^2 (r^2 + k)}\right] = \frac{Q k}{4\pi (r^2 + k)}
\end{equation}
Derivamos respecto a $r$ aplicando la regla de la cadena:
\begin{equation}
    \dv{}{r}\left(\frac{Q k}{4\pi (r^2 + k)}\right) = \frac{Q k}{4\pi} \left[ -\frac{2r}{(r^2 + k)^2} \right] = -\frac{Q k r}{2\pi (r^2 + k)^2}
\end{equation}
Por lo tanto, sustituyendo en $\rho_b(r)$:
\begin{equation}
    \rho_b(r) = -\frac{1}{r^2}\left(-\frac{Q k r}{2\pi (r^2 + k)^2}\right) = +\frac{Q k}{2\pi r (r^2 + k)^2}
\end{equation}
Nótese que debido a la inhomogeneidad espacial de $\chi_e(r)$, aparece una densidad de carga ligada en el volumen de signo positivo (mismo signo que $Q$).

Ahora calculamos las densidades superficiales de carga ligada $\sigma_b = \Pvec \cdot \nhat$ en las dos fronteras del dieléctrico:
\begin{itemize}
    \item \textbf{En la superficie interna $r = a$:} El vector unitario normal saliente \textit{respecto al dieléctrico} apunta hacia el centro de la esfera, es decir, $\nhat = -\rhat$. Evaluando $\Pvec(a)$:
    \begin{equation}
        \sigma_b(a) = \Pvec(a) \cdot (-\rhat) = -\frac{Q k}{4\pi a^2 (a^2 + k)}
    \end{equation}
    \item \textbf{En la superficie externa $r = b$:} El vector unitario normal saliente respecto al dieléctrico apunta hacia afuera, es decir, $\nhat = +\rhat$. Evaluando $\Pvec(b)$:
    \begin{equation}
        \sigma_b(b) = \Pvec(b) \cdot (+\rhat) = +\frac{Q k}{4\pi b^2 (b^2 + k)}
    \end{equation}
\end{itemize}

Procedemos a calcular la carga ligada total $Q_{b,\text{total}}$ sumando las contribuciones de las dos superficies y del volumen:
\begin{align}
    Q_{b,\text{sup}}(a) &= (4\pi a^2) \sigma_b(a) = (4\pi a^2)\left(-\frac{Q k}{4\pi a^2 (a^2 + k)}\right) = -\frac{Q k}{a^2 + k} \\
    Q_{b,\text{sup}}(b) &= (4\pi b^2) \sigma_b(b) = (4\pi b^2)\left(+\frac{Q k}{4\pi b^2 (b^2 + k)}\right) = +\frac{Q k}{b^2 + k}
\end{align}
Para la carga ligada en el volumen, integramos $\rho_b(r)$ en cascarones esféricos $d\tau = 4\pi r^2 dr$ desde $r=a$ hasta $r=b$:
\begin{equation}
    Q_{b,\text{vol}} = \int_{a}^{b} \rho_b(r) 4\pi r^2 dr = \int_{a}^{b} \left(\frac{Q k}{2\pi r (r^2 + k)^2}\right) 4\pi r^2 dr = Q k \int_{a}^{b} \frac{2r}{(r^2 + k)^2} dr
\end{equation}
Haciendo el cambio de variable elemental $u = r^2 + k \implies du = 2r\,dr$:
\begin{equation}
    Q_{b,\text{vol}} = Q k \int_{a^2+k}^{b^2+k} \frac{du}{u^2} = Q k \left[ -\frac{1}{u} \right]_{a^2+k}^{b^2+k} = Q k \left( \frac{1}{a^2 + k} - \frac{1}{b^2 + k} \right) = \frac{Q k}{a^2 + k} - \frac{Q k}{b^2 + k}
\end{equation}
Sumando las tres contribuciones:
\begin{equation}
    Q_{b,\text{total}} = Q_{b,\text{sup}}(a) + Q_{b,\text{sup}}(b) + Q_{b,\text{vol}} = -\frac{Q k}{a^2 + k} + \frac{Q k}{b^2 + k} + \frac{Q k}{a^2 + k} - \frac{Q k}{b^2 + k} = 0
\end{equation}
Queda demostrado rigurosamente que la carga ligada neta es idénticamente nula, como exige la conservación de la carga en un dieléctrico inicialmente neutro.

\medskip
\textbf{Parte c) Energía electrostática total del sistema.}

En presencia de medios dieléctricos lineales, la energía total almacenada en el campo electrostático viene dada por la integral en todo el espacio:
\begin{equation}
    W = \frac{1}{2} \int_{\text{todo el espacio}} \Dvec \cdot \Evec \, d\tau
\end{equation}
Dividimos la integral en las regiones donde los campos no son nulos ($a < r < b$ y $r > b$), usando $d\tau = 4\pi r^2 dr$:
\begin{equation}
    W = \frac{1}{2} \int_{a}^{b} \left(\frac{Q}{4\pi r^2}\right)\left(\frac{Q}{4\pi\varepsilon_0(r^2+k)}\right) 4\pi r^2 dr + \frac{1}{2} \int_{b}^{\infty} \left(\frac{Q}{4\pi r^2}\right)\left(\frac{Q}{4\pi\varepsilon_0 r^2}\right) 4\pi r^2 dr
\end{equation}
Simplificamos las constantes en ambas integrales:
\begin{equation}
    W = \frac{Q^2}{8\pi\varepsilon_0} \left[ \int_{a}^{b} \frac{dr}{r^2 + k} + \int_{b}^{\infty} \frac{dr}{r^2} \right]
\end{equation}
Resolvemos cada integral por separado. Para la primera, recordamos que $\int \frac{dr}{r^2 + (\sqrt{k})^2} = \frac{1}{\sqrt{k}}\arctan\left(\frac{r}{\sqrt{k}}\right)$:
\begin{equation}
    I_1 = \int_{a}^{b} \frac{dr}{r^2 + k} = \frac{1}{\sqrt{k}}\left[ \arctan\left(\frac{b}{\sqrt{k}}\right) - \arctan\left(\frac{a}{\sqrt{k}}\right) \right]
\end{equation}
Para la segunda integral:
\begin{equation}
    I_2 = \int_{b}^{\infty} \frac{dr}{r^2} = \left[ -\frac{1}{r} \right]_{b}^{\infty} = \frac{1}{b}
\end{equation}
Agrupando los resultados, obtenemos la energía total:
\begin{equation}
    W = \frac{Q^2}{8\pi\varepsilon_0} \left\{ \frac{1}{b} + \frac{1}{\sqrt{k}}\left[ \arctan\left(\frac{b}{\sqrt{k}}\right) - \arctan\left(\frac{a}{\sqrt{k}}\right) \right] \right\}
\end{equation}
\end{solucion}

% -------------------------------------------------------------------------
\section[Problemas de frontera con dieléctricos]{Problemas de Valor en la Frontera con Dieléctricos y Fuerzas}

\begin{ejercicio}{Esfera dieléctrica en un campo eléctrico externo uniforme}{esfera_campo_uniforme}
Una esfera homogénea, isótropa y lineal de radio $R$ y permitividad eléctrica $\varepsilon$ (constante dieléctrica $\varepsilon_r = \varepsilon/\varepsilon_0$) se coloca en una región donde existía previamente un campo eléctrico uniforme $\Evec_0 = E_0 \zhat$. No hay cargas libres ni dentro ni sobre la superficie de la esfera. Determine el potencial $V(r,\theta)$ y el campo eléctrico resultante $\Evec(\bm{r})$ tanto dentro como fuera de la esfera, utilizando el método de separación de variables y las condiciones de contorno para dieléctricos.
\end{ejercicio}

\begin{solucion}
Dado que no hay carga libre en ninguna parte del espacio ($\rho_f = 0$), dentro de la esfera ($r < R$) la ley de Gauss es $\div{\Dvec} = \div(\varepsilon\Evec) = -\varepsilon\nabla^2 V = 0$, y fuera de la esfera ($r > R$) es $-\varepsilon_0\nabla^2 V = 0$. En consecuencia, el potencial electrostático $V(r,\theta)$ satisface la ecuación de Laplace $\nabla^2 V = 0$ tanto en $r < R$ como en $r > R$.

Las condiciones de frontera que debe cumplir el potencial son:
\begin{enumerate}
    \item \textbf{Continuidad del potencial en la interfaz $r = R$:}
    \begin{equation}
        V_{\text{int}}(R,\theta) = V_{\text{ext}}(R,\theta)
    \end{equation}
    \item \textbf{Continuidad de la componente normal de $\Dvec$ en $r = R$ (ya que $\sigma_f = 0$):}
    \begin{equation}
        D_{\text{int},\perp}(R,\theta) = D_{\text{ext},\perp}(R,\theta) \implies -\left.\varepsilon \pdv{V_{\text{int}}}{r}\right|_{r=R} = -\left.\varepsilon_0 \pdv{V_{\text{ext}}}{r}\right|_{r=R} \implies \left.\varepsilon_r \pdv{V_{\text{int}}}{r}\right|_{r=R} = \left.\pdv{V_{\text{ext}}}{r}\right|_{r=R}
    \end{equation}
    \item \textbf{Comportamiento asintótico en el infinito ($r \to \infty$):}
    Muy lejos de la esfera, el campo debe recuperar el valor uniforme $\Evec_0 = E_0\zhat$. Como $\Evec_0 = -\grad V_0$, el potencial asintótico es:
    \begin{equation}
        V(r,\theta) \xrightarrow{r \to \infty} -E_0 z = -E_0 r \cos\theta = -E_0 r P_1(\cos\theta)
    \end{equation}
    \item \textbf{Regularidad en el origen ($r = 0$):}
    El potencial $V_{\text{int}}$ no puede divergir en $r=0$.
\end{enumerate}

Escribimos las soluciones generales con simetría azimutal imponiendo directamente las condiciones (3) y (4):
\begin{align}
    V_{\text{int}}(r,\theta) &= \sum_{l=0}^{\infty} A_l r^l P_l(\cos\theta), \qquad (r \le R) \\
    V_{\text{ext}}(r,\theta) &= -E_0 r P_1(\cos\theta) + \sum_{l=0}^{\infty} \frac{B_l}{r^{l+1}} P_l(\cos\theta), \qquad (r \ge R)
\end{align}

Aplicamos la primera condición de contorno ($V_{\text{int}}(R,\theta) = V_{\text{ext}}(R,\theta)$):
\begin{equation}
    \sum_{l=0}^{\infty} A_l R^l P_l(\cos\theta) = -E_0 R P_1(\cos\theta) + \sum_{l=0}^{\infty} \frac{B_l}{R^{l+1}} P_l(\cos\theta)
\end{equation}
Por la independencia lineal y ortogonalidad de los polinomios de Legendre $P_l(\cos\theta)$:
\begin{align}
    \text{Para } l \neq 1: \quad & A_l R^l = \frac{B_l}{R^{l+1}} \implies B_l = A_l R^{2l+1} \label{eq:cond1_lneq1} \\
    \text{Para } l = 1: \quad & A_1 R = -E_0 R + \frac{B_1}{R^2} \implies A_1 - \frac{B_1}{R^3} = -E_0 \label{eq:cond1_leq1}
\end{align}

Ahora derivamos $V_{\text{int}}$ y $V_{\text{ext}}$ respecto a $r$ para aplicar la segunda condición de contorno:
\begin{align}
    \pdv{V_{\text{int}}}{r} &= \sum_{l=0}^{\infty} l A_l r^{l-1} P_l(\cos\theta) \\
    \pdv{V_{\text{ext}}}{r} &= -E_0 P_1(\cos\theta) - \sum_{l=0}^{\infty} (l+1)\frac{B_l}{r^{l+2}} P_l(\cos\theta)
\end{align}
Multiplicando $\left.\pdv{V_{\text{int}}}{r}\right|_{r=R}$ por $\varepsilon_r$ e igualando a $\left.\pdv{V_{\text{ext}}}{r}\right|_{r=R}$:
\begin{equation}
    \sum_{l=0}^{\infty} \varepsilon_r l A_l R^{l-1} P_l(\cos\theta) = -E_0 P_1(\cos\theta) - \sum_{l=0}^{\infty} (l+1)\frac{B_l}{R^{l+2}} P_l(\cos\theta)
\end{equation}
Nuevamente separamos por órdenes multipolares $l$:
\begin{itemize}
    \item \textbf{Para $l \neq 1$:}
    \begin{equation}
        \varepsilon_r l A_l R^{l-1} = -(l+1)\frac{B_l}{R^{l+2}}
    \end{equation}
    Sustituyendo $B_l = A_l R^{2l+1}$ de la ecuación \eqref{eq:cond1_lneq1}:
    \begin{equation}
        \varepsilon_r l A_l R^{l-1} = -(l+1) A_l R^{l-1} \implies \big(\varepsilon_r l + l + 1\big) A_l R^{l-1} = 0
    \end{equation}
    Como $\varepsilon_r > 0$ y $l \ge 0$, el factor $(\varepsilon_r l + l + 1)$ es estrictamente positivo y nunca puede anularse. Por consiguiente:
    \begin{equation}
        A_l = 0 \quad \text{y} \quad B_l = 0 \qquad (\forall l \neq 1)
    \end{equation}

    \item \textbf{Para $l = 1$:}
    \begin{equation}
        \varepsilon_r (1) A_1 R^{0} = -E_0 - (1+1)\frac{B_1}{R^3} \implies \varepsilon_r A_1 + \frac{2 B_1}{R^3} = -E_0 \label{eq:cond2_leq1}
    \end{equation}
\end{itemize}

Resolvemos el sistema lineal de $2 \times 2$ formado por las ecuaciones \eqref{eq:cond1_leq1} y \eqref{eq:cond2_leq1}:
\begin{equation}
    \begin{cases}
        A_1 - \dfrac{B_1}{R^3} = -E_0 \\[8pt]
        \varepsilon_r A_1 + \dfrac{2 B_1}{R^3} = -E_0
    \end{cases}
\end{equation}
Multiplicamos la primera ecuación por $2$ y la sumamos a la segunda para eliminar $B_1$:
\begin{equation}
    (2 + \varepsilon_r) A_1 = -3 E_0 \implies A_1 = -\frac{3}{\varepsilon_r + 2} E_0
\end{equation}
Sustituimos $A_1$ en la primera ecuación para despejar $B_1$:
\begin{equation}
    \frac{B_1}{R^3} = A_1 + E_0 = \left( 1 - \frac{3}{\varepsilon_r + 2} \right) E_0 = \left( \frac{\varepsilon_r + 2 - 3}{\varepsilon_r + 2} \right) E_0 = \frac{\varepsilon_r - 1}{\varepsilon_r + 2} E_0 \implies B_1 = \left(\frac{\varepsilon_r - 1}{\varepsilon_r + 2}\right) R^3 E_0
\end{equation}

Finalmente, escribimos los potenciales y campos eléctricos en cada región:
\begin{enumerate}
    \item \textbf{Interior de la esfera ($r \le R$):}
    \begin{equation}
        V_{\text{int}}(r,\theta) = -\left(\frac{3}{\varepsilon_r + 2}\right) E_0 r \cos\theta = -\left(\frac{3}{\varepsilon_r + 2}\right) E_0 z
    \end{equation}
    El campo eléctrico interno es:
    \begin{equation}
        \Evec_{\text{int}} = -\grad V_{\text{int}} = \left(\frac{3}{\varepsilon_r + 2}\right) E_0 \zhat = \frac{3}{\varepsilon_r + 2} \Evec_0
    \end{equation}
    El campo dentro de la esfera es uniforme y paralelo a $\Evec_0$, pero atenuado por el factor $\frac{3}{\varepsilon_r+2} < 1$ (efecto de apantallamiento por despolarización). La polarización inducida es uniforme:
    \begin{equation}
        \Pvec = \varepsilon_0 \chi_e \Evec_{\text{int}} = \varepsilon_0 (\varepsilon_r - 1)\left(\frac{3}{\varepsilon_r + 2}\right)\Evec_0 = 3\varepsilon_0 \left(\frac{\varepsilon_r - 1}{\varepsilon_r + 2}\right)\Evec_0
    \end{equation}

    \item \textbf{Exterior de la esfera ($r \ge R$):}
    \begin{equation}
        V_{\text{ext}}(r,\theta) = -E_0 r \cos\theta + \left(\frac{\varepsilon_r - 1}{\varepsilon_r + 2}\right) E_0 \frac{R^3}{r^2}\cos\theta
    \end{equation}
    El campo eléctrico exterior es la superposición del campo uniforme original $\Evec_0$ más el campo de un dipolo puntual inducido con momento dipolar $\bm{p} = \frac{4}{3}\pi R^3 \Pvec = 4\pi\varepsilon_0 \left(\frac{\varepsilon_r - 1}{\varepsilon_r + 2}\right) R^3 \Evec_0$:
    \begin{equation}
        \Evec_{\text{ext}}(r,\theta) = \Evec_0 + \left(\frac{\varepsilon_r - 1}{\varepsilon_r + 2}\right) E_0 \frac{R^3}{r^3}\left(2\cos\theta\,\rhat + \sin\theta\,\thetahat\right)
    \end{equation}
\end{enumerate}
\end{solucion}

% =========================================================================
% CAPÍTULO 5: MAGNETOSTÁTICA
% =========================================================================
\chapter{Magnetostática}

\section{Fuerza de Lorentz, Corrientes y Ley de Biot-Savart}

\begin{ejercicio}{Campo magnético en el eje de una espira poligonal regular y límite circular}{poligono_biot_savart}
Un alambre delgado lleva una corriente estacionaria $I$ y tiene la forma de un polígono regular de $N$ lados inscrito en un círculo de radio $R$ sobre el plano $xy$, centrado en el origen.
\begin{enumerate}[label=\alph*)]
    \item Utilizando la ley de Biot-Savart, demuestre que el campo magnético en un punto sobre el eje $z$ a una altura $z$ del centro es:
    \begin{equation*}
        \Bvec(z) = \frac{\mu_0 N I}{2\pi} \frac{R^2 \cos(\pi/N)\sin(\pi/N)}{\sqrt{z^2 + R^2\cos^2(\pi/N)}\,\big(z^2 + R^2\big)} \zhat
    \end{equation*}
    \item Verifique explícitamente que en el límite $N \to \infty$ se recupera la fórmula estándar para el campo magnético en el eje de una espira circular de radio $R$.
\end{enumerate}
\end{ejercicio}

\begin{solucion}
\textbf{Parte a) Cálculo del campo mediante la ley de Biot-Savart.}

El polígono regular de $N$ lados está inscrito en una circunferencia de radio $R$, lo que significa que la distancia desde el origen $O(0,0,0)$ hasta cada uno de los $N$ vértices del polígono es $R$.

Cada lado del polígono subtiende un ángulo central $\Delta\phi = \frac{2\pi}{N}$ visto desde el origen. Consideremos uno de los lados del polígono; por conveniencia geométrica, orientemos los ejes $x$ e $y$ en el plano del polígono de modo que este segmento rectilíneo sea paralelo al eje $y$ y esté situado a una distancia perpendicular fija $x = d$ del origen (el apotema del polígono).

El ángulo que forma el radio hacia un extremo del segmento con el eje $x$ es la mitad del ángulo central, es decir, $\alpha = \frac{\Delta\phi}{2} = \frac{\pi}{N}$. Por trigonometría elemental en el triángulo rectángulo formado por el origen, el punto medio del segmento y un vértice:
\begin{align}
    \text{Apotema (distancia perpendicular al centro):} \quad d &= R \cos\left(\frac{\pi}{N}\right) \\
    \text{Semilongitud del segmento:} \quad \frac{L}{2} &= R \sin\left(\frac{\pi}{N}\right)
\end{align}
Así, el segmento va desde $y_1 = -R\sin(\pi/N)$ hasta $y_2 = +R\sin(\pi/N)$, manteniendo constante $x = d = R\cos(\pi/N)$ y $z' = 0$.

Calculemos el campo magnético $d\Bvec_1$ producido en el punto de campo $\bm{r} = (0,0,z) = z\zhat$ por un elemento diferencial de este segmento situado en $\bm{r}' = d\xhat + y'\yhat$, por el cual circula corriente en la dirección $+\yhat$ (sentido antihorario visto desde $+z$), de modo que $d\bm{l}' = dy'\,\yhat$.

El vector de separación $\rsep = \bm{r} - \bm{r}'$ y su magnitud son:
\begin{equation}
    \rsep = -d\xhat - y'\yhat + z\zhat, \qquad \mathscr{r} = \sqrt{d^2 + y'^2 + z^2}
\end{equation}
Calculamos el producto vectorial $d\bm{l}' \times \rsep$:
\begin{equation}
    d\bm{l}' \times \rsep = (dy'\,\yhat) \times (-d\xhat - y'\yhat + z\zhat) = dy'\big( d(\xhat\times\yhat)_{\text{inv}} + z(\yhat\times\zhat) \big) = dy' (z\xhat + d\zhat)
\end{equation}
(recordando que $\yhat \times (-\xhat) = +\zhat$ y $\yhat \times \zhat = +\xhat$).

Al sumar las contribuciones de los $N$ lados simétricos dispuestos alrededor del eje $z$, las componentes horizontales perpendiculares al eje $z$ (como el término proporcional a $z\xhat$) se cancelan exactamente por simetría de rotación discreta $C_N$. Únicamente sobrevive la componente a lo largo de $\zhat$, y cada uno de los $N$ lados aporta exactamente el mismo valor a $B_z$. Por lo tanto, el campo total en el eje $z$ es $N$ veces la componente $z$ de un solo lado:
\begin{equation}
    \Bvec(z) = N B_{1,z} \zhat
\end{equation}
Aplicando la ley de Biot-Savart para hallar $B_{1,z}$:
\begin{equation}
    B_{1,z} = \frac{\mu_0 I}{4\pi} \int_{-L/2}^{L/2} \frac{(d\bm{l}' \times \rsep)_z}{\mathscr{r}^3} = \frac{\mu_0 I d}{4\pi} \int_{-L/2}^{L/2} \frac{dy'}{\left(d^2 + z^2 + y'^2\right)^{3/2}}
\end{equation}
Definimos la constante $a^2 \equiv d^2 + z^2$. Para resolver la integral estándar $\int \frac{dy'}{(a^2 + y'^2)^{3/2}}$, hacemos la sustitución trigonométrica $y' = a \tan\theta \implies dy' = a \sec^2\theta\,d\theta$ y $(a^2 + y'^2)^{3/2} = a^3 \sec^3\theta$:
\begin{equation}
    \int \frac{dy'}{(a^2 + y'^2)^{3/2}} = \int \frac{a\sec^2\theta\,d\theta}{a^3\sec^3\theta} = \frac{1}{a^2}\int \cos\theta\,d\theta = \frac{\sin\theta}{a^2} = \frac{y'}{a^2\sqrt{a^2 + y'^2}}
\end{equation}
Evaluamos la integral definida entre $-L/2$ y $+L/2$:
\begin{equation}
    \int_{-L/2}^{L/2} \frac{dy'}{\left(d^2 + z^2 + y'^2\right)^{3/2}} = \frac{2(L/2)}{(d^2 + z^2)\sqrt{d^2 + z^2 + (L/2)^2}}
\end{equation}
Sustituimos ahora las relaciones geométricas del polígono: $d = R\cos(\pi/N)$ y $L/2 = R\sin(\pi/N)$. Notamos en la raíz cuadrada que:
\begin{equation}
    d^2 + \left(\frac{L}{2}\right)^2 = R^2\cos^2\left(\frac{\pi}{N}\right) + R^2\sin^2\left(\frac{\pi}{N}\right) = R^2
\end{equation}
Por lo tanto, el denominador con la raíz se simplifica limpiamente a:
\begin{equation}
    \sqrt{d^2 + z^2 + (L/2)^2} = \sqrt{z^2 + R^2}
\end{equation}
Reemplazando $d$ y $L/2$ en la expresión de $B_{1,z}$ y multiplicando por $N$:
\begin{align}
    \Bvec(z) &= N \left[ \frac{\mu_0 I \big(R\cos(\pi/N)\big)}{4\pi} \frac{2 R\sin(\pi/N)}{\big(z^2 + R^2\cos^2(\pi/N)\big)\sqrt{z^2 + R^2}} \right] \zhat \nonumber \\
    &= \frac{\mu_0 N I}{2\pi} \frac{R^2 \cos(\pi/N)\sin(\pi/N)}{\big(z^2 + R^2\cos^2(\pi/N)\big)\sqrt{z^2 + R^2}} \zhat
\end{align}
(Nota: En el enunciado original o variantes del libro, según si se inscribe o circunscribe el radio, los factores bajo la raíz permutan; con la definición de vértices a distancia $R$, el término racional es $z^2 + R^2\cos^2(\pi/N)$ y el irracional $\sqrt{z^2+R^2}$).

\medskip
\textbf{Parte b) Límite cuando $N \to \infty$.}

Cuando el número de lados tiende a infinito ($N \to \infty$), el argumento $\frac{\pi}{N} \to 0$. Utilizamos los desarrollos de Taylor de primer orden para funciones trigonométricas de ángulo pequeño:
\begin{equation}
    \cos\left(\frac{\pi}{N}\right) \xrightarrow{N\to\infty} 1, \qquad \sin\left(\frac{\pi}{N}\right) \approx \frac{\pi}{N} \implies N \sin\left(\frac{\pi}{N}\right) \xrightarrow{N\to\infty} \pi
\end{equation}
Tomando el límite en nuestra expresión del campo magnético:
\begin{align}
    \lim_{N\to\infty} \Bvec(z) &= \frac{\mu_0 I}{2\pi} \frac{R^2 (1) (\pi)}{\big(z^2 + R^2(1)^2\big)\sqrt{z^2 + R^2}} \zhat \nonumber \\
    &= \frac{\mu_0 I R^2}{2\big(z^2 + R^2\big)^{3/2}} \zhat
\end{align}
Esta es exactamente la expresión clásica del campo magnético en el eje de una espira circular de radio $R$.
\end{solucion}

% -------------------------------------------------------------------------
\section[Ley de Ampère y potencial vectorial]{Ley de Ampère y el Potencial Vectorial Magnético \texorpdfstring{$\Avec$}{A}}

\begin{ejercicio}{Esfera uniformemente cargada en rotación: Potencial Vectorial y Campo Magnético}{esfera_rotante}
Una cáscara esférica de radio $R$, que porta una densidad superficial de carga uniforme $\sigma$, gira alrededor del eje $z$ con una velocidad angular constante $\bm{\omega} = \omega\zhat$. 
\begin{enumerate}[label=\alph*)]
    \item Calcule el potencial vectorial magnético $\Avec(\bm{r})$ tanto dentro ($r < R$) como fuera ($r > R$) de la cáscara esférica mediante integración directa.
    \item A partir de $\Avec(\bm{r})$, determine el campo magnético $\Bvec(\bm{r}) = \curl{\Avec}$ en el interior y en el exterior de la esfera.
\end{enumerate}
\end{ejercicio}

\begin{solucion}
\textbf{Parte a) Cálculo del potencial vectorial $\Avec(\bm{r})$.}

La densidad superficial de corriente $\bm{K}(\bm{r}')$ sobre la cáscara esférica en un punto $\bm{r}'$ viene dada por el producto de la densidad superficial de carga $\sigma$ y la velocidad lineal $\bm{v}(\bm{r}') = \bm{\omega} \times \bm{r}'$:
\begin{equation}
    \bm{K}(\bm{r}') = \sigma \bm{v}(\bm{r}') = \sigma (\bm{\omega} \times \bm{r}')
\end{equation}
El potencial vectorial magnético generado por una corriente superficial es:
\begin{equation}
    \Avec(\bm{r}) = \frac{\mu_0}{4\pi} \oint_{\mathcal{S}'} \frac{\bm{K}(\bm{r}')}{\mathscr{r}} da'
\end{equation}
donde $\mathscr{r} = |\bm{r} - \bm{r}'|$. Para evaluar esta integral sin complicaciones trigonométricas innecesarias, empleamos el truco geométrico de Griffiths: en lugar de colocar el eje polar de integración a lo largo de $\bm{\omega}$, \textbf{alineamos temporalmente el eje $z'$ del sistema de coordenadas esféricas con el punto de campo $\bm{r}$}, de forma que $\bm{r} = r\zhat'$, mientras que el vector velocidad angular $\bm{\omega}$ yace en el plano $x'z'$ formando un ángulo $\psi$ con $\bm{r}$:
\begin{align}
    \bm{r} &= (0, 0, r) \\
    \bm{\omega} &= (\omega\sin\psi, 0, \omega\cos\psi) \\
    \bm{r}' &= (R\sin\theta'\cos\phi', R\sin\theta'\sin\phi', R\cos\theta')
\end{align}
Con esta elección, la distancia $\mathscr{r} = |\bm{r} - \bm{r}'|$ no depende del ángulo azimutal $\phi'$, pues el ángulo entre $\bm{r}$ y $\bm{r}'$ es simplemente $\theta'$:
\begin{equation}
    \mathscr{r} = \sqrt{r^2 + R^2 - 2rR\cos\theta'}
\end{equation}
El diferencial de área sobre la esfera de radio $R$ es $da' = R^2\sin\theta'\,d\theta'\,d\phi'$. Calculamos ahora explícitamente el producto vectorial $\bm{\omega} \times \bm{r}'$ en estas coordenadas $(x', y', z')$:
\begin{align}
    \bm{\omega} \times \bm{r}' &= \begin{vmatrix}
        \xhat' & \yhat' & \zhat' \\
        \omega\sin\psi & 0 & \omega\cos\psi \\
        R\sin\theta'\cos\phi' & R\sin\theta'\sin\phi' & R\cos\theta'
    \end{vmatrix} \nonumber \\
    &= \omega R \Big[ (-\cos\psi\sin\theta'\sin\phi')\xhat' + (\cos\psi\sin\theta'\cos\phi' - \sin\psi\cos\theta')\yhat' + (\sin\psi\sin\theta'\sin\phi')\zhat' \Big]
\end{align}

Al integrar respecto al ángulo azimutal $\phi'$ desde $0$ hasta $2\pi$, teniendo en cuenta que $\mathscr{r}$ no depende de $\phi'$, usamos que:
\begin{equation}
    \int_{0}^{2\pi} \sin\phi'\,d\phi' = 0, \qquad \int_{0}^{2\pi} \cos\phi'\,d\phi' = 0, \qquad \int_{0}^{2\pi} d\phi' = 2\pi
\end{equation}
Todos los términos que contienen $\sin\phi'$ o $\cos\phi'$ se anulan exactamente al integrar en $\phi'$. El único término que sobrevive es el segundo término de la componente $\yhat'$:
\begin{equation}
    \int_{0}^{2\pi} (\bm{\omega} \times \bm{r}')\,d\phi' = -2\pi\omega R \sin\psi \cos\theta'\,\yhat'
\end{equation}
Observemos qué significa geométricamente el vector $-\omega\sin\psi\,\yhat'$. Dado que $\bm{\omega} = \omega\sin\psi\xhat' + \omega\cos\psi\zhat'$ y el vector unitario en la dirección de $\bm{r}$ es $\rhat = \zhat'$, su producto vectorial es:
\begin{equation}
    \bm{\omega} \times \bm{r} = (\omega\sin\psi\xhat' + \omega\cos\psi\zhat') \times (r\zhat') = -\omega r \sin\psi\,\yhat' \implies \bm{\omega} \times \rhat = -\omega\sin\psi\,\yhat'
\end{equation}
¡Esta identificación es independiente del sistema de coordenadas auxiliar! Sustituyendo esto en la integral de $\Avec(\bm{r})$, nos queda únicamente la integral polar en $\theta'$:
\begin{equation}
    \Avec(\bm{r}) = \frac{\mu_0 \sigma R^3}{2} (\bm{\omega} \times \rhat) \int_{0}^{\pi} \frac{\cos\theta'\sin\theta'}{\sqrt{r^2 + R^2 - 2rR\cos\theta'}} d\theta'
\end{equation}

Resolvemos la integral definida $I = \int_{0}^{\pi} \frac{\cos\theta'\sin\theta'}{\sqrt{r^2 + R^2 - 2rR\cos\theta'}} d\theta'$ mediante el cambio de variable:
\begin{equation}
    u = \sqrt{r^2 + R^2 - 2rR\cos\theta'} \implies u^2 = r^2 + R^2 - 2rR\cos\theta'
\end{equation}
Diferenciando respecto a $\theta'$ y despejando $\cos\theta'$:
\begin{equation}
    2u\,du = 2rR\sin\theta'\,d\theta' \implies \sin\theta'\,d\theta' = \frac{u\,du}{rR}, \qquad \cos\theta' = \frac{r^2 + R^2 - u^2}{2rR}
\end{equation}
Cambiando los límites de integración:
\begin{itemize}
    \item Cuando $\theta' = 0$: $u_{\text{inf}} = \sqrt{r^2 + R^2 - 2rR} = |r - R|$.
    \item Cuando $\theta' = \pi$: $u_{\text{sup}} = \sqrt{r^2 + R^2 + 2rR} = r + R$.
\end{itemize}
Sustituimos todo en la integral $I$:
\begin{align}
    I &= \int_{|r-R|}^{r+R} \left(\frac{r^2 + R^2 - u^2}{2rR}\right) \frac{1}{u} \left(\frac{u\,du}{rR}\right) = \frac{1}{2r^2 R^2} \int_{|r-R|}^{r+R} \big(r^2 + R^2 - u^2\big) du \nonumber \\
    &= \frac{1}{2r^2 R^2} \left[ (r^2 + R^2)u - \frac{u^3}{3} \right]_{|r-R|}^{r+R}
\end{align}
Distinguimos ahora los dos casos debido al valor absoluto $|r - R|$:
\begin{enumerate}
    \item \textbf{Interior de la esfera ($r < R$):} Aquí $|r - R| = R - r$. Evaluamos el corchete:
    \begin{align}
        \left[ (r^2 + R^2)u - \frac{u^3}{3} \right]_{R-r}^{R+r} &= (r^2 + R^2)\big[(R+r) - (R-r)\big] - \frac{1}{3}\big[(R+r)^3 - (R-r)^3\big] \nonumber \\
        &= (r^2 + R^2)(2r) - \frac{1}{3}\big[ (R^3 + 3R^2 r + 3R r^2 + r^3) - (R^3 - 3R^2 r + 3R r^2 - r^3) \big] \nonumber \\
        &= 2r^3 + 2R^2 r - \frac{1}{3}\big(6R^2 r + 2r^3\big) = 2r^3 + 2R^2 r - 2R^2 r - \frac{2}{3}r^3 = \frac{4}{3}r^3
    \end{align}
    Por tanto, para $r < R$:
    \begin{equation}
        I_{\text{int}} = \frac{1}{2r^2 R^2}\left(\frac{4}{3}r^3\right) = \frac{2r}{3R^2}
    \end{equation}
    Multiplicando por el prefactor de $\Avec$, y volviendo al sistema de coordenadas habitual donde $\bm{\omega} = \omega\zhat$ y $\bm{\omega} \times \rhat = \omega\sin\theta\,\phihat$ (o bien $\bm{\omega}\times\bm{r} = \omega r \sin\theta\,\phihat$):
    \begin{equation}
        \Avec_{\text{int}}(\bm{r}) = \frac{\mu_0 \sigma R^3}{2} (\bm{\omega} \times \rhat) \left(\frac{2r}{3R^2}\right) = \frac{\mu_0 R \sigma}{3} (\bm{\omega} \times \bm{r}) = \frac{\mu_0 R \omega \sigma}{3} r \sin\theta\,\phihat \qquad (r \le R)
    \end{equation}

    \item \textbf{Exterior de la esfera ($r > R$):} Aquí $|r - R| = r - R$. Por la total simetría algebraica entre $r$ y $R$ en el integrando, el resultado del corchete se obtiene intercambiando $r \leftrightarrow R$ en el resultado anterior, es decir, da $\frac{4}{3}R^3$:
    \begin{equation}
        I_{\text{ext}} = \frac{1}{2r^2 R^2}\left(\frac{4}{3}R^3\right) = \frac{2R}{3r^2}
    \end{equation}
    Sustituyendo en $\Avec(\bm{r})$:
    \begin{equation}
        \Avec_{\text{ext}}(\bm{r}) = \frac{\mu_0 \sigma R^3}{2} (\bm{\omega} \times \rhat) \left(\frac{2R}{3r^2}\right) = \frac{\mu_0 R^4 \sigma}{3r^2} (\bm{\omega} \times \rhat) = \frac{\mu_0 R^4 \omega \sigma}{3} \frac{\sin\theta}{r^2}\,\phihat \qquad (r \ge R)
    \end{equation}
\end{enumerate}

\medskip
\textbf{Parte b) Cálculo del campo magnético $\Bvec = \curl{\Avec}$.}

\begin{itemize}
    \item \textbf{En el interior ($r < R$):} Podemos calcular el rotacional directamente en coordenadas cartesianas notando que $\bm{\omega} = \omega\zhat$ y $\bm{r} = x\xhat + y\yhat + z\zhat$, de modo que $\bm{\omega} \times \bm{r} = \omega(-y\xhat + x\yhat)$:
    \begin{equation}
        \Bvec_{\text{int}} = \curl{\Avec_{\text{int}}} = \frac{\mu_0 R \sigma}{3} \curl(\bm{\omega} \times \bm{r}) = \frac{\mu_0 R \omega \sigma}{3} \left(\pdv{x}{x} - \pdv{(-y)}{y}\right)\zhat = \frac{2}{3}\mu_0 \sigma R \omega \zhat = \frac{2}{3}\mu_0 \sigma R \bm{\omega}
    \end{equation}
    ¡El campo magnético en el interior de una cáscara esférica cargada en rotación es \textbf{completamente uniforme}!

    \item \textbf{En el exterior ($r > R$):} Usamos la expresión del rotacional en coordenadas esféricas para un campo que solo tiene componente $\phi$, es decir $\Avec = A_\phi(r,\theta)\phihat$:
    \begin{equation}
        \curl{\Avec} = \frac{1}{r\sin\theta}\pdv{}{\theta}\left(\sin\theta\,A_\phi\right)\rhat - \frac{1}{r}\pdv{}{r}\left(r A_\phi\right)\thetahat
    \end{equation}
    Con $A_\phi = \frac{\mu_0 R^4 \omega \sigma}{3} \frac{\sin\theta}{r^2}$, calculamos cada término:
    \begin{align}
        B_r &= \frac{\mu_0 R^4 \omega \sigma}{3 r^3 \sin\theta} \dv{}{\theta}\left(\sin^2\theta\right) = \frac{\mu_0 R^4 \omega \sigma}{3 r^3 \sin\theta} (2\sin\theta\cos\theta) = \frac{2\mu_0 R^4 \omega \sigma}{3 r^3}\cos\theta \\
        B_\theta &= -\frac{\mu_0 R^4 \omega \sigma \sin\theta}{3 r} \dv{}{r}\left(\frac{1}{r}\right) = -\frac{\mu_0 R^4 \omega \sigma \sin\theta}{3 r}\left(-\frac{1}{r^2}\right) = \frac{\mu_0 R^4 \omega \sigma}{3 r^3}\sin\theta
    \end{align}
    Agrupando componentes:
    \begin{equation}
        \Bvec_{\text{ext}}(r,\theta) = \frac{\mu_0 R^4 \omega \sigma}{3 r^3}\left(2\cos\theta\,\rhat + \sin\theta\,\thetahat\right)
    \end{equation}
    Este es exactamente el campo magnético de un dipolo magnético puntual ideal en el origen con momento dipolar $\bm{m} = \frac{4}{3}\pi R^4 \sigma \bm{\omega}$.
\end{itemize}
\end{solucion}

% =========================================================================
% CAPÍTULO 6: CAMPOS MAGNÉTICOS EN LA MATERIA
% =========================================================================
\chapter[Campos magnéticos en la materia]{Campos Magnéticos en la Materia (Magnetización y Medios Magnéticos)}

\section[Magnetización y campo auxiliar]{Magnetización \texorpdfstring{$\Mvec$}{M}, Corrientes Ligadas y el Campo Auxiliar \texorpdfstring{$\Hvec$}{H}}

\begin{ejercicio}{Cilindro largo con magnetización azimutal o transversal}{cilindro_magnetizado}
Un cilindro circular infinitamente largo de radio $R$ porta una magnetización "congelada" paralela a su eje que varía radialmente según $\Mvec(s) = k s^2 \zhat$, donde $k$ es una constante y $s$ es la distancia al eje del cilindro. No hay corrientes libres en ninguna parte del espacio ($\Jvec_f = \bm{0}$).
\begin{enumerate}[label=\alph*)]
    \item Encuentre las densidades de corriente ligada de volumen $\Jvec_b$ y de superficie $\bm{K}_b$, y a partir de ellas calcule el campo magnético $\Bvec(\bm{r})$ dentro y fuera del cilindro usando la ley de Ampère ordinaria.
    \item Utilice la ley de Ampère para el campo auxiliar $\Hvec$ ($\oint \Hvec \cdot d\bm{l} = I_{f,\text{enc}}$) para hallar $\Hvec(\bm{r})$ y deduzca $\Bvec(\bm{r})$ directamente en un solo paso, comprobando que coincide con la parte (a).
\end{enumerate}
\end{ejercicio}

\begin{solucion}
\textbf{Parte a) Método de las corrientes ligadas.}

Las corrientes ligadas equivalentes producidas por una magnetización $\Mvec(\bm{r})$ están definidas por:
\begin{equation}
    \Jvec_b = \curl{\Mvec}, \qquad \bm{K}_b = \Mvec \times \nhat
\end{equation}
En coordenadas cilíndricas $(s,\phi,z)$, la magnetización es $\Mvec = M_z(s)\zhat = k s^2 \zhat$. Calculamos su rotacional:
\begin{equation}
    \Jvec_b = \curl{\Mvec} = -\pdv{M_z}{s}\phihat = -\dv{}{s}\left(k s^2\right)\phihat = -2 k s \,\phihat \qquad (s < R)
\end{equation}
En la superficie lateral del cilindro ($s = R$), el vector normal unitario saliente es $\nhat = \shat$. Calculamos $\bm{K}_b$:
\begin{equation}
    \bm{K}_b = \Mvec(R) \times \shat = \left(k R^2 \zhat\right) \times \shat = +k R^2 \,\phihat
\end{equation}
Observamos la geometría de las corrientes ligadas: tanto $\Jvec_b$ como $\bm{K}_b$ fluyen exclusivamente en la dirección azimutal $\phihat$, exactamente igual que en un conjunto de solenoides coaxiales infinitos. Para un solenoide infinito con corrientes en $\phihat$, el campo magnético $\Bvec$ apunta en la dirección axial $\zhat$ y se anula fuera del cilindro si la corriente neta por unidad de longitud es cero. Verifiquemos la corriente ligada total por unidad de longitud axial $\Delta z$:
\begin{equation}
    I_{b,\text{por unidad de }z} = \int_{0}^{R} J_{b,\phi}(s)\,ds + K_{b,\phi} = \int_{0}^{R} (-2ks)\,ds + k R^2 = -k R^2 + k R^2 = 0
\end{equation}
Por lo tanto, para $s > R$ el campo magnético es nulo:
\begin{equation}
    \Bvec(\bm{r}) = \bm{0} \qquad (s > R)
\end{equation}
Para hallar $\Bvec(s) = B_z(s)\zhat$ en el interior ($s < R$), tomamos una espira amperiana rectangular en el plano $sz$ de longitud $L$ paralela al eje $z$, con un lado interior a una distancia $s < R$ del eje y el lado exterior fuera del cilindro en $s_{\text{ext}} > R$ (donde $B = 0$). Aplicando la ley de Ampère $\oint \Bvec \cdot d\bm{l} = \mu_0 I_{b,\text{enc}}$:
\begin{equation}
    B_z(s) L = \mu_0 I_{b,\text{enc}}
\end{equation}
La corriente ligada que atraviesa la espira amperiana entre el radio $s$ y la superficie $R$ comprende la integral de $\Jvec_b$ desde $s' = s$ hasta $s' = R$, más la corriente superficial $\bm{K}_b$ en $s'=R$:
\begin{align}
    I_{b,\text{enc}} &= L \left[ \int_{s}^{R} J_{b,\phi}(s')\,ds' + K_{b,\phi} \right] = L \left[ \int_{s}^{R} (-2k s')\,ds' + k R^2 \right] \nonumber \\
    &= L \Big[ -\big(R^2 - s^2\big)k + k R^2 \Big] = L k s^2
\end{align}
Sustituyendo en la ley de Ampère y cancelando $L$:
\begin{equation}
    B_z(s) = \mu_0 k s^2 \implies \Bvec(\bm{r}) = \mu_0 k s^2 \zhat = \mu_0 \Mvec \qquad (s < R)
\end{equation}

\medskip
\textbf{Parte b) Método del campo auxiliar $\Hvec$.}

La ley de Ampère para $\Hvec$ en forma diferencial e integral es:
\begin{equation}
    \curl{\Hvec} = \Jvec_f = \bm{0}, \qquad \oint \Hvec \cdot d\bm{l} = I_{f,\text{enc}} = 0
\end{equation}
Además, la divergencia de $\Bvec = \mu_0(\Hvec + \Mvec)$ es siempre cero:
\begin{equation}
    \div{\Bvec} = 0 \implies \div{\Hvec} = -\div{\Mvec}
\end{equation}
En nuestro problema, $\Mvec = k s^2 \zhat$ no depende de $z$, por lo que:
\begin{equation}
    \div{\Mvec} = \pdv{M_z}{z} = 0
\end{equation}
Asimismo, en la superficie del cilindro $s = R$, la componente normal de $\Mvec$ es nula ($\Mvec \cdot \shat = 0$), por lo que no hay "carga magnética ficticia" superficial ni volumétrica. Como $\div{\Hvec} = 0$ y $\curl{\Hvec} = \bm{0}$ en todo el espacio y $\Hvec \to \bm{0}$ en el infinito, la única solución es que el campo auxiliar sea idénticamente nulo en todas partes:
\begin{equation}
    \Hvec(\bm{r}) = \bm{0} \qquad (\text{en todo el espacio})
\end{equation}
A partir de la definición constitutiva general $\Bvec = \mu_0(\Hvec + \Mvec)$, obtenemos de manera inmediata:
\begin{equation}
    \Bvec(\bm{r}) = \mu_0 \Mvec(\bm{r}) = \begin{cases}
        \mu_0 k s^2 \zhat, & \text{para } s < R \\
        \bm{0}, & \text{para } s > R
    \end{cases}
\end{equation}
Lo cual coincide exactamente con el resultado de la parte (a).
\end{solucion}

% -------------------------------------------------------------------------
\section{Problemas de Contorno en Medios Magnéticos Lineales}

\begin{ejercicio}{Blindaje magnético mediante una cáscara esférica permeable}{blindaje_magnetico}
Una cáscara esférica de permeabilidad magnética lineal $\mu$ (permeabilidad relativa $\mu_r = \mu/\mu_0$), radio interno $a$ y radio externo $b$, se coloca en un campo magnético externo inicialmente uniforme $\Bvec_0 = B_0 \zhat$. El interior de la cavidad ($r < a$) y el exterior de la cáscara ($r > b$) están vacíos ($\mu = \mu_0$). 
\begin{enumerate}[label=\alph*)]
    \item Utilizando el potencial escalar magnético $W(\bm{r})$ (tal que $\Hvec = -\grad W$), plantee las condiciones de contorno en $r=a$ y $r=b$ y demuestre que el campo magnético dentro de la cavidad hueca ($r < a$) es uniforme e igual a:
    \begin{equation*}
        \Bvec_{\text{cav}} = \left[ \frac{9\mu_r}{(2\mu_r + 1)(\mu_r + 2) - 2(\mu_r - 1)^2 (a/b)^3} \right] \Bvec_0
    \end{equation*}
    \item Analice el límite de alta permeabilidad ($\mu_r \gg 1$) para explicar cómo funciona el blindaje magnético (mu-metal).
\end{enumerate}
\end{ejercicio}

\begin{solucion}
\textbf{Parte a) Determinación del campo en la cavidad mediante el potencial escalar magnético.}

Dado que no existen corrientes libres ($\Jvec_f = \bm{0}$) en ninguna región del espacio, tenemos $\curl{\Hvec} = \bm{0}$ en todas partes. Por tanto, podemos expresar $\Hvec$ como el gradiente de un potencial escalar magnético $W$:
\begin{equation}
    \Hvec = -\grad W
\end{equation}
En cada una de las tres regiones homogéneas ($r < a$, $a < r < b$, y $r > b$), la permeabilidad $\mu$ es constante, por lo que $\div{\Bvec} = \mu \div{\Hvec} = -\mu \nabla^2 W = 0$. Así, $W(r,\theta)$ satisface la ecuación de Laplace $\nabla^2 W = 0$ en cada región.

Muy lejos de la esfera ($r \to \infty$), el campo es $\Hvec_0 = \frac{B_0}{\mu_0}\zhat \equiv H_0 \zhat$, lo que implica $W(r,\theta) \to -H_0 r \cos\theta$. Dado que esta condición asintótica es proporcional únicamente a $P_1(\cos\theta) = \cos\theta$, y por la ortogonalidad de los polinomios de Legendre en las esferas concéntricas $r=a$ y $r=b$, solamente los términos con $l=1$ son distintos de cero en las tres regiones:
\begin{equation}
    W(r,\theta) = \begin{cases}
        A r \cos\theta, & \text{para } r \le a \quad (\text{Región I: Cavidad}) \\[6pt]
        \left(B r + \dfrac{C}{r^2}\right)\cos\theta, & \text{para } a \le r \le b \quad (\text{Región II: Material permeable}) \\[10pt]
        \left(-H_0 r + \dfrac{D}{r^2}\right)\cos\theta, & \text{para } r \ge b \quad (\text{Región III: Exterior})
    \end{cases}
\end{equation}
Nótese que en la Región I el campo será $\Hvec_{\text{cav}} = -\grad(A z) = -A\zhat$, de modo que el campo magnético dentro de la cavidad es $\Bvec_{\text{cav}} = \mu_0 \Hvec_{\text{cav}} = -\mu_0 A \zhat$. Nuestro objetivo es despejar la constante $A$ en función de $H_0 = B_0/\mu_0$.

Las condiciones de frontera magnetostáticas sin corrientes superficiales libres ($\bm{K}_f = \bm{0}$) exigen la continuidad de $W$ (que garantiza la continuidad de $H_\parallel$) y la continuidad de la componente normal de $\Bvec$, es decir, $\mu \pdv{W}{r}$:
\begin{enumerate}
    \item \textbf{En $r = a$ (continuidad de $W$):}
    \begin{equation}
        A a = B a + \frac{C}{a^2} \implies A = B + \frac{C}{a^3} \label{eq:mag_bc1}
    \end{equation}
    \item \textbf{En $r = a$ (continuidad de $\mu \pdv{W}{r}$):}
    \begin{equation}
        \mu_0 A = \mu \left( B - \frac{2C}{a^3} \right) \implies A = \mu_r \left( B - \frac{2C}{a^3} \right) \label{eq:mag_bc2}
    \end{equation}
    \item \textbf{En $r = b$ (continuidad de $W$):}
    \begin{equation}
        B b + \frac{C}{b^2} = -H_0 b + \frac{D}{b^2} \implies B + \frac{C}{b^3} = -H_0 + \frac{D}{b^3} \label{eq:mag_bc3}
    \end{equation}
    \item \textbf{En $r = b$ (continuidad de $\mu \pdv{W}{r}$):}
    \begin{equation}
        \mu \left( B - \frac{2C}{b^3} \right) = \mu_0 \left( -H_0 - \frac{2D}{b^3} \right) \implies \mu_r \left( B - \frac{2C}{b^3} \right) = -H_0 - \frac{2D}{b^3} \label{eq:mag_bc4}
    \end{equation}
\end{enumerate}

Procedemos a resolver este sistema algebraico de $4$ ecuaciones con $4$ incógnitas ($A,B,C,D$) de forma ordenada:

\textbf{Paso 1: Expresar $B$ y $C/a^3$ en términos de $A$ usando \eqref{eq:mag_bc1} y \eqref{eq:mag_bc2}.}
De \eqref{eq:mag_bc1} y \eqref{eq:mag_bc2} tenemos:
\begin{equation}
    \begin{cases}
        B + \dfrac{C}{a^3} = A \\[8pt]
        B - \dfrac{2C}{a^3} = \dfrac{A}{\mu_r}
    \end{cases}
\end{equation}
Restando la segunda ecuación de la primera:
\begin{equation}
    \frac{3C}{a^3} = A\left(1 - \frac{1}{\mu_r}\right) = \frac{\mu_r - 1}{\mu_r} A \implies C = \frac{\mu_r - 1}{3\mu_r} a^3 A
\end{equation}
Multiplicando la primera ecuación por $2$ y sumándola a la segunda:
\begin{equation}
    3B = A\left(2 + \frac{1}{\mu_r}\right) = \frac{2\mu_r + 1}{\mu_r} A \implies B = \frac{2\mu_r + 1}{3\mu_r} A
\end{equation}

\textbf{Paso 2: Eliminar $D/b^3$ entre las ecuaciones en $r=b$, \eqref{eq:mag_bc3} y \eqref{eq:mag_bc4}.}
Multiplicamos \eqref{eq:mag_bc3} por $2$ y le sumamos \eqref{eq:mag_bc4}:
\begin{align}
    2\left(B + \frac{C}{b^3}\right) + \mu_r\left(B - \frac{2C}{b^3}\right) &= 2\left(-H_0 + \frac{D}{b^3}\right) + \left(-H_0 - \frac{2D}{b^3}\right) \nonumber \\
    (2 + \mu_r)B + 2(1 - \mu_r)\frac{C}{b^3} &= -3H_0
\end{align}

\textbf{Paso 3: Sustituir las expresiones de $B$ y $C$ en función de $A$.}
Reemplazando $B = \frac{2\mu_r + 1}{3\mu_r} A$ y $\frac{C}{b^3} = \frac{\mu_r - 1}{3\mu_r} \left(\frac{a}{b}\right)^3 A$:
\begin{equation}
    (\mu_r + 2)\left(\frac{2\mu_r + 1}{3\mu_r}\right) A - 2(\mu_r - 1)\left(\frac{\mu_r - 1}{3\mu_r}\right)\left(\frac{a}{b}\right)^3 A = -3H_0
\end{equation}
Multiplicando ambos miembros por $3\mu_r$:
\begin{equation}
    \left[ (\mu_r + 2)(2\mu_r + 1) - 2(\mu_r - 1)^2 \left(\frac{a}{b}\right)^3 \right] A = -9\mu_r H_0
\end{equation}
Despejando $-A$:
\begin{equation}
    -A = \frac{9\mu_r}{(2\mu_r + 1)(\mu_r + 2) - 2(\mu_r - 1)^2 (a/b)^3} H_0
\end{equation}
Como $\Bvec_{\text{cav}} = -\mu_0 A \zhat$ y $\Bvec_0 = \mu_0 H_0 \zhat$, multiplicando por $\mu_0 \zhat$ obtenemos exactamente la relación buscada:
\begin{equation}
    \Bvec_{\text{cav}} = \left[ \frac{9\mu_r}{(2\mu_r + 1)(\mu_r + 2) - 2(\mu_r - 1)^2 (a/b)^3} \right] \Bvec_0
\end{equation}

\medskip
\textbf{Parte b) Análisis del límite $\mu_r \gg 1$ (Blindaje magnético).}

Si fabricamos la cáscara con un material ferromagnético blando de permeabilidad relativa muy elevada ($\mu_r \gg 1$, por ejemplo $\mu_r \sim 10^4$), podemos expandir el denominador conservando los términos dominantes de orden $\mu_r^2$:
\begin{equation}
    (2\mu_r + 1)(\mu_r + 2) - 2(\mu_r - 1)^2 \left(\frac{a}{b}\right)^3 \approx 2\mu_r^2 - 2\mu_r^2 \left(\frac{a}{b}\right)^3 = 2\mu_r^2 \left[ 1 - \left(\frac{a}{b}\right)^3 \right]
\end{equation}
En este régimen asintótico, el campo en la cavidad se reduce a:
\begin{equation}
    \Bvec_{\text{cav}} \approx \frac{9}{2\mu_r \big[1 - (a/b)^3\big]} \Bvec_0 \propto \frac{1}{\mu_r} \Bvec_0 \xrightarrow{\mu_r \to \infty} \bm{0}
\end{equation}
Las líneas de campo magnético externo son "canalizadas" preferentemente por el interior de la pared de alta permeabilidad ($a < r < b$), dejando la cavidad interna prácticamente libre de campo magnético.
\end{solucion}

% =========================================================================
% CAPÍTULO 7: ELECTRODINÁMICA
% =========================================================================
\chapter[Electrodinámica]{Electrodinámica (Ley de Inducción de Faraday y Ecuaciones de Maxwell)}

\section{Fuerza Electromotriz Motriz y Ley de Faraday}

\begin{ejercicio}{Disco de Faraday (Generador Homopolar) y Freno Electromagnético}{barra_rieles_faraday}
Una barra metálica de masa $m$, longitud $l$ y resistencia eléctrica despreciable desliza sin fricción sobre dos rieles conductores horizontales paralelos separados por una distancia $l$. Los rieles están unidos en un extremo por un resistor de resistencia $R$, y todo el sistema se encuentra inmerso en un campo magnético uniforme vertical $\Bvec = B_0 \zhat$. En el instante $t = 0$, se le imprime a la barra una velocidad inicial $\bm{v}_0 = v_0 \xhat$.
\begin{enumerate}[label=\alph*)]
    \item Utilizando la ley de Faraday y la fuerza magnética de Lorentz, plantee y resuelva la ecuación diferencial del movimiento para determinar la velocidad $v(t)$ de la barra y la corriente inducida $I(t)$ en el circuito.
    \item Calcule la energía total disipada por efecto Joule en el resistor desde $t = 0$ hasta $t \to \infty$ y demuestre explícitamente que es igual a la energía cinética inicial de la barra $\frac{1}{2}mv_0^2$.
\end{enumerate}
\end{ejercicio}

\begin{solucion}
\textbf{Parte a) Ecuación de movimiento, velocidad $v(t)$ y corriente $I(t)$.}

Sea $x(t)$ la posición de la barra medida desde el resistor en el instante $t$. El área encerrada por el circuito rectangular formado por el resistor, los dos rieles y la barra móvil es $A(t) = l \cdot x(t)$.

Como el campo magnético $\Bvec = B_0 \zhat$ es uniforme y perpendicular al plano del circuito (eligiendo la normal del circuito como $+\zhat$), el flujo magnético $\Phi_B(t)$ a través de la espira es:
\begin{equation}
    \Phi_B(t) = \int \Bvec \cdot d\bm{a} = B_0 l x(t)
\end{equation}
Por la ley de inducción de Faraday, la fuerza electromotriz (fem) $\mathcal{E}$ inducida alrededor del circuito es:
\begin{equation}
    \mathcal{E}(t) = -\dv{\Phi_B}{t} = -B_0 l \dv{x}{t} = -B_0 l v(t)
\end{equation}
El signo negativo (ley de Lenz) indica que si la barra se mueve hacia la derecha ($v > 0$, aumentando el flujo hacia $+z$), la corriente inducida circula en sentido horario vista desde $+z$ (de modo que su campo propio apunte hacia $-z$). En magnitud, la corriente que circula por el resistor $R$ es:
\begin{equation}
    I(t) = \frac{|\mathcal{E}(t)|}{R} = \frac{B_0 l v(t)}{R}
\end{equation}
Al circular en sentido horario, en la barra móvil (situada en el lado derecho del rectángulo) la corriente fluye en la dirección $-\yhat$, por lo que el vector longitud orientada de la barra es $\bm{l} = -l\yhat$. La fuerza magnética $\bm{F}_{\text{mag}}$ que actúa sobre la barra portadora de corriente es:
\begin{equation}
    \bm{F}_{\text{mag}} = I (\bm{l} \times \Bvec) = I (-l\yhat) \times (B_0 \zhat) = -I l B_0 \xhat = -\frac{B_0^2 l^2}{R} v(t) \xhat
\end{equation}
Como era de esperar por la ley de Lenz, la fuerza magnética se opone directamente al movimiento de la barra. Aplicando la segunda ley de Newton en el eje $x$:
\begin{equation}
    m \dv{v}{t} = -\frac{B_0^2 l^2}{R} v
\end{equation}
Definimos el tiempo característico de amortiguamiento electromagnético $\tau \equiv \frac{mR}{B_0^2 l^2}$. Separando variables e integrando desde $t=0$ (con $v(0) = v_0$) hasta un tiempo $t$:
\begin{equation}
    \int_{v_0}^{v(t)} \frac{dv'}{v'} = -\frac{1}{\tau} \int_{0}^{t} dt' \implies \ln\left(\frac{v(t)}{v_0}\right) = -\frac{t}{\tau}
\end{equation}
Despejando $v(t)$ e $I(t)$:
\begin{align}
    v(t) &= v_0 \exp\left(-\frac{B_0^2 l^2}{m R} t\right) \\
    I(t) &= \frac{B_0 l v_0}{R} \exp\left(-\frac{B_0^2 l^2}{m R} t\right)
\end{align}

\medskip
\textbf{Parte b) Conservación de la energía (Efecto Joule vs. Energía Cinética).}

La potencia instantánea disipada en forma de calor en el resistor de resistencia $R$ es $P(t) = I^2(t) R$. La energía térmica total $W_{\text{Joule}}$ disipada desde $t=0$ hasta $t \to \infty$ es la integral temporal de la potencia:
\begin{align}
    W_{\text{Joule}} &= \int_{0}^{\infty} I^2(t) R \, dt = \int_{0}^{\infty} \left(\frac{B_0^2 l^2 v_0^2}{R^2}\right) \exp\left(-\frac{2 B_0^2 l^2}{m R} t\right) R \, dt \nonumber \\
    &= \frac{B_0^2 l^2 v_0^2}{R} \int_{0}^{\infty} \exp\left(-\frac{2 B_0^2 l^2}{m R} t\right) dt
\end{align}
Resolvemos la integral exponencial impropia:
\begin{equation}
    \int_{0}^{\infty} e^{-\alpha t}\,dt = \left[ -\frac{1}{\alpha} e^{-\alpha t} \right]_{0}^{\infty} = \frac{1}{\alpha}, \qquad \text{con } \alpha = \frac{2 B_0^2 l^2}{m R}
\end{equation}
Sustituyendo el valor de la integral:
\begin{equation}
    W_{\text{Joule}} = \left(\frac{B_0^2 l^2 v_0^2}{R}\right) \left(\frac{m R}{2 B_0^2 l^2}\right) = \frac{1}{2} m v_0^2
\end{equation}
Queda demostrado exactamente que el $100\%$ de la energía cinética inicial de la barra se transforma en energía eléctrica y se disipa en el resistor.
\end{solucion}

% -------------------------------------------------------------------------
\section{Inductancia y Energía Magnética}

\begin{ejercicio}{Inductancia mutua entre un alambre recto y una espira triangular, y autoinductancia de un cable coaxial}{inductancia_coaxial}
Considere un cable coaxial largo de longitud $l$ que consta de un cilindro conductor macizo interior de radio $a$ por el cual circula una corriente estacionaria $I$ distribuida uniformemente en su sección transversal (en dirección $+\zhat$), y una cáscara cilíndrica conductora exterior muy delgada y concéntrica de radio $b > a$ por la que retorna la misma corriente $I$ (en dirección $-\zhat$).
\begin{enumerate}[label=\alph*)]
    \item Determine el campo magnético $\Bvec(s)$ en todas las regiones ($s < a$, $a < s < b$, y $s > b$).
    \item Calcule la autoinductancia por unidad de longitud $\mathcal{L} = L/l$ del cable coaxial utilizando el método de la energía magnética almacenada $W = \frac{1}{2}LI^2 = \frac{1}{2\mu_0}\int B^2 d\tau$. (Nota: el cálculo por flujo simple falla o requiere cuidado extremo cuando la corriente está distribuida en un volumen, mientras que el método energético es exacto).
\end{enumerate}
\end{ejercicio}

\begin{solucion}
\textbf{Parte a) Cálculo del campo magnético mediante la ley de Ampère.}

Por simetría cilíndrica, las líneas de campo magnético son circunferencias concéntricas con el eje $z$, es decir, $\Bvec(\bm{r}) = B_\phi(s)\phihat$. Tomamos circunferencias amperianas de radio $s$:
\begin{equation}
    \oint \Bvec \cdot d\bm{l} = B_\phi(s) (2\pi s) = \mu_0 I_{\text{enc}}
\end{equation}
Analizamos cada una de las tres regiones:
\begin{enumerate}[label=\roman*)]
    \item \textbf{Interior del conductor central ($0 \le s < a$):}
    La densidad de corriente es uniforme en la sección transversal de área $\pi a^2$, por lo que $\Jvec = \frac{I}{\pi a^2}\zhat$. La corriente encerrada por un círculo de radio $s < a$ es proporcional al área:
    \begin{equation}
        I_{\text{enc}}(s) = J (\pi s^2) = I \frac{s^2}{a^2}
    \end{equation}
    Sustituyendo en la ley de Ampère:
    \begin{equation}
        B_\phi(s) (2\pi s) = \mu_0 I \frac{s^2}{a^2} \implies \Bvec(s) = \frac{\mu_0 I s}{2\pi a^2}\,\phihat \qquad (s < a)
    \end{equation}

    \item \textbf{Región entre los dos conductores ($a < s < b$):}
    Aquí se encierra la totalidad de la corriente $I$ del conductor central:
    \begin{equation}
        I_{\text{enc}} = I \implies \Bvec(s) = \frac{\mu_0 I}{2\pi s}\,\phihat \qquad (a < s < b)
    \end{equation}

    \item \textbf{Exterior del cable coaxial ($s > b$):}
    La corriente encerrada neta es la suma algebraica de la corriente de ida $+I$ y la corriente de retorno $-I$ en la cáscara $s=b$:
    \begin{equation}
        I_{\text{enc}} = I - I = 0 \implies \Bvec(s) = \bm{0} \qquad (s > b)
    \end{equation}
\end{enumerate}

\medskip
\textbf{Parte b) Cálculo de la autoinductancia mediante la energía magnética.}

La energía magnetostática almacenada en un tramo de longitud $l$ del cable coaxial (asumiendo $\mu = \mu_0$ en el conductor no ferromagnético y en el espacio intermedio) es:
\begin{equation}
    W = \frac{1}{2\mu_0} \int B^2 \, d\tau
\end{equation}
En coordenadas cilíndricas, el elemento de volumen para un tubo de radio $s$, espesor $ds$ y longitud $l$ es $d\tau = 2\pi s \, l \, ds$. Separamos la integral en las dos regiones donde $B \neq 0$:
\begin{equation}
    W = \frac{1}{2\mu_0} \left[ \int_{0}^{a} \left(\frac{\mu_0 I s}{2\pi a^2}\right)^2 2\pi s l \, ds + \int_{a}^{b} \left(\frac{\mu_0 I}{2\pi s}\right)^2 2\pi s l \, ds \right]
\end{equation}
Calculamos la primera integral (contribución interna del conductor macizo):
\begin{equation}
    W_{\text{int}} = \frac{1}{2\mu_0} \frac{\mu_0^2 I^2 (2\pi l)}{4\pi^2 a^4} \int_{0}^{a} s^3 \, ds = \frac{\mu_0 I^2 l}{4\pi a^4} \left[ \frac{s^4}{4} \right]_{0}^{a} = \frac{\mu_0 l}{16\pi} I^2
\end{equation}
Notemos el hecho notable de que la energía interna del alambre por unidad de longitud es independiente de su radio $a$.

Calculamos ahora la segunda integral (contribución de la región entre $a$ y $b$):
\begin{equation}
    W_{\text{entre}} = \frac{1}{2\mu_0} \frac{\mu_0^2 I^2 (2\pi l)}{4\pi^2} \int_{a}^{b} \frac{ds}{s} = \frac{\mu_0 l}{4\pi} \ln\left(\frac{b}{a}\right) I^2
\end{equation}
Sumando ambas contribuciones obtenemos la energía magnética total:
\begin{equation}
    W = W_{\text{int}} + W_{\text{entre}} = \frac{\mu_0 l}{4\pi} \left[ \frac{1}{4} + \ln\left(\frac{b}{a}\right) \right] I^2
\end{equation}
Igualando esta expresión a la definición energética de autoinductancia $W = \frac{1}{2} L I^2$:
\begin{equation}
    \frac{1}{2} L I^2 = \frac{\mu_0 l}{4\pi} \left[ \frac{1}{4} + \ln\left(\frac{b}{a}\right) \right] I^2 \implies L = \frac{\mu_0 l}{2\pi} \left[ \ln\left(\frac{b}{a}\right) + \frac{1}{4} \right]
\end{equation}
La autoinductancia por unidad de longitud $\mathcal{L} = L/l$ es:
\begin{equation}
    \mathcal{L} = \frac{\mu_0}{2\pi} \left[ \ln\left(\frac{b}{a}\right) + \frac{1}{4} \right]
\end{equation}
\end{solucion}

% -------------------------------------------------------------------------
\section{Corriente de Desplazamiento de Maxwell}

\begin{ejercicio}{Condensador con pérdidas (medio óhmico-dieléctrico) y corriente de desplazamiento}{condensador_perdidas}
Un condensador de placas paralelas de área $A$ y separación $d$ está relleno con un material imperfecto que posee simultáneamente una conductividad eléctrica $\sigma$ y una permitividad eléctrica $\varepsilon$. Las placas se conectan a un generador de corriente alterna que suministra una corriente total $I(t) = I_0 \cos(\omega t)$ hacia la placa superior.
\begin{enumerate}[label=\alph*)]
    \item Utilizando la ecuación de continuidad y la ley de Gauss, determine el campo eléctrico $\Evec(t)$ entre las placas en el régimen estacionario sinusoidal.
    \item Calcule explícitamente la densidad de corriente de conducción $\Jvec_c(t)$ y la densidad de corriente de desplazamiento $\Jvec_d(t)$ dentro del material, y verifique que la suma de ambas multiplicada por el área $A$ reproduce exactamente $I(t)$.
\end{enumerate}
\end{ejercicio}

\begin{solucion}
\textbf{Parte a) Determinación del campo eléctrico $\Evec(t)$ en el interior.}

Despreciando los efectos de borde ($A \gg d^2$), el campo eléctrico $\Evec(t) = E(t)\zhat$ es espacialmente uniforme en la región entre las placas. Dentro del material óhmico-dieléctrico lineal se cumplen las dos relaciones constitutivas:
\begin{equation}
    \Jvec_c(t) = \sigma \Evec(t), \qquad \Dvec(t) = \varepsilon \Evec(t)
\end{equation}
Sea $Q_f(t)$ la carga libre instantánea depositada sobre la cara interna de la placa superior (de modo que $\sigma_f(t) = Q_f(t)/A$). Por la condición de contorno para el desplazamiento eléctrico en la placa conductora superior:
\begin{equation}
    D(t) = \sigma_f(t) = \frac{Q_f(t)}{A} \implies \varepsilon E(t) = \frac{Q_f(t)}{A}
\end{equation}
Por conservación de la carga aplicada a una superficie cerrada que envuelve únicamente a la placa superior, la tasa de variación de la carga libre en la placa es igual a la corriente que entra por el cable alimentador ($I(t)$) menos la corriente de conducción que se fuga a través del material hacia la placa inferior ($I_c(t) = J_c(t) A$):
\begin{equation}
    \dv{Q_f(t)}{t} = I(t) - I_c(t) = I(t) - A \sigma E(t)
\end{equation}
Sustituyendo $Q_f(t) = A \varepsilon E(t)$ en el miembro izquierdo:
\begin{equation}
    A \varepsilon \dv{E(t)}{t} + A \sigma E(t) = I(t) = I_0 \cos(\omega t)
\end{equation}
Dividiendo por $A$, obtenemos una ecuación diferencial lineal de primer orden para $E(t)$:
\begin{equation}
    \varepsilon \dv{E}{t} + \sigma E = \frac{I_0}{A} \cos(\omega t) \label{eq:edo_maxwell_disp}
\end{equation}

Para hallar la solución estacionaria armónica, pasamos a notación fasorial compleja escribiendo $E(t) = \Re\big\{\tilde{E}_0 e^{i\omega t}\big\}$ y $\cos(\omega t) = \Re\big\{e^{i\omega t}\big\}$. Sustituyendo en \eqref{eq:edo_maxwell_disp}:
\begin{equation}
    (i\omega\varepsilon + \sigma) \tilde{E}_0 e^{i\omega t} = \frac{I_0}{A} e^{i\omega t} \implies \tilde{E}_0 = \frac{I_0}{A(\sigma + i\omega\varepsilon)}
\end{equation}
Multiplicamos numerador y denominador por el complejo conjugado $(\sigma - i\omega\varepsilon) = \sqrt{\sigma^2 + \omega^2\varepsilon^2}\,e^{-i\delta}$, donde el ángulo de desfase $\delta$ se define mediante $\tan\delta = \frac{\omega\varepsilon}{\sigma}$ (con $\cos\delta = \frac{\sigma}{\sqrt{\sigma^2+\omega^2\varepsilon^2}}$ y $\sin\delta = \frac{\omega\varepsilon}{\sqrt{\sigma^2+\omega^2\varepsilon^2}}$):
\begin{equation}
    \tilde{E}_0 = \frac{I_0}{A\sqrt{\sigma^2 + \omega^2\varepsilon^2}} e^{-i\delta}
\end{equation}
Tomando la parte real de $\tilde{E}_0 e^{i\omega t}$:
\begin{equation}
    \Evec(t) = \frac{I_0}{A\sqrt{\sigma^2 + \omega^2\varepsilon^2}} \cos(\omega t - \delta)\,\zhat, \qquad \text{con } \delta = \arctan\left(\frac{\omega\varepsilon}{\sigma}\right)
\end{equation}

\medskip
\textbf{Parte b) Cálculo de $\Jvec_c(t)$, $\Jvec_d(t)$ y comprobación de la ley de Ampère-Maxwell.}

Calculamos la densidad de corriente de conducción $\Jvec_c(t) = \sigma \Evec(t)$:
\begin{equation}
    \Jvec_c(t) = \frac{\sigma I_0}{A\sqrt{\sigma^2 + \omega^2\varepsilon^2}} \cos(\omega t - \delta)\,\zhat = \frac{I_0}{A} \cos\delta \cos(\omega t - \delta)\,\zhat
\end{equation}
Calculamos la densidad de corriente de desplazamiento de Maxwell en un medio material, $\Jvec_d(t) = \pdv{\Dvec}{t} = \varepsilon \pdv{\Evec}{t}$:
\begin{align}
    \Jvec_d(t) &= \varepsilon \dv{}{t}\left[ \frac{I_0}{A\sqrt{\sigma^2 + \omega^2\varepsilon^2}} \cos(\omega t - \delta) \right] \zhat \nonumber \\
    &= -\frac{\omega\varepsilon I_0}{A\sqrt{\sigma^2 + \omega^2\varepsilon^2}} \sin(\omega t - \delta)\,\zhat = -\frac{I_0}{A} \sin\delta \sin(\omega t - \delta)\,\zhat
\end{align}
Sumamos ambas densidades de corriente y multiplicamos por el área $A$:
\begin{equation}
    A\big(J_c(t) + J_d(t)\big) = I_0 \Big[ \cos\delta \cos(\omega t - \delta) - \sin\delta \sin(\omega t - \delta) \Big]
\end{equation}
Aplicando la identidad trigonométrica del coseno de la suma de dos ángulos, $\cos(\alpha + \beta) = \cos\alpha\cos\beta - \sin\alpha\sin\beta$, con $\alpha = \delta$ y $\beta = \omega t - \delta$:
\begin{equation}
    A\big(J_c(t) + J_d(t)\big) = I_0 \cos\big(\delta + (\omega t - \delta)\big) = I_0 \cos(\omega t) = I(t)
\end{equation}
Se verifica así que la corriente total (conducción más desplazamiento) es continua a través del condensador.
\end{solucion}

% =========================================================================
% CAPÍTULO 8: LEYES DE CONSERVACIÓN
% =========================================================================
\chapter[Leyes de conservación]{Leyes de Conservación (Teorema de Poynting, Tensor de Maxwell y Momento del Campo)}

\section{El Teorema de Poynting y el Flujo de Energía Electromagnética}

\begin{ejercicio}{Carga de un condensador circular y el vector de Poynting}{poynting_condensador}
Un condensador vacío de placas circulares paralelas de radio $a$ y separación $d$ ($d \ll a$) se está cargando mediante una corriente lenta $I(t)$ que entra por el eje de las placas, de modo que el campo eléctrico entre ellas aumenta linealmente o arbitrariamente con el tiempo.
\begin{enumerate}[label=\alph*)]
    \item Determine el campo eléctrico $\Evec(\bm{r},t)$ y el campo magnético inducido $\Bvec(\bm{r},t)$ en la región interior entre las placas ($s \le a$).
    \item Calcule el vector de Poynting $\Svec(\bm{r},t)$ en el borde cilíndrico del condensador ($s = a$) y demuestre integrando el flujo de $\Svec$ que la potencia que entra radialmente desde el espacio circundante hacia el volumen entre las placas es exactamente igual a la tasa temporal de incremento de la energía electrostática almacenada $\dv{W_e}{t}$.
\end{enumerate}
\end{ejercicio}

\begin{solucion}
\textbf{Parte a) Cálculo de $\Evec$ y $\Bvec$ entre las placas.}

Situemos el condensador centrado en el eje $z$, con placas de área $A = \pi a^2$. Si $Q(t)$ es la carga instantánea en la placa inferior y $-Q(t)$ en la superior (tal que $I(t) = \dv{Q}{t}$), el campo eléctrico uniforme entre las placas es:
\begin{equation}
    \Evec(t) = \frac{\sigma(t)}{\varepsilon_0}\zhat = \frac{Q(t)}{\pi \varepsilon_0 a^2}\zhat
\end{equation}
En el espacio vacío entre las placas no hay corriente de conducción ($\Jvec_c = \bm{0}$), pero existe una densidad de corriente de desplazamiento uniforme:
\begin{equation}
    \Jvec_d(t) = \varepsilon_0 \pdv{\Evec}{t} = \frac{1}{\pi a^2}\dv{Q}{t}\zhat = \frac{I(t)}{\pi a^2}\zhat
\end{equation}
Aplicamos la ley de Ampère-Maxwell a una circunferencia amperiana de radio $s \le a$ concéntrica con el eje $z$ entre las placas:
\begin{equation}
    \oint \Bvec \cdot d\bm{l} = \mu_0 \varepsilon_0 \int \pdv{\Evec}{t} \cdot d\bm{a} \implies B_\phi(s) (2\pi s) = \mu_0 \left(\frac{I(t)}{\pi a^2}\right) (\pi s^2)
\end{equation}
Despejando $\Bvec(s,t)$:
\begin{equation}
    \Bvec(s,t) = \frac{\mu_0 I(t) s}{2\pi a^2}\,\phihat = \frac{\mu_0 \varepsilon_0 s}{2} \dv{E(t)}{t}\,\phihat \qquad (s \le a)
\end{equation}

\medskip
\textbf{Parte b) Vector de Poynting y balance energético.}

El vector de Poynting en el vacío se define como:
\begin{equation}
    \Svec = \frac{1}{\mu_0} (\Evec \times \Bvec)
\end{equation}
Sustituyendo las expresiones de $\Evec$ y $\Bvec$ en cualquier punto del interior ($s \le a$) y recordando que en coordenadas cilíndricas $\zhat \times \phihat = -\shat$:
\begin{equation}
    \Svec(s,t) = \frac{1}{\mu_0} \left( E(t)\zhat \right) \times \left( \frac{\mu_0 \varepsilon_0 s}{2} \dv{E(t)}{t}\,\phihat \right) = -\frac{1}{2}\varepsilon_0 s E(t) \dv{E(t)}{t}\,\shat
\end{equation}
El signo $-\shat$ revela que la energía electromagnética fluye radialmente \textbf{hacia adentro} desde los bordes laterales del condensador hacia el eje central.

Calculemos la potencia total $P_{\text{in}}$ que atraviesa hacia adentro la superficie lateral cilíndrica del volumen entre las placas en $s = a$ (cuya área lateral es $2\pi a d$ y vector diferencial de área saliente es $d\bm{a} = da\,\shat$):
\begin{align}
    P_{\text{in}} &= -\oint_{s=a} \Svec(a,t) \cdot d\bm{a} = -\left( -\frac{1}{2}\varepsilon_0 a E(t) \dv{E(t)}{t}\,\shat \right) \cdot \big( 2\pi a d \,\shat \big) \nonumber \\
    &= \big(\pi a^2 d\big) \varepsilon_0 E(t) \dv{E(t)}{t}
\end{align}
Por otro lado, la energía electrostática total $W_e(t)$ almacenada en el volumen $\mathcal{V} = \pi a^2 d$ del condensador es:
\begin{equation}
    W_e(t) = \frac{1}{2}\varepsilon_0 \int_{\mathcal{V}} E^2(t)\,d\tau = \frac{1}{2}\varepsilon_0 E^2(t) \big(\pi a^2 d\big)
\end{equation}
Derivando $W_e(t)$ respecto al tiempo mediante la regla de la cadena:
\begin{equation}
    \dv{W_e(t)}{t} = \big(\pi a^2 d\big) \varepsilon_0 E(t) \dv{E(t)}{t}
\end{equation}
Comparando ambas expresiones, concluimos que $P_{\text{in}} = \dv{W_e}{t}$, verificando el teorema de Poynting.
\end{solucion}

% -------------------------------------------------------------------------
\section{El Tensor de Esfuerzos de Maxwell}

\begin{ejercicio}{Fuerza neta entre los dos hemisferios de una esfera maciza uniformemente cargada mediante el Tensor de Maxwell}{tensor_maxwell_esfera}
Una esfera maciza aislante de radio $R$ porta una carga total $Q$ distribuida uniformemente en todo su volumen (densidad $\rho = \frac{Q}{\frac{4}{3}\pi R^3}$). Utilice el tensor de esfuerzos de Maxwell $\overleftrightarrow{\bm{T}}$ para calcular la fuerza electrostática de repulsión que el hemisferio sur ($z < 0$) ejerce sobre el hemisferio norte ($z > 0$).
\end{ejercicio}

\begin{solucion}
La fuerza electromagnética total $\bm{F}$ que actúa sobre las cargas contenidas dentro de un volumen $\mathcal{V}$ delimitado por una superficie cerrada $\mathcal{S}$ (en ausencia de campos dependientes del tiempo y campos magnéticos, $\Bvec = \bm{0}$) viene dada por la integral de superficie del tensor de esfuerzos de Maxwell $\overleftrightarrow{\bm{T}}$:
\begin{equation}
    \bm{F} = \oint_{\mathcal{S}} \overleftrightarrow{\bm{T}} \cdot d\bm{a}
\end{equation}
donde las componentes cartesianas de $\overleftrightarrow{\bm{T}}$ para un campo puramente electrostático son:
\begin{equation}
    T_{ij} = \varepsilon_0 \left( E_i E_j - \frac{1}{2}\delta_{ij} E^2 \right)
\end{equation}
Por simetría axial alrededor del eje $z$, la fuerza neta sobre el hemisferio norte apunta puramente en la dirección $+\zhat$, es decir, $\bm{F} = F_z \zhat$, donde:
\begin{equation}
    F_z = \oint_{\mathcal{S}} \big( T_{zx} da_x + T_{zy} da_y + T_{zz} da_z \big)
\end{equation}

Elegimos como superficie cerrada $\mathcal{S}$ la frontera del hemisferio norte, la cual consta de dos partes:
\begin{enumerate}
    \item \textbf{$\mathcal{S}_1$: El casquete hemisférico superior} de radio $r = R$ con $0 \le \theta \le \pi/2$.
    \item \textbf{$\mathcal{S}_2$: El disco ecuatorial plano} en el plano $z = 0$ con radio $0 \le s \le R$.
\end{enumerate}
Necesitamos conocer el campo eléctrico $\Evec(\bm{r})$ en ambas superficies. Por la ley de Gauss aplicada a una esfera uniformemente cargada:
\begin{equation}
    \Evec(\bm{r}) = \begin{cases}
        \dfrac{Q}{4\pi\varepsilon_0 R^3} \bm{r} = \dfrac{Q r}{4\pi\varepsilon_0 R^3}\,\rhat, & \text{para } r \le R \\[10pt]
        \dfrac{Q}{4\pi\varepsilon_0 r^2}\,\rhat, & \text{para } r \ge R
    \end{cases}
\end{equation}

\textbf{Cálculo de la contribución $F_{z,1}$ sobre el casquete hemisférico $\mathcal{S}_1$ ($r = R$):}
En $r = R$, el vector diferencial de área saliente es $d\bm{a} = R^2\sin\theta\,d\theta\,d\phi\,\rhat$, cuyas componentes cartesianas son:
\begin{equation}
    da_x = R^2\sin^2\theta\cos\phi\,d\theta\,d\phi, \quad da_y = R^2\sin^2\theta\sin\phi\,d\theta\,d\phi, \quad da_z = R^2\sin\theta\cos\theta\,d\theta\,d\phi
\end{equation}
En $r = R$, la magnitud del campo eléctrico es $E = \frac{Q}{4\pi\varepsilon_0 R^2}$, y sus componentes cartesianas son:
\begin{equation}
    E_x = E\sin\theta\cos\phi, \qquad E_y = E\sin\theta\sin\phi, \qquad E_z = E\cos\theta
\end{equation}
Calculamos la contracción $(\overleftrightarrow{\bm{T}}\cdot d\bm{a})_z = T_{zx}da_x + T_{zy}da_y + T_{zz}da_z$:
\begin{align}
    (\overleftrightarrow{\bm{T}}\cdot d\bm{a})_z &= \varepsilon_0 \left[ E_z E_x da_x + E_z E_y da_y + \left(E_z^2 - \frac{1}{2}E^2\right) da_z \right] \nonumber \\
    &= \varepsilon_0 E^2 R^2\sin\theta\,d\theta\,d\phi \Big[ \cos\theta\sin^2\theta\cos^2\phi + \cos\theta\sin^2\theta\sin^2\phi + \left(\cos^2\theta - \frac{1}{2}\right)\cos\theta \Big] \nonumber \\
    &= \varepsilon_0 E^2 R^2\sin\theta\,d\theta\,d\phi \, \cos\theta \Big[ \sin^2\theta(\cos^2\phi + \sin^2\phi) + \cos^2\theta - \frac{1}{2} \Big] \nonumber \\
    &= \varepsilon_0 E^2 R^2\sin\theta\,d\theta\,d\phi \, \cos\theta \left[ 1 - \frac{1}{2} \right] = \frac{1}{2}\varepsilon_0 E^2 R^2 \sin\theta\cos\theta\,d\theta\,d\phi
\end{align}
Integramos sobre el hemisferio norte ($\phi \in [0,2\pi]$ y $\theta \in [0,\pi/2]$) sustituyendo $E = \frac{Q}{4\pi\varepsilon_0 R^2}$:
\begin{align}
    F_{z,1} &= \frac{1}{2}\varepsilon_0 \left(\frac{Q}{4\pi\varepsilon_0 R^2}\right)^2 R^2 \int_{0}^{2\pi} d\phi \int_{0}^{\pi/2} \sin\theta\cos\theta\,d\theta \nonumber \\
    &= \frac{1}{2}\varepsilon_0 \frac{Q^2}{16\pi^2\varepsilon_0^2 R^2} (2\pi) \left[ \frac{\sin^2\theta}{2} \right]_{0}^{\pi/2} = \frac{Q^2}{32\pi\varepsilon_0 R^2} = \frac{1}{8}\left(\frac{Q^2}{4\pi\varepsilon_0 R^2}\right)
\end{align}

\textbf{Cálculo de la contribución $F_{z,2}$ sobre el disco ecuatorial $\mathcal{S}_2$ ($z = 0, s \le R$):}
Sobre el disco plano en $z = 0$ que cierra por debajo al hemisferio norte, el vector unitario normal saliente apunta hacia abajo: $\nhat = -\zhat$. Por lo tanto:
\begin{equation}
    da_x = 0, \qquad da_y = 0, \qquad da_z = -s\,ds\,d\phi
\end{equation}
En el plano ecuatorial $z = 0$ ($\theta = \pi/2$), el campo eléctrico interno $\Evec = \frac{Q}{4\pi\varepsilon_0 R^3}(x\xhat + y\yhat + z\zhat)$ tiene componente $z$ idénticamente nula:
\begin{equation}
    E_z(z=0) = 0, \qquad E^2(s) = \left(\frac{Q s}{4\pi\varepsilon_0 R^3}\right)^2
\end{equation}
Por lo tanto, el único elemento del tensor que interviene es $T_{zz}$:
\begin{equation}
    T_{zz} = \varepsilon_0 \left( E_z^2 - \frac{1}{2}E^2 \right) = -\frac{1}{2}\varepsilon_0 \left(\frac{Q}{4\pi\varepsilon_0 R^3}\right)^2 s^2
\end{equation}
Multiplicando por $da_z = -s\,ds\,d\phi$ (los dos signos negativos dan positivo) e integrando sobre el disco de radio $R$:
\begin{align}
    F_{z,2} &= \int_{0}^{2\pi} d\phi \int_{0}^{R} T_{zz} da_z = \frac{1}{2}\varepsilon_0 \left(\frac{Q}{4\pi\varepsilon_0 R^3}\right)^2 (2\pi) \int_{0}^{R} s^3\,ds \nonumber \\
    &= \pi\varepsilon_0 \left(\frac{Q^2}{16\pi^2\varepsilon_0^2 R^6}\right) \left[\frac{R^4}{4}\right] = \frac{Q^2}{64\pi\varepsilon_0 R^2} = \frac{1}{16}\left(\frac{Q^2}{4\pi\varepsilon_0 R^2}\right)
\end{align}

\textbf{Fuerza total:}
Sumando ambas contribuciones superficiales:
\begin{equation}
    \bm{F} = (F_{z,1} + F_{z,2})\zhat = \left(\frac{1}{8} + \frac{1}{16}\right)\frac{Q^2}{4\pi\varepsilon_0 R^2}\zhat = \frac{3}{16} \frac{Q^2}{4\pi\varepsilon_0 R^2} \zhat
\end{equation}
\end{solucion}

% -------------------------------------------------------------------------
\section{Momento Lineal y Momento Angular del Campo Electromagnético}

\begin{ejercicio}{Paradoja de Feynman: Momento angular almacenado en campos estáticos cruzados e inducción al desactivarlos}{paradoja_feynman}
Considere un solenoide largo de radio $a$, longitud $l$ y $n$ vueltas por unidad de longitud por el que circula una corriente $I$ en sentido antihorario, situado coaxialmente en el interior de dos cascarones cilíndricos conductores largos de longitud $l$ y radios $R_1$ y $R_2$ (con $R_1 < a < R_2$). El cilindro interior ($R_1$) porta una carga $+Q$ y el exterior ($R_2$) porta una carga $-Q$.
\begin{enumerate}[label=\alph*)]
    \item Calcule el momento angular total $\bm{L}_{\text{em}}$ almacenado en los campos electromagnéticos estáticos cruzados $\Evec$ y $\Bvec$.
    \item Si la corriente $I$ del solenoide se reduce gradualmente hasta cero ($I \to 0$), determine el campo eléctrico azimutal inducido por la ley de Faraday en $s = R_1$ y $s = R_2$, calcule el impulso angular mecánico total entregado a los dos cilindros cargados y demuestre que coincide exactamente con el $\bm{L}_{\text{em}}$ inicial.
\end{enumerate}
\end{ejercicio}

\begin{solucion}
\textbf{Parte a) Cálculo del momento angular electromagnético inicial $\bm{L}_{\text{em}}$.}

Determinemos primero los campos estáticos iniciales en cada región radial $s$:
\begin{itemize}
    \item \textbf{Campo eléctrico $\Evec(s)$:} Por la ley de Gauss en simetría cilíndrica, existe campo eléctrico únicamente en la región comprendida entre los dos cilindros cargados ($R_1 < s < R_2$):
    \begin{equation}
        \Evec(s) = \begin{cases}
            \dfrac{Q}{2\pi\varepsilon_0 l s}\,\shat, & \text{para } R_1 < s < R_2 \\[8pt]
            \bm{0}, & \text{en el resto del espacio}
        \end{cases}
    \end{equation}
    \item \textbf{Campo magnético $\Bvec(s)$:} Por la ley de Ampère para un solenoide largo de radio $a$, existe campo magnético únicamente en el interior del solenoide ($s < a$):
    \begin{equation}
        \Bvec(s) = \begin{cases}
            \mu_0 n I\,\zhat, & \text{para } s < a \\[4pt]
            \bm{0}, & \text{para } s > a
        \end{cases}
    \end{equation}
\end{itemize}

La densidad volumétrica de momento lineal del campo electromagnético es $\bm{g} = \varepsilon_0 (\Evec \times \Bvec)$. Como $\Evec \neq \bm{0}$ en $(R_1, R_2)$ y $\Bvec \neq \bm{0}$ en $(0, a)$, ambos campos coexisten únicamente en la región anular superpuesta $R_1 < s < a$:
\begin{equation}
    \bm{g}(s) = \varepsilon_0 \left(\frac{Q}{2\pi\varepsilon_0 l s}\,\shat\right) \times \big(\mu_0 n I\,\zhat\big) = -\frac{\mu_0 n I Q}{2\pi l s}\,\phihat \qquad (R_1 < s < a)
\end{equation}
(recordando que $\shat \times \zhat = -\phihat$).

La densidad de momento angular respecto al origen es $\bm{\ell} = \bm{r} \times \bm{g}$. Escribiendo $\bm{r} = s\shat + z\zhat$:
\begin{equation}
    \bm{\ell} = (s\shat + z\zhat) \times \left(-\frac{\mu_0 n I Q}{2\pi l s}\,\phihat\right) = -\frac{\mu_0 n I Q}{2\pi l}\,\zhat + \frac{\mu_0 n I Q z}{2\pi l s}\,\shat
\end{equation}
Al integrar en el volumen cilíndrico con $d\tau = s\,ds\,d\phi\,dz$, la componente radial $\shat$ se anula por simetría azimutal al integrar en $\phi \in [0, 2\pi]$. Integrando la componente $z$ entre $s = R_1$ y $s = a$:
\begin{align}
    \bm{L}_{\text{em}} &= \int \bm{\ell}\,d\tau = \int_{0}^{l} dz \int_{0}^{2\pi} d\phi \int_{R_1}^{a} \left(-\frac{\mu_0 n I Q}{2\pi l}\,\zhat\right) s\,ds \nonumber \\
    &= -\frac{\mu_0 n I Q}{2\pi l} (l) (2\pi) \left[ \frac{s^2}{2} \right]_{R_1}^{a} \zhat = -\frac{1}{2}\mu_0 n I Q \big(a^2 - R_1^2\big)\,\zhat
\end{align}

\medskip
\textbf{Parte b) Inducción de Faraday al apagar la corriente e impulso angular mecánico.}

Cuando la corriente $I(t)$ varía con el tiempo desde su valor inicial $I$ hasta $0$, el campo magnético dentro del solenoide $B(t) = \mu_0 n I(t)$ cambia, induciendo un campo eléctrico azimutal $\Evec_{\text{ind}}(s,t) = E_\phi(s,t)\phihat$ dado por la ley de Faraday en forma integral $\oint \Evec_{\text{ind}} \cdot d\bm{l} = -\dv{\Phi_B}{t}$:
\begin{itemize}
    \item \textbf{En $s = R_1$ (dentro del solenoide, $R_1 < a$):} El flujo magnético encerrado por la circunferencia de radio $R_1$ es $\Phi_B(R_1) = (\pi R_1^2) \mu_0 n I(t)$. Por lo tanto:
    \begin{equation}
        E_\phi(R_1) (2\pi R_1) = -\pi R_1^2 \mu_0 n \dv{I}{t} \implies \Evec_{\text{ind}}(R_1) = -\frac{1}{2}\mu_0 n R_1 \dv{I}{t}\,\phihat
    \end{equation}
    Este campo eléctrico ejerce un torque mecánico $\bm{N}_1$ sobre el cilindro interior de radio $R_1$ y carga $+Q$:
    \begin{equation}
        \bm{N}_1 = (R_1\shat) \times \big(Q \Evec_{\text{ind}}(R_1)\big) = -\frac{1}{2}\mu_0 n Q R_1^2 \dv{I}{t}\,\zhat
    \end{equation}

    \item \textbf{En $s = R_2$ (fuera del solenoide, $R_2 > a$):} El flujo magnético encerrado por la circunferencia de radio $R_2$ está confinado al área del solenoide $\pi a^2$, es decir, $\Phi_B(R_2) = (\pi a^2) \mu_0 n I(t)$. Por lo tanto:
    \begin{equation}
        E_\phi(R_2) (2\pi R_2) = -\pi a^2 \mu_0 n \dv{I}{t} \implies \Evec_{\text{ind}}(R_2) = -\frac{\mu_0 n a^2}{2 R_2} \dv{I}{t}\,\phihat
    \end{equation}
    Este campo ejerce un torque mecánico $\bm{N}_2$ sobre el cilindro exterior de radio $R_2$ y carga $-Q$:
    \begin{equation}
        \bm{N}_2 = (R_2\shat) \times \big(-Q \Evec_{\text{ind}}(R_2)\big) = +\frac{1}{2}\mu_0 n Q a^2 \dv{I}{t}\,\zhat
    \end{equation}
\end{itemize}

El torque mecánico total sobre ambos cilindros es $\bm{N}_{\text{tot}} = \bm{N}_1 + \bm{N}_2$:
\begin{equation}
    \bm{N}_{\text{tot}}(t) = \frac{1}{2}\mu_0 n Q \big(a^2 - R_1^2\big) \dv{I}{t}\,\zhat
\end{equation}
Integrando el torque respecto al tiempo desde $t=0$ (con $I(0) = I$) hasta $t = \infty$ (con $I(\infty) = 0$) para obtener el momento angular mecánico final $\bm{L}_{\text{mec}}$ adquirido por los cilindros:
\begin{equation}
    \bm{L}_{\text{mec}} = \int_{0}^{\infty} \bm{N}_{\text{tot}}(t)\,dt = \frac{1}{2}\mu_0 n Q \big(a^2 - R_1^2\big)\,\zhat \int_{I}^{0} dI = -\frac{1}{2}\mu_0 n I Q \big(a^2 - R_1^2\big)\,\zhat
\end{equation}
Observamos que $\bm{L}_{\text{mec}} = \bm{L}_{\text{em}}$, demostrando que el momento angular total (campo + materia) se conserva exactamente.
\end{solucion}

% =========================================================================
% CAPÍTULO 9: ONDAS ELECTROMAGNÉTICAS
% =========================================================================
\chapter[Ondas electromagnéticas]{Ondas Electromagnéticas (Propagación, Reflexión, Medios Conductores y Guías de Onda)}

\section[Reflexión y transmisión en dieléctricos]{Reflexión y Transmisión en Medios Dieléctricos (Ecuaciones de Fresnel y Ángulo de Brewster)}

\begin{ejercicio}{Deducción de las Ecuaciones de Fresnel para polarización paralela al plano de incidencia y el Ángulo de Brewster}{fresnel_brewster}
Una onda electromagnética plana monocromática incide de forma oblicua con un ángulo $\theta_I$ sobre la interfaz plana ($z = 0$) entre dos medios dieléctricos lineales, homogéneos e isótropos de índices de refracción $n_1$ y $n_2$ (y permeabilidades magnéticas $\mu_1, \mu_2$). Suponga que el campo eléctrico está polarizado \textbf{paralelamente} al plano de incidencia $xz$ (polarización $p$ o TM).
\begin{enumerate}[label=\alph*)]
    \item Aplique las cuatro condiciones de contorno electrodinámicas en $z = 0$ para deducir las ecuaciones de Fresnel para las amplitudes complejas reflejada y transmitida ($E_{0R}/E_{0I}$ y $E_{0T}/E_{0I}$).
    \item Demuestre que cuando $\mu_1 = \mu_2 = \mu_0$, existe un ángulo de incidencia $\theta_B$ (ángulo de Brewster) para el cual la onda reflejada se anula por completo, y deduzca que $\tan\theta_B = n_2/n_1$.
\end{enumerate}
\end{ejercicio}

\begin{solucion}
\textbf{Parte a) Condiciones de frontera y ecuaciones de Fresnel.}

Situamos el plano de incidencia en el plano $xz$ y la interfaz en $z = 0$. Los vectores de onda incidente ($\bm{k}_I$), reflejado ($\bm{k}_R$) y transmitido ($\bm{k}_T$) son:
\begin{align}
    \bm{k}_I &= k_I (\sin\theta_I \xhat + \cos\theta_I \zhat) \\
    \bm{k}_R &= k_I (\sin\theta_I \xhat - \cos\theta_I \zhat) \\
    \bm{k}_T &= k_T (\sin\theta_T \xhat + \cos\theta_T \zhat)
\end{align}
donde por continuidad de la fase sobre el plano $z=0$ se cumplen la ley de reflexión ($\theta_R = \theta_I$) y la ley de Snell ($k_I\sin\theta_I = k_T\sin\theta_T \implies n_1\sin\theta_I = n_2\sin\theta_T$).

Como el campo eléctrico yace en el plano $xz$ y es transversal a $\bm{k}$ ($\Evec \cdot \bm{k} = 0$), mientras que el campo magnético $\Bvec = \frac{1}{v}(\hat{\bm{k}} \times \Evec)$ apunta en la dirección $\yhat$, escribimos los tres campos como:
\begin{align}
    \tilde{\Evec}_I(\bm{r},t) &= \tilde{E}_{0I} (\cos\theta_I \xhat - \sin\theta_I \zhat) e^{i(\bm{k}_I\cdot\bm{r} - \omega t)}, & \tilde{\Bvec}_I(\bm{r},t) &= \frac{1}{v_1}\tilde{E}_{0I}\,\yhat\,e^{i(\bm{k}_I\cdot\bm{r} - \omega t)} \\
    \tilde{\Evec}_R(\bm{r},t) &= \tilde{E}_{0R} (\cos\theta_I \xhat + \sin\theta_I \zhat) e^{i(\bm{k}_R\cdot\bm{r} - \omega t)}, & \tilde{\Bvec}_R(\bm{r},t) &= -\frac{1}{v_1}\tilde{E}_{0R}\,\yhat\,e^{i(\bm{k}_R\cdot\bm{r} - \omega t)} \\
    \tilde{\Evec}_T(\bm{r},t) &= \tilde{E}_{0T} (\cos\theta_T \xhat - \sin\theta_T \zhat) e^{i(\bm{k}_T\cdot\bm{r} - \omega t)}, & \tilde{\Bvec}_T(\bm{r},t) &= \frac{1}{v_2}\tilde{E}_{0T}\,\yhat\,e^{i(\bm{k}_T\cdot\bm{r} - \omega t)}
\end{align}

Imponemos las condiciones de contorno en $z = 0$ (donde los factores exponenciales son idénticos y se cancelan):
\begin{enumerate}[label=\roman*)]
    \item \textbf{Continuidad de $D_\perp = \varepsilon E_z$:}
    \begin{equation}
        \varepsilon_1\big(-\tilde{E}_{0I}\sin\theta_I + \tilde{E}_{0R}\sin\theta_I\big) = \varepsilon_2\big(-\tilde{E}_{0T}\sin\theta_T\big)
    \end{equation}
    \item \textbf{Continuidad de $E_\parallel = E_x$:}
    \begin{equation}
        \big(\tilde{E}_{0I} + \tilde{E}_{0R}\big)\cos\theta_I = \tilde{E}_{0T}\cos\theta_T \implies \tilde{E}_{0I} + \tilde{E}_{0R} = \alpha \tilde{E}_{0T} \label{eq:fresnel_1}
    \end{equation}
    donde hemos definido el parámetro geométrico $\alpha \equiv \frac{\cos\theta_T}{\cos\theta_I}$.
    \item \textbf{Continuidad de $B_\perp = B_z$:} Se satisface trivialmente ($0 = 0$).
    \item \textbf{Continuidad de $H_\parallel = \frac{1}{\mu}B_y$:}
    \begin{equation}
        \frac{1}{\mu_1 v_1}\big(\tilde{E}_{0I} - \tilde{E}_{0R}\big) = \frac{1}{\mu_2 v_2}\tilde{E}_{0T} \implies \tilde{E}_{0I} - \tilde{E}_{0R} = \beta \tilde{E}_{0T} \label{eq:fresnel_2}
    \end{equation}
    donde hemos definido el parámetro material $\beta \equiv \frac{\mu_1 v_1}{\mu_2 v_2} = \frac{\mu_1 n_2}{\mu_2 n_1}$. (Nótese que la condición i sobre $D_\perp$ es equivalente a iv usando la ley de Snell y $v = 1/\sqrt{\varepsilon\mu}$).
\end{enumerate}

Resolvemos el sistema lineal formado por \eqref{eq:fresnel_1} y \eqref{eq:fresnel_2}:
Sumando ambas ecuaciones:
\begin{equation}
    2\tilde{E}_{0I} = (\alpha + \beta)\tilde{E}_{0T} \implies \left(\frac{\tilde{E}_{0T}}{\tilde{E}_{0I}}\right) = \frac{2}{\alpha + \beta} = \frac{2\cos\theta_I}{\cos\theta_T + \frac{\mu_1 n_2}{\mu_2 n_1}\cos\theta_I}
\end{equation}
Restando la ecuación \eqref{eq:fresnel_2} de la \eqref{eq:fresnel_1}:
\begin{equation}
    2\tilde{E}_{0R} = (\alpha - \beta)\tilde{E}_{0T} = (\alpha - \beta)\left(\frac{2}{\alpha + \beta}\right)\tilde{E}_{0I} \implies \left(\frac{\tilde{E}_{0R}}{\tilde{E}_{0I}}\right) = \frac{\alpha - \beta}{\alpha + \beta} = \frac{\cos\theta_T - \frac{\mu_1 n_2}{\mu_2 n_1}\cos\theta_I}{\cos\theta_T + \frac{\mu_1 n_2}{\mu_2 n_1}\cos\theta_I}
\end{equation}

\medskip
\textbf{Parte b) Deducción del ángulo de Brewster.}

Para medios no magnéticos ($\mu_1 = \mu_2 = \mu_0$), tenemos $\beta = \frac{n_2}{n_1}$. La amplitud reflejada $\tilde{E}_{0R}$ se hace exactamente cero cuando el numerador se anula, es decir, cuando $\alpha = \beta$:
\begin{equation}
    \frac{\cos\theta_T}{\cos\theta_B} = \frac{n_2}{n_1} \implies \cos\theta_T = \frac{n_2}{n_1}\cos\theta_B
\end{equation}
Elevamos ambos miembros al cuadrado y usamos la identidad trigonométrica $\cos^2\theta_T = 1 - \sin^2\theta_T$:
\begin{equation}
    1 - \sin^2\theta_T = \left(\frac{n_2}{n_1}\right)^2 \cos^2\theta_B
\end{equation}
Por la ley de Snell, $\sin\theta_T = \frac{n_1}{n_2}\sin\theta_B$. Sustituyendo esto en el miembro izquierdo:
\begin{equation}
    1 - \left(\frac{n_1}{n_2}\right)^2 \sin^2\theta_B = \left(\frac{n_2}{n_1}\right)^2 \big(1 - \sin^2\theta_B\big)
\end{equation}
Agrupamos los términos que contienen $\sin^2\theta_B$:
\begin{equation}
    1 - \left(\frac{n_2}{n_1}\right)^2 = \left[ \left(\frac{n_1}{n_2}\right)^2 - \left(\frac{n_2}{n_1}\right)^2 \right] \sin^2\theta_B
\end{equation}
Factorizamos el corchete poniendo común denominador:
\begin{equation}
    \frac{n_1^2 - n_2^2}{n_1^2} = \left( \frac{n_1^4 - n_2^4}{n_1^2 n_2^2} \right) \sin^2\theta_B = \frac{(n_1^2 - n_2^2)(n_1^2 + n_2^2)}{n_1^2 n_2^2} \sin^2\theta_B
\end{equation}
Cancelando el factor $\frac{n_1^2 - n_2^2}{n_1^2}$ (suponiendo $n_1 \neq n_2$):
\begin{equation}
    1 = \frac{n_1^2 + n_2^2}{n_2^2} \sin^2\theta_B \implies \sin^2\theta_B = \frac{n_2^2}{n_1^2 + n_2^2}
\end{equation}
A partir de aquí calculamos $\cos^2\theta_B$:
\begin{equation}
    \cos^2\theta_B = 1 - \sin^2\theta_B = 1 - \frac{n_2^2}{n_1^2 + n_2^2} = \frac{n_1^2}{n_1^2 + n_2^2}
\end{equation}
Dividiendo $\sin\theta_B$ entre $\cos\theta_B$:
\begin{equation}
    \tan\theta_B = \frac{\sin\theta_B}{\cos\theta_B} = \frac{n_2}{n_1}
\end{equation}
Adicionalmente, notamos que en este ángulo $\sin\theta_T = \frac{n_1}{n_2}\sin\theta_B = \cos\theta_B$, lo cual significa que $\theta_B + \theta_T = \frac{\pi}{2}$ (los rayos reflejado y refractado forman un ángulo recto entre sí).
\end{solucion}

% -------------------------------------------------------------------------
\section{Ondas Electromagnéticas en Conductores y Efecto Pelicular (Skin Depth)}

\begin{ejercicio}{Propagación, atenuación y desfase de ondas en un buen conductor}{ondas_conductor}
Considere una onda electromagnética plana monocromática de frecuencia angular $\omega$ que se propaga a lo largo del eje $z$ dentro de un medio conductor óhmico lineal con parámetros $\varepsilon$, $\mu$ y conductividad $\sigma$.
\begin{enumerate}[label=\alph*)]
    \item A partir de las ecuaciones de Maxwell con $\rho_f = 0$ y $\Jvec_f = \sigma\Evec$, deduzca la ecuación de onda amortiguada (ecuación del telegrafista) y demuestre que el número de onda complejo $\tilde{k} = k + i\kappa$ satisface $\tilde{k}^2 = \mu\varepsilon\omega^2 + i\mu\sigma\omega$.
    \item Resuelva algebraicamente para obtener las partes real ($k$) e imaginaria ($\kappa$) exactas, evalúe el límite de un buen conductor ($\sigma \gg \varepsilon\omega$), demuestre que el campo magnético $\Bvec$ está retrasado en fase $\phi = 45^\circ$ ($\pi/4$ rad) respecto al campo eléctrico $\Evec$, y calcule la razón entre las densidades medias de energía magnética y eléctrica.
\end{enumerate}
\end{ejercicio}

\begin{solucion}
\textbf{Parte a) Deducción de la ecuación de onda y relación de dispersión compleja.}

Dentro del conductor homogéneo, cualquier exceso inicial de carga libre se disipa exponencialmente en un tiempo característico de relajación $\tau = \varepsilon/\sigma$ extremadamente corto, por lo que en régimen estacionario $\rho_f = 0$. Las cuatro ecuaciones de Maxwell junto con la ley de Ohm $\Jvec_f = \sigma\Evec$ toman la forma:
\begin{equation}
    \div{\Evec} = 0, \qquad \div{\Bvec} = 0, \qquad \curl{\Evec} = -\pdv{\Bvec}{t}, \qquad \curl{\Bvec} = \mu\sigma\Evec + \mu\varepsilon\pdv{\Evec}{t}
\end{equation}
Aplicamos el rotacional a la ley de Faraday y usamos la identidad vectorial $\curl(\curl{\Evec}) = \grad(\div{\Evec}) - \nabla^2\Evec$:
\begin{equation}
    -\nabla^2\Evec = -\pdv{}{t}(\curl{\Bvec}) = -\pdv{}{t}\left( \mu\sigma\Evec + \mu\varepsilon\pdv{\Evec}{t} \right)
\end{equation}
Reordenando términos, obtenemos la ecuación de onda con disipación:
\begin{equation}
    \nabla^2\Evec = \mu\varepsilon\pdv[2]{\Evec}{t} + \mu\sigma\pdv{\Evec}{t}
\end{equation}
Proponemos una solución de onda plana monocromática que viaja en $+z$:
\begin{equation}
    \tilde{\Evec}(z,t) = \tilde{\Evec}_0 e^{i(\tilde{k}z - \omega t)}
\end{equation}
Al sustituir las derivadas $\pdv[2]{}{z} \to (i\tilde{k})^2 = -\tilde{k}^2$, $\pdv{}{t} \to -i\omega$ y $\pdv[2]{}{t} \to -\omega^2$:
\begin{equation}
    -\tilde{k}^2 \tilde{\Evec} = -\mu\varepsilon\omega^2 \tilde{\Evec} - i\mu\sigma\omega \tilde{\Evec} \implies \tilde{k}^2 = \mu\varepsilon\omega^2 + i\mu\sigma\omega
\end{equation}

\medskip
\textbf{Parte b) Separación en parte real e imaginaria y régimen de buen conductor.}

Escribimos $\tilde{k} = k + i\kappa$ (con $k, \kappa \in \mathbb{R}^+$) y elevamos al cuadrado:
\begin{equation}
    (k + i\kappa)^2 = \big(k^2 - \kappa^2\big) + 2ik\kappa = \mu\varepsilon\omega^2 + i\mu\sigma\omega
\end{equation}
Igualando las partes reales e imaginarias de ambos miembros, además del módulo $|\tilde{k}^2| = k^2 + \kappa^2$:
\begin{align}
    k^2 - \kappa^2 &= \mu\varepsilon\omega^2 \label{eq:cond_re} \\
    k^2 + \kappa^2 &= \sqrt{(\mu\varepsilon\omega^2)^2 + (\mu\sigma\omega)^2} = \mu\varepsilon\omega^2 \sqrt{1 + \left(\frac{\sigma}{\varepsilon\omega}\right)^2} \label{eq:cond_mod}
\end{align}
Sumando y restando \eqref{eq:cond_mod} y \eqref{eq:cond_re} despejamos exactamente $k$ y $\kappa$:
\begin{equation}
    k = \omega\sqrt{\frac{\varepsilon\mu}{2}}\left[ \sqrt{1 + \left(\frac{\sigma}{\varepsilon\omega}\right)^2} + 1 \right]^{1/2}, \qquad \kappa = \omega\sqrt{\frac{\varepsilon\mu}{2}}\left[ \sqrt{1 + \left(\frac{\sigma}{\varepsilon\omega}\right)^2} - 1 \right]^{1/2}
\end{equation}

\textbf{Límite de un buen conductor ($\sigma \gg \varepsilon\omega$):}
En este régimen $\frac{\sigma}{\varepsilon\omega} \gg 1$, por lo que $\sqrt{1 + \left(\frac{\sigma}{\varepsilon\omega}\right)^2} \approx \frac{\sigma}{\varepsilon\omega} \gg 1$. Despreciando el $\pm 1$ frente a $\frac{\sigma}{\varepsilon\omega}$:
\begin{equation}
    k \approx \kappa \approx \omega\sqrt{\frac{\varepsilon\mu}{2}\left(\frac{\sigma}{\varepsilon\omega}\right)} = \sqrt{\frac{\omega\mu\sigma}{2}}
\end{equation}
La profundidad de penetración o espesor pelicular (*skin depth*) $d \equiv \frac{1}{\kappa}$ es por tanto $d = \sqrt{\frac{2}{\omega\mu\sigma}}$.

En forma polar compleja, como $k = \kappa$:
\begin{equation}
    \tilde{k} = k + i\kappa \approx \sqrt{\frac{\omega\mu\sigma}{2}}(1 + i) = \sqrt{\omega\mu\sigma}\,e^{i\pi/4}
\end{equation}
Por la ley de Faraday para una onda plana polarizada en $\xhat$ ($\tilde{\Evec} = \tilde{E}_0 e^{-\kappa z} e^{i(kz - \omega t)}\xhat$):
\begin{equation}
    \curl{\tilde{\Evec}} = -\pdv{\tilde{\Bvec}}{t} \implies i\tilde{k}\tilde{E}_0\yhat = i\omega\tilde{B}_0\yhat \implies \tilde{B}_0 = \frac{\tilde{k}}{\omega}\tilde{E}_0 \approx \sqrt{\frac{\mu\sigma}{\omega}} e^{i\pi/4} \tilde{E}_0
\end{equation}
Si elegimos el origen de tiempos de modo que $\tilde{E}_0 = E_0 \in \mathbb{R}$, las partes reales físicas de los campos son:
\begin{align}
    \Evec(z,t) &= E_0 e^{-\kappa z} \cos(kz - \omega t)\,\xhat \\
    \Bvec(z,t) &= E_0 \sqrt{\frac{\mu\sigma}{\omega}} e^{-\kappa z} \cos\left(kz - \omega t + \frac{\pi}{4}\right)\yhat
\end{align}
Como la fase temporal es $-\omega t$, un término $+\pi/4$ dentro del argumento $(kz - \omega t + \pi/4) = kz - \omega(t - \frac{\pi}{4\omega})$ significa que los máximos de $\Bvec$ ocurren un tiempo $\Delta t = \frac{\pi}{4\omega}$ más tarde que los de $\Evec$, es decir, $\Bvec$ está retrasado $45^\circ$ respecto a $\Evec$.

Finalmente, calculemos el cociente entre el promedio temporal de la densidad de energía magnética $\langle u_{\text{mag}} \rangle = \frac{1}{2\mu}\langle B^2 \rangle = \frac{1}{4\mu}|\tilde{B}_0|^2 e^{-2\kappa z}$ y la densidad de energía eléctrica $\langle u_{\text{elec}} \rangle = \frac{1}{2}\varepsilon\langle E^2 \rangle = \frac{1}{4}\varepsilon|\tilde{E}_0|^2 e^{-2\kappa z}$:
\begin{equation}
    \frac{\langle u_{\text{mag}} \rangle}{\langle u_{\text{elec}} \rangle} = \frac{\frac{1}{4\mu} \left(\frac{\mu\sigma}{\omega}\right) E_0^2 e^{-2\kappa z}}{\frac{1}{4}\varepsilon E_0^2 e^{-2\kappa z}} = \frac{\sigma}{\varepsilon\omega} \gg 1
\end{equation}
En un buen conductor, la energía de la onda electromagnética reside casi en su totalidad en el campo magnético.
\end{solucion}

% -------------------------------------------------------------------------
\section{Guías de Onda Metálicas}

\begin{ejercicio}{Modos Transversales Eléctricos ($\text{TE}_{mn}$), Frecuencias de Corte y Velocidad de Grupo en una Guía Rectangular}{guia_onda_rectangular}
Una guía de onda metálica hueca tiene sección transversal rectangular de dimensiones $a \times b$ ($0 \le x \le a$, $0 \le y \le b$, con $a > b$) y paredes perfectamente conductoras.
\begin{enumerate}[label=\alph*)]
    \item Resuelva la ecuación de Helmholtz bidimensional para la componente longitudinal $B_z(x,y)$ de un modo Transversal Eléctrico (TE, donde $E_z = 0$), aplicando las condiciones de contorno en las cuatro paredes para determinar los modos $\text{TE}_{mn}$ y sus frecuencias de corte $\omega_{mn}$.
    \item Demuestre que la velocidad de fase $v_p = \omega/k$ supera a $c$, pero la velocidad de grupo $v_g = d\omega/dk$ (velocidad a la que viajan la energía y la información) es siempre menor que $c$ y cumple $v_p v_g = c^2$.
\end{enumerate}
\end{ejercicio}

\begin{solucion}
\textbf{Parte a) Resolución para los modos $\text{TE}_{mn}$ y frecuencias de corte.}

Para una onda guiada a lo largo del eje $z$ de la forma $\tilde{\Evec}(x,y,z,t) = \Evec_0(x,y)e^{i(kz-\omega t)}$ y $\tilde{\Bvec}(x,y,z,t) = \Bvec_0(x,y)e^{i(kz-\omega t)}$, todas las componentes transversales ($E_x, E_y, B_x, B_y$) quedan determinadas por $E_z$ y $B_z$. Para modos TE ($E_z = 0$), las ecuaciones de Maxwell relacionan las componentes transversales con las derivadas de $B_z$:
\begin{align}
    E_x &= \frac{i\omega}{(\omega/c)^2 - k^2}\pdv{B_z}{y}, & E_y &= -\frac{i\omega}{(\omega/c)^2 - k^2}\pdv{B_z}{x} \label{eq:te_trans_e} \\
    B_x &= \frac{ik}{(\omega/c)^2 - k^2}\pdv{B_z}{x}, & B_y &= \frac{ik}{(\omega/c)^2 - k^2}\pdv{B_z}{y}
\end{align}
mientras que $B_z(x,y)$ obedece la ecuación de Helmholtz transversal:
\begin{equation}
    \left[ \pdv[2]{}{x} + \pdv[2]{}{y} + \left(\frac{\omega}{c}\right)^2 - k^2 \right] B_z(x,y) = 0
\end{equation}
Proponemos separación de variables $B_z(x,y) = X(x)Y(y)$. Dividiendo por $XY$:
\begin{equation}
    \frac{1}{X}\dv[2]{X}{x} + \frac{1}{Y}\dv[2]{Y}{y} = k^2 - \left(\frac{\omega}{c}\right)^2
\end{equation}
Llamando $-k_x^2$ al primer término y $-k_y^2$ al segundo, con:
\begin{equation}
    k_x^2 + k_y^2 = \left(\frac{\omega}{c}\right)^2 - k^2 \label{eq:disp_guia}
\end{equation}
las soluciones generales son:
\begin{equation}
    X(x) = A\sin(k_x x) + B\cos(k_x x), \qquad Y(y) = C\sin(k_y y) + D\cos(k_y y)
\end{equation}

En las paredes metálicas perfectamente conductoras, la componente tangencial de $\Evec$ y la componente normal de $\Bvec$ deben ser nulas ($\Evec_\parallel = \bm{0}$, $B_\perp = 0$). De las ecuaciones \eqref{eq:te_trans_e}, vemos que $E_y \propto \pdv{B_z}{x}$ y $B_x \propto \pdv{B_z}{x}$, mientras que $E_x \propto \pdv{B_z}{y}$ y $B_y \propto \pdv{B_z}{y}$. Por tanto, las condiciones de contorno son condiciones de Neumann homogéneas:
\begin{enumerate}
    \item En $x = 0$ y $x = a$, debe cumplirse $E_y = 0$ y $B_x = 0 \implies \left.\pdv{B_z}{x}\right|_{x=0,a} = 0$:
    \begin{equation}
        X'(0) = k_x A = 0 \implies A = 0
    \end{equation}
    \begin{equation}
        X'(a) = -k_x B \sin(k_x a) = 0 \implies k_x = \frac{m\pi}{a}, \qquad (m = 0, 1, 2, \dots)
    \end{equation}
    \item En $y = 0$ y $y = b$, debe cumplirse $E_x = 0$ y $B_y = 0 \implies \left.\pdv{B_z}{y}\right|_{y=0,b} = 0$:
    \begin{equation}
        Y'(0) = k_y C = 0 \implies C = 0
    \end{equation}
    \begin{equation}
        Y'(b) = -k_y D \sin(k_y b) = 0 \implies k_y = \frac{n\pi}{b}, \qquad (n = 0, 1, 2, \dots)
    \end{equation}
\end{enumerate}
Así, el perfil espacial del modo $\text{TE}_{mn}$ es:
\begin{equation}
    B_z(x,y) = B_0 \cos\left(\frac{m\pi x}{a}\right)\cos\left(\frac{n\pi y}{b}\right)
\end{equation}
(donde $m$ y $n$ no pueden ser simultáneamente cero, pues si $m=n=0$, $B_z = \text{cte}$ y por \eqref{eq:te_trans_e} todos los campos transversales se anulan y $\div{\Bvec} = ikB_0 = 0 \implies B_0 = 0$).

Despejando el número de onda longitudinal $k$ de \eqref{eq:disp_guia}:
\begin{equation}
    k = \sqrt{\left(\frac{\omega}{c}\right)^2 - \pi^2\left[\left(\frac{m}{a}\right)^2 + \left(\frac{n}{b}\right)^2\right]} = \frac{1}{c}\sqrt{\omega^2 - \omega_{mn}^2}
\end{equation}
donde hemos definido la frecuencia de corte $\omega_{mn}$ del modo $\text{TE}_{mn}$ como:
\begin{equation}
    \omega_{mn} = c\pi\sqrt{\left(\frac{m}{a}\right)^2 + \left(\frac{n}{b}\right)^2}
\end{equation}
Si $\omega < \omega_{mn}$, $k$ se vuelve puramente imaginario y la onda es evanescente (no se propaga). Como $a > b$, el modo fundamental de menor frecuencia de corte es el $\text{TE}_{10}$, con $\omega_{10} = \frac{c\pi}{a}$.

\medskip
\textbf{Parte b) Velocidad de fase y velocidad de grupo.}

A partir de la relación de dispersión $\omega(k) = \sqrt{c^2 k^2 + \omega_{mn}^2}$, calculamos la velocidad de fase $v_p$:
\begin{equation}
    v_p = \frac{\omega}{k} = \frac{\omega}{\frac{1}{c}\sqrt{\omega^2 - \omega_{mn}^2}} = \frac{c}{\sqrt{1 - (\omega_{mn}/\omega)^2}} > c
\end{equation}
Calculamos ahora la velocidad de grupo $v_g = \dv{\omega}{k}$ derivando implícitamente $\omega^2 = c^2 k^2 + \omega_{mn}^2$:
\begin{equation}
    2\omega\,d\omega = 2c^2 k\,dk \implies v_g = \dv{\omega}{k} = \frac{c^2 k}{\omega} = c\sqrt{1 - \left(\frac{\omega_{mn}}{\omega}\right)^2} < c
\end{equation}
Multiplicando ambas velocidades:
\begin{equation}
    v_p v_g = \left(\frac{\omega}{k}\right)\left(\frac{c^2 k}{\omega}\right) = c^2
\end{equation}
\end{solucion}

% =========================================================================
% CAPÍTULO 10: POTENCIALES Y CAMPOS
% =========================================================================
\chapter{Potenciales Dinámicos y Campos de Cargas en Movimiento}

\section[Gauge y potenciales retardados]{Transformaciones de Gauge, Potenciales Retardados y Ecuaciones de Jefimenko}

\begin{ejercicio}{Potenciales retardados de un alambre infinito cuya corriente se enciende súbitamente}{alambre_retardado}
Un alambre recto infinitamente largo situado a lo largo del eje $z$ es eléctricamente neutro en todo instante ($\lambda = 0$). Para $t < 0$ no circula corriente por él, pero en el instante $t = 0$ se conecta de forma instantánea una corriente constante $I_0$, de modo que $I(t) = I_0 \Theta(t)$ (donde $\Theta(t)$ es la función escalón de Heaviside).
\begin{enumerate}[label=\alph*)]
    \item Calcule el potencial escalar retardado $V(\bm{r},t)$ y el potencial vectorial retardado $\Avec(\bm{r},t)$ a una distancia radial $s$ del alambre para cualquier instante $t$.
    \item A partir de los potenciales hallados en (a), determine los campos $\Evec(s,t)$ y $\Bvec(s,t)$.
\end{enumerate}
\end{ejercicio}

\begin{solucion}
\textbf{Parte a) Cálculo de los potenciales retardados $V$ y $\Avec$.}

Dado que el alambre es eléctricamente neutro en todo momento ($\rho(\bm{r}',t_r) = 0$), el potencial escalar retardado en el gauge de Lorenz es idénticamente nulo:
\begin{equation}
    V(s,t) = \frac{1}{4\pi\varepsilon_0} \int \frac{\rho(\bm{r}', t_r)}{\mathscr{r}} d\tau' = 0
\end{equation}
Para el potencial vectorial retardado de un alambre filiforme a lo largo del eje $z$:
\begin{equation}
    \Avec(s,t) = \frac{\mu_0}{4\pi}\zhat \int_{-\infty}^{\infty} \frac{I(t_r)}{\mathscr{r}} dz'
\end{equation}
En coordenadas cilíndricas, para un punto de observación a distancia perpendicular $s$ del eje $z$ (podemos elegir $z = 0$ por invariancia traslacional) y un elemento del alambre en $z'$, la distancia es:
\begin{equation}
    \mathscr{r} = \sqrt{s^2 + z'^2}
\end{equation}
El tiempo retardado correspondiente a ese elemento es:
\begin{equation}
    t_r = t - \frac{\mathscr{r}}{c} = t - \frac{\sqrt{s^2 + z'^2}}{c}
\end{equation}
Como la corriente $I(t_r) = I_0 \Theta(t_r)$ es distinta de cero únicamente cuando el argumento $t_r \ge 0$, la condición para que un segmento $z'$ del alambre contribuya en el instante $t$ es:
\begin{equation}
    t - \frac{\sqrt{s^2 + z'^2}}{c} \ge 0 \implies \sqrt{s^2 + z'^2} \le ct \implies z'^2 \le (ct)^2 - s^2
\end{equation}
De aquí se desprenden dos casos físicos fundamentales:
\begin{enumerate}
    \item \textbf{Si $t < \frac{s}{c}$ (es decir, $ct < s$):} Ningún punto del alambre ha tenido tiempo aún de enviar su señal luminosa hasta la distancia $s$. Por tanto:
    \begin{equation}
        \Avec(s,t) = \bm{0} \qquad \left(t \le \frac{s}{c}\right)
    \end{equation}
    \item \textbf{Si $t > \frac{s}{c}$ (es decir, $ct > s$):} Contribuye el tramo simétrico del alambre comprendido entre $-z_{\text{máx}}$ y $+z_{\text{máx}}$, donde $z_{\text{máx}} = \sqrt{(ct)^2 - s^2}$:
    \begin{equation}
        \Avec(s,t) = \frac{\mu_0 I_0}{4\pi}\zhat \int_{-\sqrt{(ct)^2 - s^2}}^{+\sqrt{(ct)^2 - s^2}} \frac{dz'}{\sqrt{s^2 + z'^2}} = \frac{\mu_0 I_0}{2\pi}\zhat \int_{0}^{\sqrt{(ct)^2 - s^2}} \frac{dz'}{\sqrt{s^2 + z'^2}}
    \end{equation}
\end{enumerate}
Para resolver esta integral, usamos la sustitución hiperbólica $z' = s \sinh u \implies dz' = s \cosh u\,du$ y $\sqrt{s^2 + z'^2} = s \cosh u$, cuya primitiva es $\ln\big(z' + \sqrt{s^2 + z'^2}\big)$:
\begin{align}
    \int_{0}^{\sqrt{(ct)^2 - s^2}} \frac{dz'}{\sqrt{s^2 + z'^2}} &= \left[ \ln\left(z' + \sqrt{s^2 + z'^2}\right) \right]_{0}^{\sqrt{(ct)^2 - s^2}} \nonumber \\
    &= \ln\left( \sqrt{(ct)^2 - s^2} + ct \right) - \ln(s) = \ln\left( \frac{ct + \sqrt{(ct)^2 - s^2}}{s} \right)
\end{align}
Por consiguiente, para $t > s/c$:
\begin{equation}
    \Avec(s,t) = \frac{\mu_0 I_0}{2\pi} \ln\left( \frac{ct + \sqrt{(ct)^2 - s^2}}{s} \right) \zhat \qquad \left(t > \frac{s}{c}\right)
\end{equation}

\medskip
\textbf{Parte b) Cálculo de los campos eléctrico $\Evec$ y magnético $\Bvec$.}

Para $t > s/c$, calculamos el campo eléctrico mediante $\Evec = -\grad V - \pdv{\Avec}{t} = -\pdv{\Avec}{t}$:
\begin{align}
    \Evec(s,t) &= -\pdv{}{t}\left[ \frac{\mu_0 I_0}{2\pi} \ln\left( \frac{ct + \sqrt{(ct)^2 - s^2}}{s} \right) \right] \zhat \nonumber \\
    &= -\frac{\mu_0 I_0}{2\pi} \left( \frac{1}{ct + \sqrt{(ct)^2 - s^2}} \right) \left( c + \frac{c^2 t}{\sqrt{(ct)^2 - s^2}} \right) \zhat \nonumber \\
    &= -\frac{\mu_0 I_0}{2\pi} \left( \frac{1}{ct + \sqrt{(ct)^2 - s^2}} \right) c \left( \frac{\sqrt{(ct)^2 - s^2} + ct}{\sqrt{(ct)^2 - s^2}} \right) \zhat \nonumber \\
    &= -\frac{\mu_0 I_0 c}{2\pi \sqrt{(ct)^2 - s^2}}\,\zhat
\end{align}
Nótese que cuando $t \to \infty$, $\Evec(s,t) \to \bm{0}$, recuperando la ausencia de campo eléctrico en magnetostática.

Calculamos el campo magnético mediante $\Bvec = \curl{\Avec}$ en coordenadas cilíndricas. Como $\Avec = A_z(s,t)\zhat$:
\begin{align}
    \Bvec(s,t) &= -\pdv{A_z}{s}\,\phihat = -\frac{\mu_0 I_0}{2\pi} \pdv{}{s}\Big[ \ln\left(ct + \sqrt{(ct)^2 - s^2}\right) - \ln(s) \Big]\,\phihat \nonumber \\
    &= -\frac{\mu_0 I_0}{2\pi} \left[ \frac{1}{ct + \sqrt{(ct)^2 - s^2}}\left(\frac{-s}{\sqrt{(ct)^2 - s^2}}\right) - \frac{1}{s} \right] \phihat \nonumber \\
    &= \frac{\mu_0 I_0}{2\pi s} \left[ 1 + \frac{s^2}{\sqrt{(ct)^2 - s^2}\,\big(ct + \sqrt{(ct)^2 - s^2}\big)} \right] \phihat
\end{align}
Racionalizamos el segundo término multiplicando numerador y denominador por $\big(ct - \sqrt{(ct)^2 - s^2}\big)$, notando que en el denominador la diferencia de cuadrados da $(ct)^2 - \big((ct)^2 - s^2\big) = s^2$:
\begin{align}
    \Bvec(s,t) &= \frac{\mu_0 I_0}{2\pi s} \left[ 1 + \frac{s^2\big(ct - \sqrt{(ct)^2 - s^2}\big)}{s^2 \sqrt{(ct)^2 - s^2}} \right] \phihat \nonumber \\
    &= \frac{\mu_0 I_0}{2\pi s} \left[ 1 + \frac{ct}{\sqrt{(ct)^2 - s^2}} - 1 \right] \phihat = \frac{\mu_0 I_0}{2\pi s} \frac{ct}{\sqrt{(ct)^2 - s^2}}\,\phihat
\end{align}
En el límite estacionario $t \to \infty$, el factor $\frac{ct}{\sqrt{(ct)^2 - s^2}} \to 1$, recuperando la ley de Biot-Savart para un hilo infinito: $\Bvec \to \frac{\mu_0 I_0}{2\pi s}\phihat$.
\end{solucion}

% -------------------------------------------------------------------------
\section[Potenciales de Liénard-Wiechert]{Potenciales de Liénard-Wiechert y Campos de una Carga Puntual en Movimiento}

\begin{ejercicio}{Potenciales de Liénard-Wiechert y campo electromagnético de una carga puntual con velocidad constante}{lienard_wiechert_rectilineo}
Una carga puntual $q$ se mueve en línea recta con velocidad constante $\bm{v} = v\xhat$ a lo largo del eje $x$, pasando por el origen en $t = 0$, de forma que su trayectoria es $\bm{w}(t) = vt\xhat$.
\begin{enumerate}[label=\alph*)]
    \item Determine explícitamente el tiempo retardado $t_r$ en función de las coordenadas espaciales $(x,y,z)$ del punto de observación y del tiempo actual $t$, y evalúe los potenciales de Liénard-Wiechert $V(\bm{r},t)$ y $\Avec(\bm{r},t)$.
    \item Demuestre que el campo eléctrico resultante apunta exactamente hacia la posición \textbf{actual} (instantánea) de la carga en el instante $t$ (no hacia la posición retardada) y tiene la expresión:
    \begin{equation*}
        \Evec(\bm{r},t) = \frac{q}{4\pi\varepsilon_0} \frac{1 - v^2/c^2}{\big(1 - \frac{v^2}{c^2}\sin^2\theta\big)^{3/2}} \frac{\bm{R}}{R^3}
    \end{equation*}
    donde $\bm{R} = \bm{r} - \bm{v}t$ es el vector desde la posición actual de la carga hasta el punto de campo y $\theta$ es el ángulo entre $\bm{R}$ y $\bm{v}$.
\end{enumerate}
\end{ejercicio}

\begin{solucion}
\textbf{Parte a) Cálculo explícito del tiempo retardado $t_r$ y potenciales de Liénard-Wiechert.}

El tiempo retardado $t_r$ se define implícitamente por la condición del cono de luz pasado:
\begin{equation}
    \mathscr{r} = |\bm{r} - \bm{w}(t_r)| = c(t - t_r)
\end{equation}
Como $\bm{r} = (x,y,z)$ y $\bm{w}(t_r) = (vt_r, 0, 0)$, elevando al cuadrado ambos miembros obtenemos una ecuación cuadrática para $t_r$:
\begin{align}
    (x - vt_r)^2 + y^2 + z^2 &= c^2(t - t_r)^2 \\
    x^2 - 2xvt_r + v^2 t_r^2 + y^2 + z^2 &= c^2 t^2 - 2c^2 t t_r + c^2 t_r^2 \\
    (c^2 - v^2)t_r^2 - 2(c^2 t - xv)t_r + (c^2 t^2 - r^2) &= 0 \label{eq:cuad_tr}
\end{align}
donde $r^2 = x^2 + y^2 + z^2$. Resolvemos para $t_r$ usando la fórmula cuadrática:
\begin{equation}
    t_r = \frac{(c^2 t - xv) \pm \sqrt{(c^2 t - xv)^2 - (c^2 - v^2)(c^2 t^2 - r^2)}}{c^2 - v^2}
\end{equation}
Simplifiquemos algebraicamente el discriminante $\Delta$ bajo la raíz cuadrada:
\begin{align}
    \Delta &= c^4 t^2 - 2c^2 xvt + x^2 v^2 - \big(c^4 t^2 - c^2 r^2 - v^2 c^2 t^2 + v^2 r^2\big) \nonumber \\
    &= c^2\big(r^2 - 2xvt + v^2 t^2\big) - v^2(r^2 - x^2) \nonumber \\
    &= c^2\big[(x - vt)^2 + y^2 + z^2\big] - v^2(y^2 + z^2) \nonumber \\
    &= c^2\left[ (x - vt)^2 + \left(1 - \frac{v^2}{c^2}\right)(y^2 + z^2) \right]
\end{align}
Para elegir el signo correcto ($\pm$) delante de la raíz, evaluamos el límite no relativista o carga en reposo ($v = 0$), donde sabemos que $t_r = t - r/c$: con $v=0$, $t_r = \frac{c^2 t \pm c r}{c^2} = t \pm \frac{r}{c}$. Por causalidad ($t_r < t$), debemos tomar el signo negativo:
\begin{equation}
    t_r = \frac{c^2 t - xv - c\sqrt{(x - vt)^2 + (1 - v^2/c^2)(y^2 + z^2)}}{c^2 - v^2}
\end{equation}

Los potenciales de Liénard-Wiechert para una carga puntual son:
\begin{equation}
    V(\bm{r},t) = \frac{1}{4\pi\varepsilon_0} \frac{q c}{\big(\mathscr{r}c - \rsep\cdot\bm{v}\big)}, \qquad \Avec(\bm{r},t) = \frac{\bm{v}}{c^2} V(\bm{r},t)
\end{equation}
Calculemos el denominador relativista $\mathscr{r}c - \rsep\cdot\bm{v}$. Como $\mathscr{r} = c(t - t_r)$ y $\rsep\cdot\bm{v} = (x - vt_r)v$:
\begin{equation}
    \mathscr{r}c - \rsep\cdot\bm{v} = c^2(t - t_r) - v(x - vt_r) = (c^2 t - xv) - (c^2 - v^2)t_r
\end{equation}
Pero de la solución de $t_r$ obtenida arriba, vemos inmediatamente que $(c^2 - v^2)t_r = (c^2 t - xv) - \sqrt{\Delta}$. ¡Los términos $(c^2 t - xv)$ se cancelan de forma exacta! Así:
\begin{equation}
    \mathscr{r}c - \rsep\cdot\bm{v} = \sqrt{\Delta} = c\sqrt{(x - vt)^2 + \left(1 - \frac{v^2}{c^2}\right)(y^2 + z^2)}
\end{equation}
Sustituyendo en $V(\bm{r},t)$ y $\Avec(\bm{r},t)$:
\begin{equation}
    V(\bm{r},t) = \frac{1}{4\pi\varepsilon_0} \frac{q}{\sqrt{(x - vt)^2 + (1 - v^2/c^2)(y^2 + z^2)}}, \qquad \Avec(\bm{r},t) = \frac{\bm{v}}{c^2} V(\bm{r},t)
\end{equation}

\medskip
\textbf{Parte b) Cálculo del campo eléctrico $\Evec(\bm{r},t)$.}

Podemos obtener $\Evec = -\grad V - \pdv{\Avec}{t}$ derivando directamente los potenciales obtenidos en (a). Sea $u \equiv (x - vt)^2 + (1 - \beta^2)(y^2 + z^2)$ con $\beta = v/c$, de modo que $V = \frac{q}{4\pi\varepsilon_0} u^{-1/2}$ y $\Avec = \frac{v}{c^2}V\xhat$:
\begin{align}
    \pdv{V}{x} &= -\frac{q}{4\pi\varepsilon_0} \frac{x - vt}{u^{3/2}}, \qquad \pdv{V}{y} = -\frac{q}{4\pi\varepsilon_0} \frac{(1 - \beta^2)y}{u^{3/2}}, \qquad \pdv{V}{z} = -\frac{q}{4\pi\varepsilon_0} \frac{(1 - \beta^2)z}{u^{3/2}} \\
    \pdv{A_x}{t} &= \frac{v}{c^2}\pdv{V}{t} = \frac{v}{c^2}\left( +\frac{q}{4\pi\varepsilon_0}\frac{v(x - vt)}{u^{3/2}} \right) = \frac{q}{4\pi\varepsilon_0} \frac{\beta^2(x - vt)}{u^{3/2}}
\end{align}
Por tanto, la componente $x$ del campo eléctrico es:
\begin{equation}
    E_x = -\pdv{V}{x} - \pdv{A_x}{t} = \frac{q}{4\pi\varepsilon_0} \frac{(1 - \beta^2)(x - vt)}{u^{3/2}}
\end{equation}
Agrupando las tres componentes $(E_x, E_y, E_z)$ y definiendo el vector desde la posición actual $\bm{R} \equiv \bm{r} - \bm{v}t = (x - vt)\xhat + y\yhat + z\zhat$:
\begin{equation}
    \Evec(\bm{r},t) = \frac{q}{4\pi\varepsilon_0} \frac{1 - \beta^2}{\big[(x - vt)^2 + (1 - \beta^2)(y^2 + z^2)\big]^{3/2}} \bm{R}
\end{equation}
Si $\theta$ es el ángulo entre $\bm{R}$ y el eje $x$ (dirección de $\bm{v}$), entonces $x - vt = R\cos\theta$ y $y^2 + z^2 = R^2\sin^2\theta$. El corchete en el denominador se transforma en:
\begin{equation}
    (x - vt)^2 + (1 - \beta^2)(y^2 + z^2) = R^2\big(\cos^2\theta + \sin^2\theta - \beta^2\sin^2\theta\big) = R^2\big(1 - \beta^2\sin^2\theta\big)
\end{equation}
Sustituyendo en $\Evec(\bm{r},t)$ con $\beta = v/c$:
\begin{equation}
    \Evec(\bm{r},t) = \frac{q}{4\pi\varepsilon_0} \frac{1 - v^2/c^2}{\left(1 - \frac{v^2}{c^2}\sin^2\theta\right)^{3/2}} \frac{\bm{R}}{R^3}
\end{equation}
\end{solucion}

% =========================================================================
% CAPÍTULO 11: RADIACIÓN
% =========================================================================
\chapter{Radiación Electromagnética}

\section{Radiación de Dipolo Eléctrico y Dipolo Magnético}

\begin{ejercicio}{Esfera con polarización o dipolo rotante: Campos de radiación y potencia total radiada}{dipolo_rotante}
Un dipolo eléctrico de magnitud constante $p_0$ situado en el origen gira en el plano $xy$ con velocidad angular uniforme $\omega$, de forma que su momento dipolar es $\bm{p}(t) = p_0\big[\cos(\omega t)\xhat + \sin(\omega t)\yhat\big]$.
\begin{enumerate}[label=\alph*)]
    \item Determine los campos eléctrico y magnético en la zona de radiación ($r \gg c/\omega$), el vector de Poynting promediado en el tiempo $\langle\Svec\rangle$ y la distribución angular de potencia radiada $\dv{P}{\Omega}$.
    \item Integre sobre toda la esfera para encontrar la potencia total radiada $P$.
\end{enumerate}
\end{ejercicio}

\begin{solucion}
\textbf{Parte a) Campos en la zona de radiación y distribución angular de potencia.}

Para un dipolo eléctrico puntual arbitrario $\bm{p}(t)$ en el origen, las fórmulas generales para los campos eléctrico y magnético en la zona de radiación ($r \to \infty$, conservando términos de orden $1/r$) son:
\begin{equation}
    \Evec_{\text{rad}}(\bm{r},t) = \frac{\mu_0}{4\pi r} \Big[ \big(\ddot{\bm{p}}(t_r) \cdot \thetahat\big)\thetahat + \big(\ddot{\bm{p}}(t_r) \cdot \phihat\big)\phihat \Big], \qquad \Bvec_{\text{rad}}(\bm{r},t) = -\frac{1}{c}\big(\rhat \times \Evec_{\text{rad}}\big)
\end{equation}
donde $t_r = t - r/c$ y los puntos denotan derivadas respecto al tiempo retardado. Derivamos dos veces $\bm{p}(t_r)$:
\begin{equation}
    \ddot{\bm{p}}(t_r) = -\omega^2 p_0 \big[ \cos(\omega t_r)\xhat + \sin(\omega t_r)\yhat \big]
\end{equation}
Expresamos los productos escalares de $\xhat$ e $\yhat$ con los vectores unitarios esféricos $\thetahat$ y $\phihat$:
\begin{align}
    \thetahat &= \cos\theta\cos\phi\,\xhat + \cos\theta\sin\phi\,\yhat - \sin\theta\,\zhat \\
    \phihat &= -\sin\phi\,\xhat + \cos\phi\,\yhat
\end{align}
Usando las identidades del coseno y seno de la diferencia de ángulos:
\begin{align}
    \ddot{\bm{p}}(t_r) \cdot \thetahat &= -\omega^2 p_0 \cos\theta \big(\cos(\omega t_r)\cos\phi + \sin(\omega t_r)\sin\phi\big) = -\omega^2 p_0 \cos\theta \cos(\omega t_r - \phi) \\
    \ddot{\bm{p}}(t_r) \cdot \phihat &= -\omega^2 p_0 \big(-\cos(\omega t_r)\sin\phi + \sin(\omega t_r)\cos\phi\big) = -\omega^2 p_0 \sin(\omega t_r - \phi)
\end{align}
Por lo tanto, el campo eléctrico radiado es:
\begin{equation}
    \Evec_{\text{rad}}(r,\theta,\phi,t) = -\frac{\mu_0 p_0 \omega^2}{4\pi r} \Big[ \cos\theta \cos\big(\omega(t - r/c) - \phi\big)\,\thetahat + \sin\big(\omega(t - r/c) - \phi\big)\,\phihat \Big]
\end{equation}
Notemos que a lo largo del eje $z$ ($\theta = 0$), $\cos\theta = 1$ y la onda radiada está circularmente polarizada, mientras que en el plano ecuatorial ($\theta = \pi/2$), $\cos\theta = 0$ y la onda está linealmente polarizada a lo largo de $\phihat$.

El vector de Poynting en la zona de radiación es $\Svec = \frac{1}{\mu_0 c} E_{\text{rad}}^2 \,\rhat = \varepsilon_0 c \big(E_\theta^2 + E_\phi^2\big)\rhat$:
\begin{equation}
    \Svec(\bm{r},t) = \frac{\mu_0 p_0^2 \omega^4}{16\pi^2 c r^2} \Big[ \cos^2\theta \cos^2(\omega t_r - \phi) + \sin^2(\omega t_r - \phi) \Big] \rhat
\end{equation}
Promediando sobre un período completo $T = 2\pi/\omega$, dado que $\langle \cos^2(\omega t_r - \phi) \rangle = \langle \sin^2(\omega t_r - \phi) \rangle = \frac{1}{2}$:
\begin{equation}
    \langle \Svec \rangle = \frac{\mu_0 p_0^2 \omega^4}{32\pi^2 c r^2} \big(1 + \cos^2\theta\big)\,\rhat
\end{equation}
La potencia media radiada por unidad de ángulo sólido es:
\begin{equation}
    \dv{P}{\Omega} = r^2 \big(\langle \Svec \rangle \cdot \rhat\big) = \frac{\mu_0 p_0^2 \omega^4}{32\pi^2 c} \big(1 + \cos^2\theta\big)
\end{equation}

\medskip
\textbf{Parte b) Potencia total radiada $P$.}

Integramos $\dv{P}{\Omega}$ sobre toda la esfera con $d\Omega = \sin\theta\,d\theta\,d\phi$:
\begin{align}
    P &= \int \dv{P}{\Omega}\,d\Omega = \frac{\mu_0 p_0^2 \omega^4}{32\pi^2 c} \int_{0}^{2\pi} d\phi \int_{0}^{\pi} \big(1 + \cos^2\theta\big)\sin\theta\,d\theta \nonumber \\
    &= \frac{\mu_0 p_0^2 \omega^4}{32\pi^2 c} (2\pi) \left[ -\cos\theta - \frac{\cos^3\theta}{3} \right]_{0}^{\pi} \nonumber \\
    &= \frac{\mu_0 p_0^2 \omega^4}{16\pi c} \left( 2 + \frac{2}{3} \right) = \frac{\mu_0 p_0^2 \omega^4}{16\pi c}\left(\frac{8}{3}\right) = \frac{\mu_0 p_0^2 \omega^4}{6\pi c}
\end{align}
(Nota: Esto es exactamente el doble de la potencia $\frac{\mu_0 p_0^2 \omega^4}{12\pi c}$ de un dipolo lineal oscilante, ya que el dipolo rotante equivale a la superposición de dos dipolos lineales ortogonales en cuadratura de fase).
\end{solucion}

% -------------------------------------------------------------------------
\section[Radiación de cargas aceleradas]{Radiación de Cargas Puntuales Aceleradas (Fórmula de Larmor y Generalización de Liénard)}

\begin{ejercicio}{Colapso radiativo del átomo clásico de Bohr y radiación de sincrotrón relativista}{colapso_atomo_larmor}
En el modelo planetario clásico del átomo de hidrógeno (Rutherford), un electrón de carga $-e$ y masa $m$ orbita alrededor de un protón estacionario de carga $+e$.
\begin{enumerate}[label=\alph*)]
    \item Asumiendo que en cada vuelta la órbita es aproximadamente circular de radio $r(t)$ y que la pérdida de energía mecánica total $U = -\frac{e^2}{8\pi\varepsilon_0 r}$ se debe a la radiación electromagnética dada por la fórmula de Larmor $P = \frac{\mu_0 e^2 a^2}{6\pi c}$, plantee la ecuación diferencial para $r(t)$ y calcule el tiempo de vida clásico $t_{\text{colapso}}$ que tarda el electrón en caer en espiral desde el radio de Bohr $r_0 = 0.529\,\text{\AA}$ hasta el núcleo ($r = 0$).
    \item Si una carga puntual se mueve en una órbita circular con velocidad relativista ($\bm{v} \perp \bm{a}$), utilice la fórmula general de Liénard para obtener la potencia radiada en función del factor de Lorentz $\gamma = (1 - v^2/c^2)^{-1/2}$.
\end{enumerate}
\end{ejercicio}

\begin{solucion}
\textbf{Parte a) Ecuación diferencial del decaimiento orbital y tiempo de colapso.}

Para un electrón en una órbita circular de radio $r$ alrededor del protón, la fuerza electrostática de Coulomb proporciona la aceleración centrípeta $a = \frac{v^2}{r}$:
\begin{equation}
    m a = \frac{m v^2}{r} = \frac{e^2}{4\pi\varepsilon_0 r^2} \implies a = \frac{e^2}{4\pi\varepsilon_0 m r^2}
\end{equation}
La energía mecánica total (cinética más potencial) del electrón en una órbita circular de radio $r$ es:
\begin{equation}
    U = K + V = \frac{1}{2}mv^2 - \frac{e^2}{4\pi\varepsilon_0 r} = \frac{e^2}{8\pi\varepsilon_0 r} - \frac{e^2}{4\pi\varepsilon_0 r} = -\frac{e^2}{8\pi\varepsilon_0 r}
\end{equation}
Derivando $U(r)$ respecto al tiempo $t$:
\begin{equation}
    \dv{U}{t} = \frac{e^2}{8\pi\varepsilon_0 r^2} \dv{r}{t}
\end{equation}
Por conservación de la energía, la tasa de pérdida de energía mecánica $-\dv{U}{t}$ debe ser igual a la potencia radiada según la fórmula de Larmor (usando $\mu_0 = \frac{1}{\varepsilon_0 c^2}$):
\begin{equation}
    -\dv{U}{t} = P = \frac{\mu_0 e^2 a^2}{6\pi c} = \frac{e^2}{6\pi\varepsilon_0 c^3} \left(\frac{e^2}{4\pi\varepsilon_0 m r^2}\right)^2
\end{equation}
Igualando ambas expresiones:
\begin{equation}
    -\frac{e^2}{8\pi\varepsilon_0 r^2} \dv{r}{t} = \frac{e^2}{6\pi\varepsilon_0 c^3} \left(\frac{e^2}{4\pi\varepsilon_0 m}\right)^2 \frac{1}{r^4}
\end{equation}
Multiplicando por $r^2$ y despejando $r^2 \dv{r}{t}$:
\begin{equation}
    r^2 \dv{r}{t} = -\frac{4}{3 c^3}\left(\frac{e^2}{4\pi\varepsilon_0 m}\right)^2 = -\frac{4}{3} r_e^2 c
\end{equation}
donde hemos identificado el radio clásico del electrón $r_e \equiv \frac{e^2}{4\pi\varepsilon_0 m c^2} \approx 2.818 \times 10^{-15}\,\text{m}$.

Separamos variables e integramos desde $t = 0$ (con $r = r_0$) hasta $t = t_{\text{colapso}}$ (con $r = 0$):
\begin{equation}
    \int_{r_0}^{0} r^2\,dr = -\frac{4}{3} r_e^2 c \int_{0}^{t_{\text{colapso}}} dt \implies -\frac{r_0^3}{3} = -\frac{4}{3} r_e^2 c \, t_{\text{colapso}}
\end{equation}
Despejando el tiempo de colapso:
\begin{equation}
    t_{\text{colapso}} = \frac{r_0^3}{4 r_e^2 c} = \frac{\big(4\pi\varepsilon_0 m\big)^2 c^3 r_0^3}{4 e^4}
\end{equation}
Sustituyendo los valores numéricos ($r_0 = 5.29 \times 10^{-11}\,\text{m}$, $r_e = 2.818 \times 10^{-15}\,\text{m}$, $c = 3.00 \times 10^8\,\text{m/s}$):
\begin{equation}
    t_{\text{colapso}} = \frac{(5.29 \times 10^{-11}\,\text{m})^3}{4 (2.818 \times 10^{-15}\,\text{m})^2 (3.00 \times 10^8\,\text{m/s})} \approx 1.56 \times 10^{-11}\,\text{s}
\end{equation}
El átomo clásico colapsaría en apenas $15.6$ picosegundos.

\medskip
\textbf{Parte b) Radiación de sincrotrón mediante la fórmula de Liénard.}

La generalización relativista de la fórmula de Larmor obtenida por Liénard para una carga puntual con velocidad arbitraria $\bm{v}$ y aceleración $\bm{a}$ es:
\begin{equation}
    P = \frac{\mu_0 q^2 \gamma^6}{6\pi c} \left( a^2 - \left|\frac{\bm{v} \times \bm{a}}{c}\right|^2 \right)
\end{equation}
En un movimiento circular (sincrotrón), la velocidad $\bm{v}$ es perpendicular a la aceleración centrípeta $\bm{a}$ ($\bm{v} \perp \bm{a}$), por lo que $|\bm{v} \times \bm{a}| = v a$. Sustituyendo esto en el paréntesis:
\begin{equation}
    a^2 - \left|\frac{\bm{v} \times \bm{a}}{c}\right|^2 = a^2\left(1 - \frac{v^2}{c^2}\right) = \frac{a^2}{\gamma^2}
\end{equation}
Por lo tanto, la potencia de radiación de sincrotrón es:
\begin{equation}
    P_{\text{sincrotrón}} = \frac{\mu_0 q^2 \gamma^4 a^2}{6\pi c}
\end{equation}
creciendo con la cuarta potencia del factor de Lorentz $\gamma^4$.
\end{solucion}

% =========================================================================
% CAPÍTULO 12: ELECTRODINÁMICA Y RELATIVIDAD
% =========================================================================
\chapter{Electrodinámica y Relatividad Especial}

\section[Cinemática y dinámica relativista]{Cinemática y Dinámica Relativista (Cuadrivectores y Efecto Compton / Umbral)}

\begin{ejercicio}{Movimiento hiperbólico relativista bajo una fuerza constante (Campo Eléctrico Uniforme)}{movimiento_hiperbolico}
Una partícula cargada de masa en reposo $m$ y carga $q$ se encuentra inicialmente en reposo en el origen ($x(0) = 0, \bm{u}(0) = \bm{0}$) y en $t = 0$ se enciende un campo eléctrico uniforme $\Evec = E_0 \xhat$, el cual ejerce una fuerza relativista constante $\bm{F} = \dv{\bm{p}}{t} = qE_0\xhat \equiv F_0\xhat$.
\begin{enumerate}[label=\alph*)]
    \item Integre la ecuación relativista del movimiento para hallar la velocidad $u(t)$ y la posición $x(t)$ en el sistema de referencia del laboratorio, y demuestre que la trayectoria en el diagrama espacio-tiempo $(x, ct)$ es una hipérbola.
    \item Calcule el tiempo propio $\tau$ transcurrido en el reloj de la partícula en función del tiempo coordenado $t$, y exprese $x(\tau)$ y $t(\tau)$ en términos de funciones hiperbólicas de $\tau$.
\end{enumerate}
\end{ejercicio}

\begin{solucion}
\textbf{Parte a) Velocidad $u(t)$ y trayectoria hiperbólica $x(t)$.}

El momento lineal relativista de una partícula que se mueve a lo largo del eje $x$ con velocidad $u = \dv{x}{t}$ es:
\begin{equation}
    p = \gamma m u = \frac{m u}{\sqrt{1 - u^2/c^2}}
\end{equation}
La ecuación de movimiento relativista $\dv{p}{t} = F_0$ con la condición inicial $p(0) = 0$ se integra inmediatamente respecto a $t$:
\begin{equation}
    p(t) = F_0 t \implies \frac{m u}{\sqrt{1 - u^2/c^2}} = F_0 t
\end{equation}
Definamos la aceleración constante característica $a_0 \equiv \frac{F_0}{m} = \frac{q E_0}{m}$. Elevando ambos miembros al cuadrado para despejar $u(t)$:
\begin{equation}
    \frac{u^2}{1 - u^2/c^2} = (a_0 t)^2 \implies u^2 \left[ 1 + \left(\frac{a_0 t}{c}\right)^2 \right] = (a_0 t)^2 \implies u(t) = \dv{x}{t} = \frac{a_0 t}{\sqrt{1 + (a_0 t / c)^2}}
\end{equation}
Nótese que cuando $t \to \infty$, $u(t) \to c$, sin superar jamás la velocidad de la luz.

Integramos ahora $u(t)$ desde $t' = 0$ hasta $t' = t$ con $x(0) = 0$:
\begin{equation}
    x(t) = \int_{0}^{t} \frac{a_0 t'}{\sqrt{1 + (a_0 t' / c)^2}}\,dt' = \frac{c^2}{a_0} \left[ \sqrt{1 + \left(\frac{a_0 t'}{c}\right)^2} \right]_{0}^{t} = \frac{c^2}{a_0} \left[ \sqrt{1 + \left(\frac{a_0 t}{c}\right)^2} - 1 \right]
\end{equation}
Reordenando esta ecuación:
\begin{equation}
    \left( x + \frac{c^2}{a_0} \right)^2 - (ct)^2 = \left(\frac{c^2}{a_0}\right)^2
\end{equation}
lo que constituye la ecuación canónica de una hipérbola en el plano espacio-temporal $(x, ct)$.

\medskip
\textbf{Parte b) Parametrización en términos del tiempo propio $\tau$.}

El diferencial de tiempo propio $d\tau$ se relaciona con $dt$ mediante el factor de Lorentz instantáneo $\gamma(t)$:
\begin{equation}
    \gamma(t) = \frac{1}{\sqrt{1 - u^2(t)/c^2}} = \sqrt{1 + \left(\frac{a_0 t}{c}\right)^2}
\end{equation}
Por lo tanto:
\begin{equation}
    d\tau = \frac{dt}{\gamma(t)} = \frac{dt}{\sqrt{1 + (a_0 t / c)^2}}
\end{equation}
Integrando desde $t=0$ ($\tau = 0$) mediante el cambio $a_0 t / c = \sinh\xi$:
\begin{equation}
    \tau = \int_{0}^{t} \frac{dt'}{\sqrt{1 + (a_0 t' / c)^2}} = \frac{c}{a_0} \operatorname{arcsinh}\left(\frac{a_0 t}{c}\right)
\end{equation}
Invirtiendo la función seno hiperbólico para despejar $t(\tau)$ y sustituyendo en $x(t)$ (usando $1 + \sinh^2\xi = \cosh^2\xi$):
\begin{equation}
    t(\tau) = \frac{c}{a_0} \sinh\left(\frac{a_0 \tau}{c}\right), \qquad x(\tau) = \frac{c^2}{a_0} \left[ \cosh\left(\frac{a_0 \tau}{c}\right) - 1 \right]
\end{equation}
\end{solucion}

% -------------------------------------------------------------------------
\section{Transformación de Lorentz de los Campos \texorpdfstring{$\Evec$}{E} y \texorpdfstring{$\Bvec$}{B}}

\begin{ejercicio}{Origen relativista del magnetismo: Hilo conductor con corriente visto desde otro sistema inercial y Ley de Biot-Savart relativista}{transformacion_campos_hilo}
\begin{enumerate}[label=\alph*)]
    \item En el sistema de referencia del laboratorio $\mathcal{S}$, un alambre recto infinito a lo largo del eje $x$ es eléctricamente neutro ($\lambda = 0$) y lleva una corriente constante $I = \lambda_+ v_0$ hacia $+x$ (formada por iones positivos en reposo con densidad lineal $-\lambda_-$ y electrones moviéndose con velocidad $-v_0\xhat$, o viceversa). Una carga de prueba $q$ se mueve paralelamente al alambre con velocidad $\bm{u} = u\xhat$ a una distancia $y = s$. Utilizando la contracción de Lorentz de las longitudes para hallar la densidad de carga neta $\bar{\lambda}$ en el sistema de reposo $\bar{\mathcal{S}}$ de la carga $q$, calcule la fuerza eléctrica en $\bar{\mathcal{S}}$, transfórmela de vuelta a $\mathcal{S}$ y demuestre que coincide exactamente con la fuerza magnética $q\bm{u}\times\Bvec$.
    \item A partir de las leyes de transformación de Lorentz de los campos electromagnéticos ($\Evec, \Bvec$), deduzca el campo magnético $\Bvec(\bm{r},t)$ de una carga puntual $q$ que se mueve con velocidad constante $\bm{v} = v\xhat$.
\end{enumerate}
\end{ejercicio}

\begin{solucion}
\textbf{Parte a) Contracción de Lorentz de densidades de carga y fuerza magnética.}

Modelamos el alambre en el sistema $\mathcal{S}$ como dos trenes de cargas superpuestos que se mueven en sentidos opuestos con rapidez $v_0$ a lo largo del eje $x$: cargas positivas con densidad lineal $+\lambda_0$ moviéndose con velocidad $\bm{v}_+ = +v_0\xhat$, y cargas negativas con densidad lineal $-\lambda_0$ moviéndose con velocidad $\bm{v}_- = -v_0\xhat$.
En $\mathcal{S}$, la densidad de carga total es nula y la corriente neta es:
\begin{equation}
    \lambda_{\text{tot}} = \lambda_0 - \lambda_0 = 0, \qquad I = \lambda_0 v_0 + (-\lambda_0)(-v_0) = 2\lambda_0 v_0
\end{equation}
Si las densidades propias (en el sistema de reposo de cada tren de cargas) son $\pm\lambda_{0,\text{propia}}$, debido a la contracción de Lorentz en $\mathcal{S}$ se cumple $\lambda_0 = \gamma_0 \lambda_{0,\text{propia}}$, donde $\gamma_0 = (1 - v_0^2/c^2)^{-1/2}$.

Pasemos ahora al sistema inercial $\bar{\mathcal{S}}$ que se mueve con velocidad $u\xhat$ respecto a $\mathcal{S}$ (de modo que la carga de prueba $q$ está en reposo en $\bar{\mathcal{S}}$, $\bar{\bm{u}} = \bm{0}$). Por la fórmula relativista de adición de velocidades de Einstein, las velocidades $\bar{v}_\pm$ de las cargas positivas y negativas del alambre vistas desde $\bar{\mathcal{S}}$ son:
\begin{equation}
    \bar{v}_\pm = \frac{\pm v_0 - u}{1 \mp u v_0 / c^2}
\end{equation}
Calculemos los factores de Lorentz correspondientes $\bar{\gamma}_\pm = \big(1 - \bar{v}_\pm^2/c^2\big)^{-1/2}$:
\begin{align}
    1 - \frac{\bar{v}_\pm^2}{c^2} &= 1 - \frac{(v_0 \mp u)^2}{c^2(1 \mp u v_0 / c^2)^2} = \frac{(c^2 \mp u v_0)^2 - (c v_0 \mp c u)^2}{c^2(c^2 \mp u v_0)^2 / c^4} \nonumber \\
    &= \frac{c^4 \mp 2c^2 u v_0 + u^2 v_0^2 - c^2 v_0^2 \pm 2c^2 u v_0 - c^2 u^2}{(c^2 \mp u v_0)^2} = \frac{(c^2 - u^2)(c^2 - v_0^2)}{(c^2 \mp u v_0)^2} \nonumber \\
    &= \frac{(1 - u^2/c^2)(1 - v_0^2/c^2)}{(1 \mp u v_0 / c^2)^2} = \frac{1}{\gamma_u^2 \gamma_0^2 (1 \mp u v_0 / c^2)^2}
\end{align}
donde $\gamma_u = (1 - u^2/c^2)^{-1/2}$. Tomando la raíz cuadrada recíproca:
\begin{equation}
    \bar{\gamma}_\pm = \gamma_u \gamma_0 \left(1 \mp \frac{u v_0}{c^2}\right)
\end{equation}
En el sistema $\bar{\mathcal{S}}$, las densidades lineales de carga están contraídas por $\bar{\gamma}_\pm$ respecto a la densidad propia $\lambda_{0,\text{propia}} = \lambda_0 / \gamma_0$:
\begin{equation}
    \bar{\lambda}_\pm = \pm \bar{\gamma}_\pm \lambda_{0,\text{propia}} = \pm \bar{\gamma}_\pm \left(\frac{\lambda_0}{\gamma_0}\right) = \pm \gamma_u \lambda_0 \left(1 \mp \frac{u v_0}{c^2}\right)
\end{equation}
¡Como $\bar{v}_- \neq \bar{v}_+$, las cargas negativas están más contraídas que las positivas! La densidad lineal de carga neta $\bar{\lambda}_{\text{tot}}$ en $\bar{\mathcal{S}}$ ya no es cero:
\begin{equation}
    \bar{\lambda}_{\text{tot}} = \bar{\lambda}_+ + \bar{\lambda}_- = \gamma_u \lambda_0 \left[ \left(1 - \frac{u v_0}{c^2}\right) - \left(1 + \frac{u v_0}{c^2}\right) \right] = -\frac{2 \gamma_u \lambda_0 v_0 u}{c^2} = -\frac{\gamma_u I u}{c^2}
\end{equation}
En $\bar{\mathcal{S}}$, la carga $q$ está en reposo, por lo que no experimenta ninguna fuerza magnética; solo siente el campo electrostático $\bar{\Evec} = \frac{\bar{\lambda}_{\text{tot}}}{2\pi\varepsilon_0 s}\yhat$ generado por el alambre cargado:
\begin{equation}
    \bar{\bm{F}} = q\bar{\Evec} = -\frac{q \gamma_u I u}{2\pi\varepsilon_0 c^2 s}\yhat = -\gamma_u \left(\frac{\mu_0 I q u}{2\pi s}\right)\yhat
\end{equation}
(usando $\frac{1}{\varepsilon_0 c^2} = \mu_0$). Finalmente, según la ley de transformación relativista de la fuerza transversal para una partícula en reposo en $\bar{\mathcal{S}}$ ($F_\perp = \bar{F}_\perp / \gamma_u$):
\begin{equation}
    \bm{F} = \frac{\bar{\bm{F}}}{\gamma_u} = -\frac{\mu_0 I}{2\pi s} q u \,\yhat = q(u\xhat) \times \left(\frac{\mu_0 I}{2\pi s}\zhat\right) = q \bm{u} \times \Bvec
\end{equation}
Esto demuestra que la fuerza magnética es una consecuencia necesaria de la electrostática combinada con la relatividad especial.

\medskip
\textbf{Parte b) Campo magnético de una carga en movimiento uniforme mediante transformación de Lorentz.}

Sea $\bar{\mathcal{S}}$ el sistema de referencia inercial en el cual la carga puntual $q$ está en reposo en el origen $\bar{\bm{r}} = \bm{0}$, el cual se mueve con velocidad $\bm{v} = v\xhat$ respecto al laboratorio $\mathcal{S}$. En $\bar{\mathcal{S}}$, solo hay campo electrostático de Coulomb y el campo magnético es nulo:
\begin{equation}
    \bar{\Evec} = \frac{q}{4\pi\varepsilon_0 \bar{r}^3}\big(\bar{x}\xhat + \bar{y}\yhat + \bar{z}\zhat\big), \qquad \bar{\Bvec} = \bm{0}
\end{equation}
Las reglas generales de transformación inversa de Lorentz (de $\bar{\mathcal{S}}$ a $\mathcal{S}$) para los campos son:
\begin{align}
    E_x &= \bar{E}_x, & E_y &= \gamma\big(\bar{E}_y + v\bar{B}_z\big) = \gamma\bar{E}_y, & E_z &= \gamma\big(\bar{E}_z - v\bar{B}_y\big) = \gamma\bar{E}_z \\
    B_x &= \bar{B}_x = 0, & B_y &= \gamma\left(\bar{B}_y - \frac{v}{c^2}\bar{E}_z\right) = -\frac{\gamma v}{c^2}\bar{E}_z, & B_z &= \gamma\left(\bar{B}_z + \frac{v}{c^2}\bar{E}_y\right) = +\frac{\gamma v}{c^2}\bar{E}_y
\end{align}
Observando las expresiones de $(B_x, B_y, B_z)$ en función de $(E_x, E_y, E_z)$:
\begin{equation}
    B_x = 0, \qquad B_y = -\frac{v}{c^2} E_z, \qquad B_z = +\frac{v}{c^2} E_y
\end{equation}
Esto equivale vectorialmente a:
\begin{equation}
    \Bvec(\bm{r},t) = \frac{1}{c^2}(\bm{v} \times \Evec(\bm{r},t))
\end{equation}
Sustituyendo las coordenadas transformadas $\bar{x} = \gamma(x - vt)$, $\bar{y} = y$, $\bar{z} = z$ y $\bar{r} = \sqrt{\gamma^2(x-vt)^2 + y^2 + z^2} = \gamma R \sqrt{1 - \frac{v^2}{c^2}\sin^2\theta}$, recuperamos tanto el campo eléctrico del Capítulo 10 como la ley de Biot-Savart relativista:
\begin{equation}
    \Bvec(\bm{r},t) = \frac{\mu_0 q}{4\pi} \frac{1 - v^2/c^2}{\left(1 - \frac{v^2}{c^2}\sin^2\theta\right)^{3/2}} \frac{\bm{v} \times \bm{R}}{R^3}
\end{equation}
\end{solucion}

% -------------------------------------------------------------------------
\section{Formulación Covariante: El Tensor de Campo Electromagnético $F^{\mu\nu}$}

Para dotar a la electrodinámica de una estructura manifiestamente invariante Lorentz y preparar el terreno hacia la electrodinámica cuántica (QED), definimos el cuadripotencial electromagnético contravariante $A^\mu = (\Phi/c, \mathbf{A})$ y el operador cuadrigradiente $\partial_\mu = \left(\frac{1}{c}\frac{\partial}{\partial t}, \nabla\right)$, adoptando la métrica de Minkowski con signatura $\eta_{\mu\nu} = \text{diag}(+1, -1, -1, -1)$.

\begin{definicion}[Tensor de Campo Electromagnético]
El tensor antisimétrico de rango 2 de Faraday-Maxwell se define como la derivada exterior del cuadripotencial:
\begin{equation}
F_{\mu\nu} = \partial_\mu A_\nu - \partial_\nu A_\mu
\end{equation}
Sus componentes explícitas en forma matricial contravariante $F^{\mu\nu} = \eta^{\mu\alpha}\eta^{\nu\beta}F_{\alpha\beta}$ vienen dadas por:
\begin{equation}
F^{\mu\nu} = \begin{pmatrix}
0 & -E_x/c & -E_y/c & -E_z/c \\
E_x/c & 0 & -B_z & B_y \\
E_y/c & B_z & 0 & -B_x \\
E_z/c & -B_y & B_x & 0
\end{pmatrix}
\end{equation}
\end{definicion}

La densidad lagrangiana del campo electromagnético libre en el vacío es el escalar de Lorentz:
\begin{equation}
\mathcal{L}_{\text{EM}} = -\frac{1}{4\mu_0} F_{\mu\nu}F^{\mu\nu} = \frac{1}{2}\left(\varepsilon_0 \mathbf{E}^2 - \frac{1}{\mu_0}\mathbf{B}^2\right)
\end{equation}

Las ecuaciones de Euler-Lagrange aplicadas a $\mathcal{L}_{\text{EM}}$ respecto a $A_\nu$:
\begin{equation}
\partial_\mu \left( \frac{\partial \mathcal{L}_{\text{EM}}}{\partial(\partial_\mu A_\nu)} \right) - \frac{\partial \mathcal{L}_{\text{EM}}}{\partial A_\nu} = 0 \implies \partial_\mu F^{\mu\nu} = 0
\end{equation}
reproducen exactamente las ecuaciones de Maxwell inhomogéneas en ausencia de fuentes. Sin embargo, para efectuar la cuantización canónica evitando grados de libertad redundantes (como la componente temporal $\Phi$, cuyo momento conjugado $\pi^0 = \frac{\partial \mathcal{L}}{\partial \dot{A}_0} = 0$ se anula idénticamente), fijamos el **gauge de radiación o de Coulomb**:
\begin{equation}
\nabla \cdot \mathbf{A} = 0, \qquad \Phi = 0
\end{equation}
Bajo esta condición de transversalidad, el campo eléctrico es puramente transversal $\mathbf{E}_\perp = -\frac{\partial \mathbf{A}}{\partial t}$, y el potencial vector satisface la ecuación de onda d'Alembertiana:
\begin{equation}
\Box \mathbf{A}(\mathbf{r}, t) = \left( \frac{1}{c^2}\frac{\partial^2}{\partial t^2} - \nabla^2 \right) \mathbf{A}(\mathbf{r}, t) = 0
\end{equation}

% =========================================================================
\section{Descomposición Modal en una Cavidad y Osciladores Armónicos Desacoplados}

Consideremos el campo electromagnético confinado en un volumen cúbico de cuantización $V = L^3$ con condiciones de contorno periódicas de Born-von Karman $\mathbf{A}(\mathbf{r} + L\hat{\mathbf{e}}_i, t) = \mathbf{A}(\mathbf{r}, t)$. Los vectores de onda permitidos quedan discretizados según:
\begin{equation}
\mathbf{k} = \frac{2\pi}{L}(n_x, n_y, n_z), \qquad n_x, n_y, n_z \in \mathbb{Z}
\end{equation}

Para cada vector de propagación $\mathbf{k}$, la condición de transversalidad $\nabla \cdot \mathbf{A} = 0$ se traduce en el espacio recíproco como $\mathbf{k} \cdot \bm{\varepsilon}_{\mathbf{k},s} = 0$, donde $\bm{\varepsilon}_{\mathbf{k},s}$ ($s \in \{1, 2\}$) son dos vectores unitarios de polarización ortonormales entre sí y ortogonales a $\mathbf{k}$:
\begin{equation}
\bm{\varepsilon}_{\mathbf{k},s}^* \cdot \bm{\varepsilon}_{\mathbf{k},s'} = \delta_{s,s'}, \qquad \sum_{s=1}^2 (\bm{\varepsilon}_{\mathbf{k},s})_i (\bm{\varepsilon}_{\mathbf{k},s}^*)_j = \delta_{ij} - \frac{k_i k_j}{|\mathbf{k}|^2}
\end{equation}

Expandiendo el potencial vector en serie de Fourier de modos normales con frecuencia angular $\omega_k = c|\mathbf{k}|$:
\begin{equation}
\mathbf{A}(\mathbf{r}, t) = \sum_{\mathbf{k}, s} \sqrt{\frac{\hbar}{2\varepsilon_0 V \omega_k}} \left[ \alpha_{\mathbf{k},s}(t) \bm{\varepsilon}_{\mathbf{k},s} e^{i\mathbf{k}\cdot\mathbf{r}} + \alpha_{\mathbf{k},s}^*(t) \bm{\varepsilon}_{\mathbf{k},s}^* e^{-i\mathbf{k}\cdot\mathbf{r}} \right]
\end{equation}
donde $\alpha_{\mathbf{k},s}(t) = \alpha_{\mathbf{k},s}(0) e^{-i\omega_k t}$ son las amplitudes complejas adimensionales de los modos normales. Los campos eléctrico y magnético clásicos resultan:
\begin{align}
\mathbf{E}(\mathbf{r}, t) &= i \sum_{\mathbf{k}, s} \sqrt{\frac{\hbar \omega_k}{2\varepsilon_0 V}} \left[ \alpha_{\mathbf{k},s}(t) \bm{\varepsilon}_{\mathbf{k},s} e^{i\mathbf{k}\cdot\mathbf{r}} - \alpha_{\mathbf{k},s}^*(t) \bm{\varepsilon}_{\mathbf{k},s}^* e^{-i\mathbf{k}\cdot\mathbf{r}} \right] \\
\mathbf{B}(\mathbf{r}, t) &= \frac{i}{c} \sum_{\mathbf{k}, s} \sqrt{\frac{\hbar \omega_k}{2\varepsilon_0 V}} \left[ \alpha_{\mathbf{k},s}(t) (\hat{\mathbf{k}} \times \bm{\varepsilon}_{\mathbf{k},s}) e^{i\mathbf{k}\cdot\mathbf{r}} - \alpha_{\mathbf{k},s}^*(t) (\hat{\mathbf{k}} \times \bm{\varepsilon}_{\mathbf{k},s}^*) e^{-i\mathbf{k}\cdot\mathbf{r}} \right]
\end{align}

Sustituyendo estas expansiones en el Hamiltoniano clásico del campo electromagnético y utilizando la relación de ortogonalidad integral $\frac{1}{V}\int_V e^{i(\mathbf{k}-\mathbf{k}')\cdot\mathbf{r}} d^3r = \delta_{\mathbf{k},\mathbf{k}'}$, obtenemos el **Teorema de Equivalencia Mecánica de Dirac**:
\begin{equation}
H_{\text{EM}} = \frac{1}{2} \int_V \left( \varepsilon_0 |\mathbf{E}(\mathbf{r},t)|^2 + \frac{1}{\mu_0}|\mathbf{B}(\mathbf{r},t)|^2 \right) d^3r = \frac{1}{2} \sum_{\mathbf{k}, s} \hbar \omega_k \left( \alpha_{\mathbf{k},s}^*(t)\alpha_{\mathbf{k},s}(t) + \alpha_{\mathbf{k},s}(t)\alpha_{\mathbf{k},s}^*(t) \right)
\end{equation}

Definiendo las variables canónicas reales de cuadratura de posición $Q_{\mathbf{k},s}$ y momento $P_{\mathbf{k},s}$ para cada modo $(\mathbf{k}, s)$:
\begin{equation}
Q_{\mathbf{k},s} = \sqrt{\frac{\hbar}{2\omega_k}} \left( \alpha_{\mathbf{k},s} + \alpha_{\mathbf{k},s}^* \right), \qquad P_{\mathbf{k},s} = -i\sqrt{\frac{\hbar \omega_k}{2}} \left( \alpha_{\mathbf{k},s} - \alpha_{\mathbf{k},s}^* \right)
\end{equation}
el Hamiltoniano del campo libre se reduce exactamente a una suma infinita de osciladores armónicos independientes:
\begin{equation}
H_{\text{EM}} = \sum_{\mathbf{k}, s} \frac{1}{2} \left( P_{\mathbf{k},s}^2 + \omega_k^2 Q_{\mathbf{k},s}^2 \right)
\end{equation}

% =========================================================================
\section{Cuantización Canónica del Campo Electromagnético y el Concepto de Fotón}

Aplicando el procedimiento de cuantización canónica de Dirac, promovemos las cuadraturas clásicas a operadores autoadjuntos $\hat{Q}_{\mathbf{k},s}$ y $\hat{P}_{\mathbf{k},s}$ en el espacio de Hilbert que satisfacen las relaciones de conmutación a tiempos iguales:
\begin{equation}
[\hat{Q}_{\mathbf{k},s}, \hat{P}_{\mathbf{k}',s'}] = i\hbar \delta_{\mathbf{k},\mathbf{k}'} \delta_{s,s'}, \qquad [\hat{Q}_{\mathbf{k},s}, \hat{Q}_{\mathbf{k}',s'}] = [\hat{P}_{\mathbf{k},s}, \hat{P}_{\mathbf{k}',s'}] = 0
\end{equation}

En consecuencia, las amplitudes modales $\alpha_{\mathbf{k},s}$ y $\alpha_{\mathbf{k},s}^*$ se transforman en los operadores no hermíticos de **aniquilación** $\hat{a}_{\mathbf{k},s}$ y **creación** $\hat{a}_{\mathbf{k},s}^\dagger$ de fotones en el modo $(\mathbf{k}, s)$:
\begin{equation}
[\hat{a}_{\mathbf{k},s}, \hat{a}_{\mathbf{k}',s'}^\dagger] = \delta_{\mathbf{k},\mathbf{k}'} \delta_{s,s'}, \qquad [\hat{a}_{\mathbf{k},s}, \hat{a}_{\mathbf{k}',s'}] = [\hat{a}_{\mathbf{k},s}^\dagger, \hat{a}_{\mathbf{k}',s'}^\dagger] = 0
\end{equation}

\begin{definicion}[Operadores de Campo Cuántico y Hamiltoniano Cuantizado]
En la imagen de Heisenberg, los operadores de campo eléctrico y magnético se descomponen en sus partes de frecuencia positiva $\hat{\mathbf{E}}^{(+)}$ (que contiene solo operadores de aniquilación) y frecuencia negativa $\hat{\mathbf{E}}^{(-)} = [\hat{\mathbf{E}}^{(+)}]^\dagger$ (que contiene solo operadores de creación):
\begin{align}
\hat{\mathbf{E}}(\mathbf{r}, t) &= \hat{\mathbf{E}}^{(+)}(\mathbf{r}, t) + \hat{\mathbf{E}}^{(-)}(\mathbf{r}, t) \\
\hat{\mathbf{E}}^{(+)}(\mathbf{r}, t) &= i \sum_{\mathbf{k}, s} \mathcal{E}_k \bm{\varepsilon}_{\mathbf{k},s} \hat{a}_{\mathbf{k},s} e^{i(\mathbf{k}\cdot\mathbf{r} - \omega_k t)}, \qquad \mathcal{E}_k \equiv \sqrt{\frac{\hbar \omega_k}{2\varepsilon_0 V}}
\end{align}
El Hamiltoniano y el operador momento lineal total del campo electromagnético vienen dados por:
\begin{equation}
\hat{H}_{\text{EM}} = \sum_{\mathbf{k}, s} \hbar \omega_k \left( \hat{a}_{\mathbf{k},s}^\dagger \hat{a}_{\mathbf{k},s} + \frac{1}{2} \right), \qquad \hat{\mathbf{P}}_{\text{EM}} = \sum_{\mathbf{k}, s} \hbar \mathbf{k} \hat{a}_{\mathbf{k},s}^\dagger \hat{a}_{\mathbf{k},s}
\end{equation}
\end{definicion}

\begin{ejemplo}[Energía de Punto Cero del Vacío Cuántico y Efecto Casimir]
El estado fundamental del campo electromagnético es el vacío cuántico $|\{0\}\rangle = \bigotimes_{\mathbf{k},s} |0_{\mathbf{k},s}\rangle$, definido por $\hat{a}_{\mathbf{k},s}|\{0\}\rangle = 0$ para todo $(\mathbf{k}, s)$. El valor esperado del campo eléctrico en el vacío es nulo:
\[
\langle \{0\} | \hat{\mathbf{E}}(\mathbf{r}, t) | \{0\} \rangle = 0
\]
Sin embargo, el valor esperado del cuadrado del campo (fluctuaciones de vacío) es estrictamente positivo y diverge si no se introduce un corte ultravioleta:
\[
\langle \{0\} | \hat{\mathbf{E}}^2(\mathbf{r}, t) | \{0\} \rangle = \langle \{0\} | \hat{\mathbf{E}}^{(+)}(\mathbf{r}, t) \cdot \hat{\mathbf{E}}^{(-)}(\mathbf{r}, t) | \{0\} \rangle = \sum_{\mathbf{k}, s} \frac{\hbar \omega_k}{2\varepsilon_0 V}
\]
Si colocamos dos placas metálicas perfectamente conductoras, descargadas y paralelas a una distancia $d$ en el vacío, las condiciones de contorno imponen que la componente tangencial del campo eléctrico se anule en las placas, discretizando la componente $k_z = \frac{n\pi}{d}$ ($n \in \mathbb{N}$). La diferencia regularizada (usando la función Zeta de Riemann $\zeta(-3) = \frac{1}{120}$ o regularización exponencial de Euler-Maclaurin) entre la energía de punto cero $\langle \hat{H}_{\text{EM}} \rangle_0 = \frac{1}{2}\sum \hbar\omega_k$ dentro de las placas y en el espacio libre genera una fuerza atractiva neta por unidad de área $A$ conocida como **fuerza de Casimir**:
\begin{equation}
\frac{F_{\text{Casimir}}}{A} = -\frac{1}{A}\frac{\partial \Delta E_0(d)}{\partial d} = -\frac{\pi^2 \hbar c}{240 d^4}
\end{equation}
Esta fuerza macroscópica medible experimentalmente demuestra la realidad física de las fluctuaciones de punto cero $\frac{1}{2}\hbar\omega_k$ del campo electromagnético cuantizado.
\end{ejemplo}

% =========================================================================
\section{Estados Cuánticos de la Luz}

Para simplificar la notación sin pérdida de generalidad, analizaremos un único modo del campo electromagnético de frecuencia $\omega$, con operadores $\hat{a}$ y $\hat{a}^\dagger$ tales que $[\hat{a}, \hat{a}^\dagger] = 1$, operador número de fotones $\hat{n} = \hat{a}^\dagger \hat{a}$, y operadores adimensionales de cuadratura de amplitud y fase:
\begin{equation}
\hat{X}_1 = \frac{\hat{a} + \hat{a}^\dagger}{2}, \qquad \hat{X}_2 = \frac{\hat{a} - \hat{a}^\dagger}{2i}, \qquad [\hat{X}_1, \hat{X}_2] = \frac{i}{2}
\end{equation}
cuya relación de incertidumbre de Heisenberg impone $\Delta X_1 \Delta X_2 \geq \frac{1}{4}$.

\subsection{Estados Número o Estados de Fock $|n\rangle$}

Los estados de Fock $|n\rangle$ ($n \in \mathbb{N}_0$) son los autoestados del operador número $\hat{n}|n\rangle = n|n\rangle$, generados a partir del vacío mediante:
\begin{equation}
|n\rangle = \frac{(\hat{a}^\dagger)^n}{\sqrt{n!}} |0\rangle, \qquad \hat{a}|n\rangle = \sqrt{n}|n-1\rangle, \qquad \hat{a}^\dagger|n\rangle = \sqrt{n+1}|n+1\rangle
\end{equation}
Forman una base ortonormal completa del espacio de Hilbert unimodal: $\langle m|n\rangle = \delta_{mn}$ y $\sum_{n=0}^\infty |n\rangle\langle n| = \hat{\mathbb{I}}$.

\begin{propiedad}[Propiedades Estadísticas y de Fase de un Estado de Fock]
En un estado de Fock $|n\rangle$:
\begin{enumerate}
    \item El número de fotones está completamente determinado: $\langle \hat{n} \rangle = n$, con varianza nula $(\Delta n)^2 = \langle \hat{n}^2 \rangle - \langle \hat{n} \rangle^2 = 0$ (estadística **sub-poissoniana** extrema para $n \geq 1$).
    \item La fase óptica es completamente aleatoria y el valor esperado del campo eléctrico es idénticamente nulo para cualquier número de fotones:
    \[
    \langle n | \hat{E}(t) | n \rangle \propto \langle n | (\hat{a}e^{-i\omega t} - \hat{a}^\dagger e^{i\omega t}) | n \rangle = \sqrt{n}\langle n|n-1\rangle e^{-i\omega t} - \sqrt{n+1}\langle n|n+1\rangle e^{i\omega t} = 0
    \]
    \item Las fluctuaciones de cuadratura crecen linealmente con $n$:
    \[
    (\Delta X_1)^2 = (\Delta X_2)^2 = \frac{2n + 1}{4}
    \]
\end{enumerate}
\end{propiedad}

\subsection[Estados coherentes de Glauber]{Estados Coherentes de Glauber $|\alpha\rangle$}

A diferencia de los estados de Fock, un haz láser estabilizado bien por encima del umbral de oscilación tiene una fase bien definida. Roy J. Glauber formalizó estos estados en 1963.

\begin{definicion}[Tres Definiciones Equivalentes de un Estado Coherente $|\alpha\rangle$]
Para cualquier número complejo $\alpha = |\alpha|e^{i\theta} \in \mathbb{C}$, el estado coherente $|\alpha\rangle$ se define de tres maneras equivalentes:
\begin{enumerate}
    \item \textbf{Como autoestado derecho del operador de aniquilación:}
    \begin{equation}
    \hat{a}|\alpha\rangle = \alpha|\alpha\rangle, \qquad \langle \alpha | \hat{a}^\dagger = \langle \alpha | \alpha^*
    \end{equation}
    \item \textbf{Mediante el Operador Unitario de Desplazamiento de Weyl-Glauber $\hat{D}(\alpha)$ actuando sobre el vacío:}
    \begin{equation}
    |\alpha\rangle = \hat{D}(\alpha)|0\rangle, \qquad \hat{D}(\alpha) \equiv \exp\left(\alpha \hat{a}^\dagger - \alpha^* \hat{a}\right) = e^{-|\alpha|^2/2} e^{\alpha \hat{a}^\dagger} e^{-\alpha^* \hat{a}}
    \end{equation}
    donde se ha utilizado el Lema de Baker-Campbell-Hausdorff $e^{\hat{A}+\hat{B}} = e^{\hat{A}}e^{\hat{B}}e^{-[\hat{A},\hat{B}]/2}$ válido cuando $[\hat{A}, [\hat{A},\hat{B}]] = [\hat{B}, [\hat{A},\hat{B}]] = 0$.
    \item \textbf{Como superposición coherente en la base de Fock:}
    \begin{equation}
    |\alpha\rangle = e^{-|\alpha|^2/2} \sum_{n=0}^{\infty} \frac{\alpha^n}{\sqrt{n!}} |n\rangle
    \end{equation}
\end{enumerate}
\end{definicion}

\begin{ejemplo}[Cálculo de la Estadística de Fotones, Cuadraturas y Sobrecompletitud de $|\alpha\rangle$]
Analicemos exhaustivamente las propiedades físicas de un estado coherente $|\alpha\rangle$:
\begin{itemize}
    \item \textbf{Distribución de Poisson del número de fotones:} La probabilidad $P(n)$ de detectar exactamente $n$ fotones en el modo viene dada por la regla de Born:
    \begin{equation}
    P(n) = |\langle n | \alpha \rangle|^2 = \left| e^{-|\alpha|^2/2} \frac{\alpha^n}{\sqrt{n!}} \right|^2 = \frac{|\alpha|^{2n} e^{-|\alpha|^2}}{n!} = \frac{\langle \hat{n} \rangle^n e^{-\langle \hat{n} \rangle}}{n!}
    \end{equation}
    El número medio de fotones es $\langle \hat{n} \rangle = \langle \alpha | \hat{a}^\dagger \hat{a} | \alpha \rangle = |\alpha|^2$. Calculando el segundo momento usando $[\hat{a}, \hat{a}^\dagger]=1$:
    \[
    \langle \hat{n}^2 \rangle = \langle \alpha | \hat{a}^\dagger \hat{a} \hat{a}^\dagger \hat{a} | \alpha \rangle = \langle \alpha | \hat{a}^\dagger (\hat{a}^\dagger \hat{a} + 1) \hat{a} | \alpha \rangle = |\alpha|^4 + |\alpha|^2 = \langle \hat{n} \rangle^2 + \langle \hat{n} \rangle
    \]
    Por tanto, la varianza es igual a la media: $(\Delta n)^2 = \langle \hat{n} \rangle$ y la desviación estándar relativa $\frac{\Delta n}{\langle \hat{n} \rangle} = \frac{1}{\sqrt{\langle \hat{n} \rangle}} \to 0$ en el límite clásico de alta intensidad $|\alpha| \gg 1$ (ruido de disparo o \textit{shot noise}).
    \item \textbf{Estado de Mínima Incertidumbre Simétrica:} Calculando las varianzas de las cuadraturas en $|\alpha\rangle$:
    \[
    \langle \alpha | \hat{X}_1 | \alpha \rangle = \text{Re}(\alpha), \quad \langle \alpha | \hat{X}_2 | \alpha \rangle = \text{Im}(\alpha), \quad \Delta X_1 = \Delta X_2 = \frac{1}{2} \implies \Delta X_1 \Delta X_2 = \frac{1}{4}
    \]
    Un estado coherente es un paquete de ondas gaussiano desplazado en el espacio de fases cuyas fluctuaciones cuánticas son idénticas a las del estado de vacío $|0\rangle$ en todas las direcciones.
    \item \textbf{No-ortogonalidad y Relación de Clausura (Sobrecompletitud):} Dos estados coherentes distintos nunca son estrictamente ortogonales:
    \begin{equation}
    \langle \beta | \alpha \rangle = e^{-\frac{1}{2}(|\beta|^2 + |\alpha|^2 - 2\beta^* \alpha)} \implies |\langle \beta | \alpha \rangle|^2 = e^{-|\beta - \alpha|^2} > 0
    \end{equation}
    Sin embargo, forman un conjunto **sobrecompleto** que resuelve la identidad integrando sobre todo el plano complejo $d^2\alpha = d(\text{Re}\,\alpha)\,d(\text{Im}\,\alpha)$:
    \begin{equation}
    \frac{1}{\pi} \int_{\mathbb{C}} |\alpha\rangle\langle\alpha| \, d^2\alpha = \sum_{n=0}^{\infty} |n\rangle\langle n| = \hat{\mathbb{I}}
    \end{equation}
\end{itemize}
\end{ejemplo}

\subsection[Estados comprimidos]{Estados Comprimidos (\textit{Squeezed States})}

Aunque el principio de incertidumbre de Heisenberg prohíbe que el producto $\Delta X_1 \Delta X_2$ sea menor que $1/4$, nada impide reducir las fluctuaciones de una cuadratura por debajo del límite cuántico estándar del vacío ($\Delta X_1 < 1/2$) a expensas de incrementar las fluctuaciones en la cuadratura conjugada ($\Delta X_2 > 1/2$). Esta técnica es fundamental en interferometría gravitacional (como LIGO/Virgo) para superar el límite de ruido cuántico de disparo.

\begin{definicion}[Operador de Compresión (\textit{Squeezing Operator}) yVacío Comprimido]
Dado el parámetro complejo de compresión $\xi = r e^{i\theta}$ ($r \geq 0$ factor de compresión, $\theta \in [0, 2\pi)$ ángulo de compresión), el operador unitario de compresión unimodal se define como:
\begin{equation}
\hat{S}(\xi) = \exp\left[ \frac{1}{2}\left( \xi^* \hat{a}^2 - \xi (\hat{a}^\dagger)^2 \right) \right], \qquad \hat{S}^\dagger(\xi) = \hat{S}^{-1}(\xi) = \hat{S}(-\xi)
\end{equation}
Un estado coherente comprimido se obtiene aplicando compresión y luego desplazamiento al vacío: $|\alpha, \xi\rangle = \hat{D}(\alpha)\hat{S}(\xi)|0\rangle$.
\end{definicion}

Mediante el lema de trenzado de operadores $e^{\hat{A}}\hat{B}e^{-\hat{A}} = \hat{B} + [\hat{A},\hat{B}] + \frac{1}{2!}[\hat{A},[\hat{A},\hat{B}]] + \dots$, obtenemos la **Transformación Unitaria de Bogoliubov**:
\begin{align}
\hat{S}^\dagger(\xi) \hat{a} \hat{S}(\xi) &= \hat{a} \cosh r - \hat{a}^\dagger e^{i\theta} \sinh r \\
\hat{S}^\dagger(\xi) \hat{a}^\dagger \hat{S}(\xi) &= \hat{a}^\dagger \cosh r - \hat{a} e^{-i\theta} \sinh r
\end{align}

Si rotamos las cuadraturas un ángulo $\theta/2$ definiendo $\hat{Y}_1 + i\hat{Y}_2 = (\hat{X}_1 + i\hat{X}_2)e^{-i\theta/2}$, las desviaciones estándar en un estado comprimido $|\alpha, \xi\rangle$ cumplen:
\begin{equation}
\Delta Y_1 = \frac{1}{2} e^{-r} < \frac{1}{2}, \qquad \Delta Y_2 = \frac{1}{2} e^{+r} > \frac{1}{2}, \qquad \Delta Y_1 \Delta Y_2 = \frac{1}{4}
\end{equation}
En el espacio de fases $(X_1, X_2)$, el círculo de incertidumbre del estado coherente se deforma en una **elipse de incertidumbre** de igual área pero semiejes comprimidos y estirados exponencialmente por $e^{\mp r}$. En la base de Fock, el vacío comprimido $|\xi\rangle = \hat{S}(\xi)|0\rangle$ contiene únicamente un número **par** de fotones (debido a los términos cuadráticos $(\hat{a}^\dagger)^2$ que crean fotones en pares):
\begin{equation}
|\xi\rangle = \frac{1}{\sqrt{\cosh r}} \sum_{n=0}^{\infty} (-1)^n \frac{\sqrt{(2n)!}}{2^n n!} e^{in\theta} (\tanh r)^n |2n\rangle
\end{equation}

\subsection{Estados Térmicos (Luz Caótica o de Cuerpo Negro)}

Una fuente térmica en equilibrio termodinámico a temperatura $T$ no se describe mediante un estado puro $|\psi\rangle$, sino mediante un operador densidad mixto $\hat{\rho}_{\text{th}}$ que maximiza la entropía de von Neumann $S(\hat{\rho}) = -k_B \text{Tr}(\hat{\rho}\ln\hat{\rho})$ sujeto a una energía media fija $\langle \hat{H} \rangle$:
\begin{equation}
\hat{\rho}_{\text{th}} = \frac{e^{-\beta \hat{H}}}{\text{Tr}\left(e^{-\beta \hat{H}}\right)} = \left(1 - e^{-\beta\hbar\omega}\right) \sum_{n=0}^{\infty} e^{-n\beta\hbar\omega} |n\rangle\langle n|, \qquad \beta \equiv \frac{1}{k_B T}
\end{equation}

En términos del número medio de fotones térmicos dado por la **distribución de Bose-Einstein**:
\begin{equation}
\bar{n} = \langle \hat{n} \rangle_{\text{th}} = \text{Tr}(\hat{\rho}_{\text{th}} \hat{n}) = \frac{1}{e^{\hbar\omega / k_B T} - 1}
\end{equation}
la probabilidad $P_{\text{th}}(n) = \langle n | \hat{\rho}_{\text{th}} | n \rangle$ de encontrar $n$ fotones sigue la **distribución geométrica de Bose-Einstein**:
\begin{equation}
P_{\text{th}}(n) = \frac{\bar{n}^n}{(1 + \bar{n})^{n+1}}
\end{equation}
cuya varianza es $(\Delta n)^2_{\text{th}} = \bar{n}^2 + \bar{n} > \bar{n}$ (**estadística super-poissoniana**), reflejando las fluctuaciones térmicas de intensidad propias de las ondas caóticas.

% =========================================================================
\section[Representaciones cuánticas de fase]{Representaciones en el Espacio de Fases Cuántico: Cuasi-Probabilidades $P$, $Q$ y Wigner $W$}

En mecánica estadística clásica, cualquier estado se describe mediante una distribución de probabilidad conjunta positiva $f(q, p) \geq 0$ en el espacio de fases. En mecánica cuántica, la no conmutatividad $[\hat{a}, \hat{a}^\dagger]=1$ impide definir una probabilidad genuina simultánea para ambas cuadraturas, dando lugar a **distribuciones de cuasi-probabilidad** asociadas al ordenamiento de operadores:

\begin{enumerate}
    \item \textbf{Representación $P(\alpha)$ de Glauber-Sudarshan (Orden Normal $\hat{a}^\dagger \hat{a}$):}
    Expresa cualquier matriz densidad $\hat{\rho}$ como una mezcla diagonal de proyectores sobre estados coherentes:
    \begin{equation}
    \hat{\rho} = \int_{\mathbb{C}} P(\alpha) |\alpha\rangle\langle\alpha| \, d^2\alpha, \qquad \int_{\mathbb{C}} P(\alpha) d^2\alpha = 1
    \end{equation}
    Para un estado coherente $|\alpha_0\rangle$, $P(\alpha) = \delta^{(2)}(\alpha - \alpha_0)$. Para luz térmica, $P(\alpha) = \frac{1}{\pi\bar{n}}e^{-|\alpha|^2/\bar{n}} > 0$. Para estados genuinamente cuánticos (como estados de Fock $|n\rangle$ o comprimidos), $P(\alpha)$ toma valores negativos o se vuelve más singular que una delta de Dirac (contiene derivadas de orden superior $\partial_\alpha^k \partial_{\alpha^*}^k \delta^{(2)}(\alpha)$), lo que constituye el **criterio riguroso de no-clasicidad de la luz**.
    
    \item \textbf{Función $Q(\alpha)$ de Husimi-Kano (Orden Antinormal $\hat{a} \hat{a}^\dagger$):}
    Se define como el valor esperado de $\hat{\rho}$ en un estado coherente normalizado por $\pi$:
    \begin{equation}
    Q(\alpha) = \frac{1}{\pi} \langle \alpha | \hat{\rho} | \alpha \rangle
    \end{equation}
    Siempre está bien comportada, es acotada y estrictamente no negativa ($0 \leq Q(\alpha) \leq \frac{1}{\pi}$), pero ensancha las distribuciones debido al ruido de punto cero del proyectil $|\alpha\rangle$.
    
    \item \textbf{Función de Wigner $W(\alpha)$ (Orden Simétrico de Weyl):}
    Es la transformada de Fourier bidimensional de la función característica simétrica $\chi_W(\lambda) = \text{Tr}\left[\hat{\rho} \exp(\lambda \hat{a}^\dagger - \lambda^* \hat{a})\right]$:
    \begin{equation}
    W(\alpha) = \frac{1}{\pi^2} \int_{\mathbb{C}} \text{Tr}\left[ \hat{\rho} \hat{D}(\lambda) \right] e^{\alpha\lambda^* - \alpha^*\lambda} \, d^2\lambda
    \end{equation}
    Sus distribuciones marginales reproducen exactamente las probabilidades cuánticas de cualquier cuadratura medida mediante detección homodina equilibrada. La presencia de **regiones negativas** ($W(\alpha) < 0$), como ocurre en el origen para los estados de Fock impares ($W_{|n\rangle}(0) = \frac{2}{\pi}(-1)^n$), es una firma directa de interferencia cuántica en el espacio de fases.
\end{enumerate}

% =========================================================================
\section[Coherencia cuántica y fotoconteo]{Teoría de la Coherencia Cuántica de Glauber y Estadística de Fotoconteo}

La teoría de detección óptica de Glauber se basa en que los fotodetectores operan por absorción de fotones (efecto fotoeléctrico), la cual está gobernada por la parte de frecuencia positiva del campo $\hat{E}^{(+)}(\mathbf{r},t)$.

\begin{definicion}[Funciones de Correlación de Primer y Segundo Orden de Glauber]
Para un campo estacionario, las funciones de correlación normalizadas de primer orden $g^{(1)}(\tau)$ (coherencia de amplitud/fase medida en interferómetros de Michelson o Mach-Zehnder) y de segundo orden $g^{(2)}(\tau)$ (coherencia de intensidad medida en el interferómetro de Hanbury Brown y Twiss - HBT) se definen en orden normal como:
\begin{align}
g^{(1)}(\tau) &= \frac{\langle \hat{E}^{(-)}(t) \hat{E}^{(+)}(t+\tau) \rangle}{\langle \hat{E}^{(-)}(t) \hat{E}^{(+)}(t) \rangle} = \frac{\langle \hat{a}^\dagger(t) \hat{a}(t+\tau) \rangle}{\langle \hat{a}^\dagger(t) \hat{a}(t) \rangle} \\
g^{(2)}(\tau) &= \frac{\langle \hat{E}^{(-)}(t) \hat{E}^{(-)}(t+\tau) \hat{E}^{(+)}(t+\tau) \hat{E}^{(+)}(t) \rangle}{\langle \hat{E}^{(-)}(t) \hat{E}^{(+)}(t) \rangle^2}
\end{align}
A retraso temporal nulo ($\tau = 0$) para un único modo, $g^{(2)}(0)$ se reduce a:
\begin{equation}
g^{(2)}(0) = \frac{\langle \hat{a}^\dagger \hat{a}^\dagger \hat{a} \hat{a} \rangle}{\langle \hat{a}^\dagger \hat{a} \rangle^2} = \frac{\langle \hat{n}(\hat{n} - 1) \rangle}{\langle \hat{n} \rangle^2} = 1 + \frac{(\Delta n)^2 - \langle \hat{n} \rangle}{\langle \hat{n} \rangle^2} = 1 + \frac{Q_M}{\langle \hat{n} \rangle}
\end{equation}
donde $Q_M \equiv \frac{(\Delta n)^2 - \langle \hat{n} \rangle}{\langle \hat{n} \rangle}$ es el **parámetro de Mandel**.
\end{definicion}

\begin{teorema}[Límites Clásicos vs. Cuánticos de $g^{(2)}(\tau)$]
Para cualquier campo electromagnético clásico descrito por una intensidad estocástica positiva $I(t)$ (o equivalentemente, cuya representación $P(\alpha) \geq 0$ sea una distribución de probabilidad genuina), la desigualdad de Cauchy-Schwarz impone dos restricciones fundamentales:
\begin{equation}
\text{(1) } g^{(2)}(0) \geq 1, \qquad \text{(2) } g^{(2)}(0) \geq g^{(2)}(\tau) \quad \forall \tau
\end{equation}
La violación de la primera condición ($g^{(2)}(0) < 1 \iff Q_M < 0$) se denomina **estadística sub-poissoniana**, mientras que la violación de la segunda condición ($g^{(2)}(0) < g^{(2)}(\tau)$ para $\tau > 0$) se denomina **antiagrupamiento de fotones (\textit{photon antibunching})**. Ambos fenómenos son estrictamente imposibles en electrodinámica clásica y prueban inequívocamente la naturaleza corpuscular cuantizada del campo.
\end{teorema}

\begin{ejemplo}[Comparación de $g^{(2)}(0)$ para Luz Térmica, Coherente y de Fock]
Evaluemos $g^{(2)}(0)$ para los tres estados arquetípicos:
\begin{enumerate}
    \item \textbf{Luz Térmica o Caótica ($\hat{\rho}_{\text{th}}$):} Como $(\Delta n)^2 = \bar{n}^2 + \bar{n}$, tenemos:
    \[
    g^{(2)}_{\text{th}}(0) = 1 + \frac{\bar{n}^2 + \bar{n} - \bar{n}}{\bar{n}^2} = 2 \quad (\text{Agrupamiento de fotones o \textit{photon bunching}: } g^{(2)}(\tau) = 1 + |g^{(1)}(\tau)|^2)
    \]
    Los fotones térmicos tienden a llegar en ráfagas o pares al detector con el doble de probabilidad a $\tau=0$ que a tiempos largos $\tau \gg \tau_c$.
    \item \textbf{Luz Coherente Láser ($|\alpha\rangle$):} Como $\hat{a}|\alpha\rangle = \alpha|\alpha\rangle$:
    \[
    g^{(2)}_{\text{coh}}(0) = \frac{\langle \alpha | (\hat{a}^\dagger)^2 \hat{a}^2 | \alpha \rangle}{\langle \alpha | \hat{a}^\dagger \hat{a} | \alpha \rangle^2} = \frac{|\alpha|^4}{|\alpha|^4} = 1 \quad (\text{Fotones estadísticamente independientes, } Q_M = 0)
    \]
    \item \textbf{Estado de Fock $|n\rangle$ ($n \geq 1$):} Como $(\Delta n)^2 = 0$:
    \[
    g^{(2)}_{|n\rangle}(0) = \frac{n(n-1)}{n^2} = 1 - \frac{1}{n} < 1 \quad (\text{Luz no clásica sub-poissoniana, } Q_M = -1)
    \]
    En particular, para una fuente de **un solo fotón** ($|1\rangle$, como un átomo individual en fluorescencia de resonancia o un punto cuántico semiconductor), obtenemos el límite cuántico absoluto:
    \[
    g^{(2)}_{|1\rangle}(0) = 0
    \]
    Un fotón individual no puede dividirse en dos en un divisor de haz para activar simultáneamente dos fotodetectores a $\tau=0$ (experimento de Grangier, Roger y Aspect, 1986).
\end{enumerate}
\end{ejemplo}

% =========================================================================
\section[Efecto Hong-Ou-Mandel]{Interferencia de Dos Fotones Indistinguibles: El Efecto Hong-Ou-Mandel (HOM)}

El divisor de haz (\textit{beam splitter}, BS) sin pérdidas es el elemento fundamental de la óptica cuántica lineal. Tiene dos modos espaciales de entrada con operadores $\hat{a}_1, \hat{a}_2$ y dos modos de salida con operadores $\hat{b}_1, \hat{b}_2$. La conservación de la energía ($\hat{b}_1^\dagger \hat{b}_1 + \hat{b}_2^\dagger \hat{b}_2 = \hat{a}_1^\dagger \hat{a}_1 + \hat{a}_2^\dagger \hat{a}_2$) y las relaciones de conmutación bosónicas $[\hat{b}_i, \hat{b}_j^\dagger] = \delta_{ij}$ exigen que la matriz de transformación sea unitaria ($\mathbf{U} \in SU(2)$). Para un divisor de haz simétrico 50:50 ($t = 1/\sqrt{2}$, $r = i/\sqrt{2}$):
\begin{equation}
\begin{pmatrix} \hat{b}_1 \\ \hat{b}_2 \end{pmatrix} = \frac{1}{\sqrt{2}} \begin{pmatrix} 1 & i \\ i & 1 \end{pmatrix} \begin{pmatrix} \hat{a}_1 \\ \hat{a}_2 \end{pmatrix} \implies \begin{cases} \hat{a}_1^\dagger = \frac{1}{\sqrt{2}}(\hat{b}_1^\dagger - i\hat{b}_2^\dagger) \\ \hat{a}_2^\dagger = \frac{1}{\sqrt{2}}(-i\hat{b}_1^\dagger + \hat{b}_2^\dagger) \end{cases}
\end{equation}

\begin{ejemplo}[Demostración Algebraica del Valle de Hong-Ou-Mandel (1987)]
Supongamos que inyectamos simultáneamente dos fotones exactamente indistinguibles (misma frecuencia, polarización, perfil transversal y tiempo de llegada), uno por cada puerto de entrada del divisor de haz 50:50:
\[
|\psi_{\text{in}}\rangle = |1_1, 1_2\rangle_{\text{in}} = \hat{a}_1^\dagger \hat{a}_2^\dagger |0, 0\rangle
\]
Sustituyendo la relación inversa de los operadores de creación en términos de los modos de salida:
\begin{align}
|\psi_{\text{out}}\rangle &= \frac{1}{2} \left( \hat{b}_1^\dagger - i\hat{b}_2^\dagger \right) \left( -i\hat{b}_1^\dagger + \hat{b}_2^\dagger \right) |0, 0\rangle \\
&= \frac{1}{2} \left[ -i(\hat{b}_1^\dagger)^2 + \hat{b}_1^\dagger \hat{b}_2^\dagger - \hat{b}_2^\dagger \hat{b}_1^\dagger - i(\hat{b}_2^\dagger)^2 \right] |0, 0\rangle
\end{align}
Dado que los fotones son bosones, sus operadores de creación en distintos modos conmutan: $[\hat{b}_1^\dagger, \hat{b}_2^\dagger] = 0 \implies \hat{b}_1^\dagger \hat{b}_2^\dagger = \hat{b}_2^\dagger \hat{b}_1^\dagger$. ¡Los dos caminos cuánticos donde ambos fotones se transmiten ($\hat{b}_1^\dagger \hat{b}_2^\dagger$) o ambos se reflejan ($-\hat{b}_2^\dagger \hat{b}_1^\dagger$) interfieren destructivamente y se cancelan de forma exacta! Usando $(\hat{b}_j^\dagger)^2|0\rangle = \sqrt{2}|2_j\rangle$, el estado de salida resulta ser un estado entrelazado NOON con $N=2$:
\begin{equation}
|\psi_{\text{out}}\rangle = -\frac{i}{\sqrt{2}} \left( |2_1, 0_2\rangle_{\text{out}} + |0_1, 2_2\rangle_{\text{out}} \right)
\end{equation}
La probabilidad de coincidencia $P_{\text{coinc}}(1_1, 1_2) = |\langle 1_1, 1_2 | \psi_{\text{out}} \rangle|^2 = 0$ cae a cero (el célebre **valle o \textit{dip} de Hong-Ou-Mandel**). Ambos fotones salen siempre juntos (*coalescencia bosónica*) por el puerto 1 o por el puerto 2 con igual probabilidad del $50\%$.
\end{ejemplo}

% =========================================================================
\section[Interacción radiación-materia]{Interacción Radiación-Materia y QED de Cavidades}

Consideremos la interacción fundamental entre un átomo de dos niveles —con estado fundamental $|g\rangle$, estado excitado $|e\rangle$ y frecuencia de transición $\omega_0 = (E_e - E_g)/\hbar$— y un único modo cuantizado de una cavidad electromagnética de alto factor de calidad $Q$ con frecuencia $\omega_c$.

Definimos los operadores atómicos de pseudoespín (matrices de Pauli atómicas y operadores de subida/bajada):
\begin{equation}
\hat{\sigma}_z = |e\rangle\langle e| - |g\rangle\langle g|, \qquad \hat{\sigma}_+ = |e\rangle\langle g|, \qquad \hat{\sigma}_- = |g\rangle\langle e|, \qquad [\hat{\sigma}_+, \hat{\sigma}_-] = \hat{\sigma}_z
\end{equation}

En la aproximación dipolar eléctrica ($k_c a_0 \ll 1$), el Hamiltoniano de interacción completo (modelo de Rabi cuántico) es $\hat{H}_{\text{int}} = -\hat{\mathbf{d}}\cdot\hat{\mathbf{E}} = \hbar g (\hat{\sigma}_+ + \hat{\sigma}_-)(\hat{a} + \hat{a}^\dagger)$, donde la constante de acoplamiento vacío-átomo (frecuencia de Rabi de un solo fotón) es $g = -\frac{\mathbf{d}_{eg}\cdot\bm{\varepsilon}}{\hbar}\sqrt{\frac{\hbar\omega_c}{2\varepsilon_0 V}}$. Pasando a la imagen de interacción, los términos $\hat{\sigma}_+ \hat{a} e^{i(\omega_0 - \omega_c)t}$ y $\hat{\sigma}_- \hat{a}^\dagger e^{-i(\omega_0 - \omega_c)t}$ conservan la energía en resonancia, mientras que los términos contrarrotantes $\hat{\sigma}_+ \hat{a}^\dagger e^{i(\omega_0 + \omega_c)t}$ y $\hat{\sigma}_- \hat{a} e^{-i(\omega_0 + \omega_c)t}$ oscilan rápidamente a $\sim 2\omega_0$ y violan la conservación de excitaciones. Bajo la **Aproximación de Onda Rotante (RWA)**, válida cuando $|g| \ll \omega_0 + \omega_c$ y $|\Delta| = |\omega_0 - \omega_c| \ll \omega_0 + \omega_c$, descartamos los términos contrarrotantes.

\begin{definicion}[Hamiltoniano de Jaynes-Cummings (1963)]
El Hamiltoniano que gobierna la electrodinámica cuántica de cavidades en régimen de acoplamiento fuerte bajo la RWA es:
\begin{equation}
\hat{H}_{\text{JC}} = \frac{1}{2}\hbar\omega_0 \hat{\sigma}_z + \hbar\omega_c \hat{a}^\dagger \hat{a} + \hbar g \left( \hat{\sigma}_+ \hat{a} + \hat{\sigma}_- \hat{a}^\dagger \right)
\end{equation}
\end{definicion}

El operador número total de excitaciones $\hat{N}_{\text{exc}} = \hat{a}^\dagger\hat{a} + |e\rangle\langle e|$ conmuta con el Hamiltoniano: $[\hat{H}_{\text{JC}}, \hat{N}_{\text{exc}}] = 0$. En consecuencia, el espacio de Hilbert $\mathcal{H}_{\text{átomo}} \otimes \mathcal{H}_{\text{campo}}$ se descompone en una suma directa de subespacios invariantes bidimensionales $\mathcal{H}^{(n)} = \text{span}\{|e, n\rangle, |g, n+1\rangle\}$ para cada $n \geq 0$ (más el estado fundamental no acoplado $|g, 0\rangle$).

En la base $\{|e, n\rangle, |g, n+1\rangle\}$, el bloque matricial $2\times 2$ del Hamiltoniano de Jaynes-Cummings es:
\begin{equation}
\hat{H}_{\text{JC}}^{(n)} = \hbar\omega_c\left(n + \frac{1}{2}\right)\mathbb{I}_2 + \frac{\hbar}{2} \begin{pmatrix} \Delta & 2g\sqrt{n+1} \\ 2g\sqrt{n+1} & -\Delta \end{pmatrix}, \qquad \Delta \equiv \omega_0 - \omega_c
\end{equation}

Diagonalizando esta matriz, obtenemos los autovalores de energía y los autoestados entrelazados átomo-fotón llamados **estados vestidos (\textit{dressed states}) o polaritones de cavidad** $|+, n\rangle$ y $|-, n\rangle$:
\begin{equation}
E_{\pm, n} = \hbar\omega_c\left(n + \frac{1}{2}\right) \pm \frac{1}{2}\hbar\Omega_n(\Delta), \qquad \Omega_n(\Delta) = \sqrt{\Delta^2 + 4g^2(n+1)}
\end{equation}
donde $\Omega_n(\Delta)$ es la **frecuencia de Rabi cuántica generalizada** para $n$ fotones. En resonancia exacta ($\Delta = 0$), el desdoblamiento de niveles (\textit{Vacuum Rabi Splitting} para $n=0$) escala de forma no lineal con $\sqrt{n+1}$ ($\Omega_n(0) = 2g\sqrt{n+1}$), una firma directa de la discretización del campo en fotones:
\begin{equation}
|\pm, n\rangle = \frac{1}{\sqrt{2}} \left( |e, n\rangle \pm |g, n+1\rangle \right) \qquad (\text{para } \Delta = 0)
\end{equation}

\begin{ejemplo}[Oscilaciones de Rabi de Vacío y Colapsos y Resurgimientos Cuánticos (\textit{Collapses and Revivals})]
\textbf{Caso 1 (Emisión espontánea reversible en el vacío $|e, 0\rangle$):} Si el átomo entra inicialmente excitado en una cavidad perfectamente vacía $|\psi(0)\rangle = |e, 0\rangle$ en resonancia ($\Delta = 0$), el estado a tiempo $t$ evoluciona unitariamente como:
\begin{equation}
|\psi(t)\rangle = \cos(gt)|e, 0\rangle - i\sin(gt)|g, 1\rangle
\end{equation}
La probabilidad de encontrar el átomo excitado oscila sinusoidalmente como $P_e(t) = \cos^2(gt) = \frac{1}{2}[1 + \cos(2gt)]$. ¡Incluso en ausencia inicial de fotones ($n=0$), las fluctuaciones del vacío cuántico inducen oscilaciones coherentes reversibles a la frecuencia de Rabi de vacío $2g$! En el instante $t = \frac{\pi}{4g}$, el sistema alcanza el estado máximamente entrelazado átomo-fotón $|\psi\rangle = \frac{1}{\sqrt{2}}(|e,0\rangle - i|g,1\rangle)$.

\textbf{Caso 2 (Campo inicialmente en un estado coherente $|\alpha\rangle$):} Si $|\psi(0)\rangle = |e\rangle \otimes |\alpha\rangle = \sum_{n=0}^\infty c_n |e, n\rangle$ con $|c_n|^2 = e^{-\bar{n}}\frac{\bar{n}^n}{n!}$ ($\bar{n} = |\alpha|^2 \gg 1$), la inversión atómica evoluciona como una suma discreta de frecuencias inconmensurables $\Omega_n = 2g\sqrt{n+1}$:
\begin{equation}
W(t) \equiv \langle \hat{\sigma}_z(t) \rangle = e^{-\bar{n}} \sum_{n=0}^{\infty} \frac{\bar{n}^n}{n!} \cos\left(2gt\sqrt{n+1}\right)
\end{equation}
A tiempos cortos $t_{\text{col}} \sim \frac{1}{g}$, las distintas componentes de Fock se desfasan destructivamente provocando el **colapso de Cummings** ($W(t) \to 0$). Sin embargo, debido a la naturaleza discreta de los números enteros $n$ en la distribución de fotones, cuando los términos vecinos en torno a $\bar{n}$ acumulan una diferencia de fase de $2\pi$:
\[
\left( 2g\sqrt{\bar{n}+1} - 2g\sqrt{\bar{n}} \right) T_{\text{rev}} \approx \frac{g}{\sqrt{\bar{n}}} T_{\text{rev}} = 2\pi \implies T_{\text{rev}} \approx \frac{2\pi\sqrt{\bar{n}}}{g}
\]
todas las fases se realinean constructivamente produciendo un **resurgimiento cuántico (\textit{quantum revival})** de las oscilaciones atómicas, fenómeno observado experimentalmente por Rempe, Walther y Klein (1987) y por el grupo de Serge Haroche (Premio Nobel 2012).
\end{ejemplo}

% =========================================================================
\section[Sistemas abiertos y decoherencia]{Sistemas Cuánticos Abiertos, Decoherencia y la Ecuación Maestra de Lindblad}

En cualquier experimento real de óptica cuántica, la cavidad pierde fotones a través de sus espejos con una tasa $\kappa$ y el átomo decae espontáneamente hacia modos del espacio libre con una tasa $\gamma$. El sistema de interés $S$ interactúa con un reservorio o baño térmico $B$ de infinitos grados de libertad a través del Hamiltoniano total $\hat{H}_{\text{tot}} = \hat{H}_S + \hat{H}_B + \hat{H}_{SB}$.

Trazando parcialmente sobre el baño ($\hat{\rho}_S(t) = \text{Tr}_B[\hat{\rho}_{\text{tot}}(t)]$) y aplicando la **Aproximación de Born** (acoplamiento débil sistema-baño, $\hat{\rho}_{\text{tot}}(t) \approx \hat{\rho}_S(t) \otimes \hat{\rho}_B(0)$) y la **Aproximación de Markov** (el tiempo de correlación de memoria del baño $\tau_B$ es infinitesimal comparado con la escala de evolución del sistema, $\tau_B \ll \gamma^{-1}, \kappa^{-1}$), Gorini, Kossakowski, Sudarshan y Lindblad demostraron que el generador más general de un semigrupo dinámico cuántico completamente positivo y que preserva la traza (CPTP) tiene la forma siguiente.

\begin{teorema}[Ecuación Maestra de Gorini-Kossakowski-Sudarshan-Lindblad (GKSL)]
La evolución temporal no unitaria de la matriz densidad reducida $\hat{\rho}$ de un sistema cuántico abierto markoviano obedece:
\begin{equation}
\frac{d\hat{\rho}}{dt} = -\frac{i}{\hbar}[\hat{H}_S, \hat{\rho}] + \sum_{k} \gamma_k \mathcal{D}[\hat{L}_k]\hat{\rho}
\end{equation}
donde $\gamma_k \geq 0$ son las tasas de relajación o decoherencia, $\hat{L}_k$ son los **operadores de salto cuántico de Lindblad**, y el superoperador disipador $\mathcal{D}[\hat{L}]$ actúa como:
\begin{equation}
\mathcal{D}[\hat{L}]\hat{\rho} \equiv \hat{L}\hat{\rho}\hat{L}^\dagger - \frac{1}{2}\left\{ \hat{L}^\dagger \hat{L}, \hat{\rho} \right\} = \hat{L}\hat{\rho}\hat{L}^\dagger - \frac{1}{2}\hat{L}^\dagger \hat{L}\hat{\rho} - \frac{1}{2}\hat{\rho}\hat{L}^\dagger \hat{L}
\end{equation}
\end{teorema}

\begin{ejemplo}[Ecuación Maestra de una Cavidad con Pérdidas en un Baño Térmico y Decoherencia de Estados Gato de Schrödinger]
Para un modo de cavidad acoplado a un baño térmico con número medio de fotones térmicos $\bar{n}_{\text{th}}$ y tasa de amortiguamiento $\kappa$, existen dos canales de salto de Lindblad: pérdida de fotones ($\hat{L}_1 = \hat{a}$ con tasa $\kappa(\bar{n}_{\text{th}} + 1)$) e inyección térmica de fotones ($\hat{L}_2 = \hat{a}^\dagger$ con tasa $\kappa \bar{n}_{\text{th}}$):
\begin{equation}
\frac{d\hat{\rho}}{dt} = -i\omega_c [\hat{a}^\dagger \hat{a}, \hat{\rho}] + \kappa(\bar{n}_{\text{th}} + 1)\left( \hat{a}\hat{\rho}\hat{a}^\dagger - \frac{1}{2}\{\hat{a}^\dagger\hat{a}, \hat{\rho}\} \right) + \kappa\bar{n}_{\text{th}}\left( \hat{a}^\dagger\hat{\rho}\hat{a} - \frac{1}{2}\{\hat{a}\hat{a}^\dagger, \hat{\rho}\} \right)
\end{equation}
Si preparamos a $T=0$ ($\bar{n}_{\text{th}}=0$) una superposición mesoscópica de dos estados coherentes con fases opuestas —un **estado Gato de Schrödinger óptico** $|\psi_{\text{cat}}\rangle = \mathcal{N}(|\alpha_0\rangle + |-\alpha_0\rangle)$—, los términos fuera de la diagonal de coherencia cuántica $|\alpha_0\rangle\langle -\alpha_0|$ en la matriz densidad decaen con una tasa de **decoherencia** $\Gamma_{\text{dec}}$ proporcional al cuadrado de la distancia en el espacio de fases:
\begin{equation}
\Gamma_{\text{dec}} = 2 \kappa |\alpha_0|^2 = 2 \kappa \langle \hat{n} \rangle
\end{equation}
Mientras que la energía de la cavidad se relaja lentamente en la escala $\tau_{\text{rel}} = \kappa^{-1}$, la superposición cuántica se destruye casi instantáneamente en el tiempo $\tau_{\text{dec}} = \frac{1}{2\kappa \langle \hat{n} \rangle} \ll \tau_{\text{rel}}$ tan pronto como un solo fotón escapa al entorno revelando información de "cuál fase" tenía el campo, transformando la superposición coherente en una mezcla estadística clásica $\hat{\rho}_{\text{mix}} = \frac{1}{2}(|\alpha_0 e^{-\kappa t/2}\rangle\langle\alpha_0 e^{-\kappa t/2}| + |-\alpha_0 e^{-\kappa t/2}\rangle\langle -\alpha_0 e^{-\kappa t/2}|)$.
\end{ejemplo}

% =========================================================================
% CAPÍTULO 4: ENTRELAZAMIENTO, PARADOJA EPR Y DESIGUALDADES DE BELL
% =========================================================================
\chapter[Entrelazamiento cuántico]{Entrelazamiento Cuántico, No-Localidad y Desigualdades de Bell}

\section[Paradoja EPR y realismo local]{La Paradoja de Einstein-Podolsky-Rosen (EPR, 1935) y el Realismo Local}

En su histórico artículo de 1935 (\textit{"Can Quantum-Mechanical Description of Physical Reality Be Considered Complete?"}), Albert Einstein, Boris Podolsky y Nathan Rosen formularon una crítica profunda a la completitud de la mecánica cuántica basándose en tres premisas filosófico-físicas que hoy denominamos **Realismo Local**:

\begin{enumerate}
    \item \textbf{Criterio de Realidad Física de EPR:} \textit{"Si, sin perturbar de modo alguno un sistema, podemos predecir con certeza (es decir, con probabilidad igual a la unidad) el valor de una cantidad física, entonces existe un elemento de realidad física correspondiente a esta cantidad física."}
    \item \textbf{Principio de Localidad Relativista (Causalidad de Einstein):} Ninguna acción o medición ejecutada en una región espacial $A$ puede alterar instantáneamente los elementos de realidad física de un sistema situado en una región $B$ separada por un intervalo relativista de tipo espacio ($\Delta s^2 = c^2\Delta t^2 - |\Delta\mathbf{r}|^2 < 0$).
    \item \textbf{Condición de Completitud:} \textit{"Todo elemento de la realidad física debe tener una contrapartida en la teoría física completa."}
\end{enumerate}

En la formulación de **David Bohm (1951)** para sistemas discretos de espín $1/2$ (o polarización de dos fotones), consideremos una fuente que emite un par de partículas hacia dos observadores distantes, Alice ($A$) y Bob ($B$), en el **estado singlete antisimétrico de Bell** $|\Psi^-\rangle$:
\begin{equation}
|\Psi^-\rangle = \frac{1}{\sqrt{2}} \left( |\uparrow_{\mathbf{n}}\rangle_A \otimes |\downarrow_{\mathbf{n}}\rangle_B - |\downarrow_{\mathbf{n}}\rangle_A \otimes |\uparrow_{\mathbf{n}}\rangle_B \right)
\end{equation}
Debido a que $|\Psi^-\rangle$ tiene espín total $S=0$, es rotacionalmente invariante ($(\hat{U}_R \otimes \hat{U}_R)|\Psi^-\rangle = |\Psi^-\rangle$) y mantiene exactamente la misma forma matemática independientemente del eje unitario $\mathbf{n} \in \mathbb{R}^3$ elegido.

Si Alice decide medir la componente de espín a lo largo del eje $\hat{\mathbf{z}}$ ($\hat{\sigma}_z^A$) y obtiene $+1$, el estado de Bob colapsa instantáneamente a $|\downarrow_z\rangle_B$, permitiendo a Alice predecir con $100\%$ de certeza que Bob obtendrá $\sigma_z^B = -1$. Por localidad (Bob está a años luz de distancia), la medición de Alice no pudo perturbar a Bob; luego $\sigma_z^B$ es un elemento de realidad preexistente. Pero Alice podría haber optado por medir a lo largo del eje $\hat{\mathbf{x}}$ ($\hat{\sigma}_x^A$), prediciendo con certeza el valor de $\sigma_x^B$, por lo que $\sigma_x^B$ también debe ser un elemento de realidad simultáneo. Como en mecánica cuántica $[\hat{\sigma}_z^B, \hat{\sigma}_x^B] = 2i\hat{\sigma}_y^B \neq 0$ no poseen autoestados comunes, EPR concluyeron que la función de onda cuántica es **incompleta** y debe ser suplementada por **variables ocultas locales** $\lambda$.

% =========================================================================
\section[Formalismo del entrelazamiento]{Formalismo General del Entrelazamiento}

\begin{definicion}[Estados Separables vs. Estados Entrelazados]
Sea un sistema compuesto bipartito con espacio de Hilbert $\mathcal{H}_{AB} = \mathcal{H}_A \otimes \mathcal{H}_B$:
\begin{itemize}
    \item Un **estado puro** $|\Psi\rangle_{AB} \in \mathcal{H}_{AB}$ es **separable (factorizable)** si y solo si puede escribirse como un único producto tensorial $|\Psi\rangle_{AB} = |\phi\rangle_A \otimes |\chi\rangle_B$. En caso contrario, se dice que está **entrelazado**.
    \item Un **estado mixto** general $\hat{\rho}_{AB}$ es **separable** (Werner, 1989) si y solo si puede expresarse como una combinación convexa de estados producto con probabilidades clásicas $p_k \geq 0$ ($\sum_k p_k = 1$):
    \begin{equation}
    \hat{\rho}_{AB}^{\text{sep}} = \sum_{k} p_k \, \hat{\rho}_A^{(k)} \otimes \hat{\rho}_B^{(k)}
    \end{equation}
    Si no admite tal descomposición, $\hat{\rho}_{AB}$ está **entrelazado**.
\end{itemize}
\end{definicion}

\begin{definicion}[La Base de Estados de Bell para Dos Qubits]
En el espacio $\mathbb{C}^2 \otimes \mathbb{C}^2$ (donde $|0\rangle \equiv |H\rangle$ o $|\uparrow\rangle$, y $|1\rangle \equiv |V\rangle$ o $|\downarrow\rangle$), los cuatro estados máximamente entrelazados de Bell forman una base ortonormal completa:
\begin{align}
|\Phi^+\rangle &= \frac{1}{\sqrt{2}} \left( |00\rangle + |11\rangle \right), & |\Phi^-\rangle &= \frac{1}{\sqrt{2}} \left( |00\rangle - |11\rangle \right) \\
|\Psi^+\rangle &= \frac{1}{\sqrt{2}} \left( |01\rangle + |10\rangle \right), & |\Psi^-\rangle &= \frac{1}{\sqrt{2}} \left( |01\rangle - |10\rangle \right)
\end{align}
Cualquiera de los cuatro estados de Bell tiene estados reducidos máximamente mixtos: $\hat{\rho}_A = \text{Tr}_B(|\text{Bell}\rangle\langle\text{Bell}|) = \frac{1}{2}\hat{\mathbb{I}}_2 = \hat{\rho}_B$.
\end{definicion}

\begin{teorema}[Descomposición de Schmidt y Entropía de Entrelazamiento]
Para cualquier estado puro bipartito $|\Psi\rangle_{AB} \in \mathcal{H}_A \otimes \mathcal{H}_B$ con $d = \min(\dim\mathcal{H}_A, \dim\mathcal{H}_B)$, existen bases ortonormales $\{|u_i\rangle_A\}$ y $\{|v_i\rangle_B\}$ tales que $|\Psi\rangle_{AB}$ admite la **Descomposición de Schmidt** con una única suma:
\begin{equation}
|\Psi\rangle_{AB} = \sum_{i=1}^{r_S} \sqrt{\lambda_i} \, |u_i\rangle_A \otimes |v_i\rangle_B, \qquad \lambda_i > 0, \quad \sum_{i=1}^{r_S} \lambda_i = 1
\end{equation}
donde los $\lambda_i$ son los autovalores no nulos comunes de $\hat{\rho}_A$ y $\hat{\rho}_B$, y el entero $r_S \leq d$ es el **rango de Schmidt**. El estado $|\Psi\rangle_{AB}$ es separable si y solo si $r_S = 1$, y está entrelazado si y solo si $r_S \geq 2$. La medida canónica de entrelazamiento para estados puros es la **entropía de entrelazamiento de von Neumann**:
\begin{equation}
E(|\Psi\rangle_{AB}) = S(\hat{\rho}_A) = S(\hat{\rho}_B) = -\sum_{i=1}^{r_S} \lambda_i \log_2 \lambda_i \in [0, \log_2 d]
\end{equation}
\end{teorema}

\begin{teorema}[Criterio de Transposición Parcial Positiva (PPT) de Peres-Horodecki]
Sea $\hat{\rho}_{AB}$ una matriz densidad en la base producto $\{|i\rangle_A |j\rangle_B\}$ con elementos $(\hat{\rho}_{AB})_{ij, kl} = \langle i, j | \hat{\rho}_{AB} | k, l \rangle$. Definimos la **transpuesta parcial** respecto al subsistema $B$, denotada $\hat{\rho}_{AB}^{T_B}$, intercambiando únicamente los índices de Bob:
\begin{equation}
\left(\hat{\rho}_{AB}^{T_B}\right)_{ij, kl} \equiv \left(\hat{\rho}_{AB}\right)_{il, kj}
\end{equation}
\begin{itemize}
    \item \textbf{Condición Necesaria de Peres (1996):} Si $\hat{\rho}_{AB}$ es separable, entonces $\hat{\rho}_{AB}^{T_B} = \sum_k p_k \hat{\rho}_A^{(k)} \otimes (\hat{\rho}_B^{(k)})^T \geq 0$ es necesariamente un operador semidefinido positivo (todos sus autovalores son no negativos). Por contraposición, si $\hat{\rho}_{AB}^{T_B}$ posee al menos un **autovalor negativo**, entonces $\hat{\rho}_{AB}$ está garantizadamente **entrelazado**.
    \item \textbf{Teorema de los Horodecki (1996):} Para sistemas de dimensión baja $\mathbb{C}^2 \otimes \mathbb{C}^2$ (dos qubits) y $\mathbb{C}^2 \otimes \mathbb{C}^3$ (qubit-qutrit), la condición PPT ($\hat{\rho}_{AB}^{T_B} \geq 0$) es **necesaria y suficiente** para la separabilidad.
\end{itemize}
\end{teorema}

% =========================================================================
\section[Teorema de Bell y desigualdad CHSH]{Teorema de Bell y Desigualdad CHSH}

En 1964, John Stewart Bell transformó el debate metafísico entre Einstein y Bohr en un teorema matemático verificable experimentalmente, demostrando que **ninguna teoría de variables ocultas locales puede reproducir todas las predicciones estadísticas de la mecánica cuántica**. Presentaremos la generalización de Clauser, Horne, Shimony y Holt (CHSH, 1969), directamente aplicable a experimentos ópticos reales.

\begin{teorema}[Demostración Completa de la Desigualdad de Bell-CHSH bajo Realismo Local]
Consideremos una fuente que emite pares de partículas hacia Alice y Bob. Alice puede elegir libremente medir a lo largo de una de dos direcciones unitarias $\mathbf{a}$ o $\mathbf{a}'$, obteniendo resultados dicotómicos $A(\mathbf{a}, \lambda) \in \{+1, -1\}$ o $A(\mathbf{a}', \lambda) \in \{+1, -1\}$. Análogamente, Bob elige entre dos direcciones $\mathbf{b}$ o $\mathbf{b}'$, obteniendo resultados $B(\mathbf{b}, \lambda) \in \{+1, -1\}$ o $B(\mathbf{b}', \lambda) \in \{+1, -1\}$.

Bajo la hipótesis de **Realismo Local**:
\begin{enumerate}
    \item \textbf{Realismo (Variables Ocultas $\lambda \in \Lambda$):} Los valores $A$ y $B$ están determinados (o sus probabilidades factorizan) por un conjunto de variables ocultas $\lambda$ distribuidas según una densidad de probabilidad normalizada $\rho(\lambda)$:
    \begin{equation}
    \rho(\lambda) \geq 0, \qquad \int_{\Lambda} \rho(\lambda) \, d\lambda = 1
    \end{equation}
    \item \textbf{Localidad de Bell:} El resultado de Alice $A(\mathbf{a}, \lambda)$ no depende de la orientación $\mathbf{b}$ del analizador de Bob, ni $B(\mathbf{b}, \lambda)$ depende de $\mathbf{a}$. Asimismo, $\rho(\lambda)$ no depende de las elecciones de medición $(\mathbf{a}, \mathbf{b})$ (\textit{libertad de elección} o ausencia de superdeterminismo).
\end{enumerate}

La función de correlación estadística medida experimentalmente es el promedio sobre $\Lambda$:
\begin{equation}
E(\mathbf{a}, \mathbf{b}) = \int_{\Lambda} A(\mathbf{a}, \lambda) B(\mathbf{b}, \lambda) \rho(\lambda) \, d\lambda
\end{equation}

Para cada realización individual de la variable oculta $\lambda$, consideremos la combinación algebraica:
\begin{equation}
s(\lambda) \equiv A(\mathbf{a}, \lambda)B(\mathbf{b}, \lambda) - A(\mathbf{a}, \lambda)B(\mathbf{b}', \lambda) + A(\mathbf{a}', \lambda)B(\mathbf{b}, \lambda) + A(\mathbf{a}', \lambda)B(\mathbf{b}', \lambda)
\end{equation}
Factorizando $A(\mathbf{a}, \lambda)$ y $A(\mathbf{a}', \lambda)$:
\begin{equation}
s(\lambda) = A(\mathbf{a}, \lambda)\left[ B(\mathbf{b}, \lambda) - B(\mathbf{b}', \lambda) \right] + A(\mathbf{a}', \lambda)\left[ B(\mathbf{b}, \lambda) + B(\mathbf{b}', \lambda) \right]
\end{equation}
Como $B(\mathbf{b}, \lambda), B(\mathbf{b}', \lambda) \in \{+1, -1\}$, hay solo dos casos posibles para cada $\lambda$:
\begin{itemize}
    \item Si $B(\mathbf{b}, \lambda) = B(\mathbf{b}', \lambda)$, entonces $B(\mathbf{b}, \lambda) - B(\mathbf{b}', \lambda) = 0$ y $B(\mathbf{b}, \lambda) + B(\mathbf{b}', \lambda) = \pm 2$. Como $A(\mathbf{a}', \lambda) = \pm 1$, resulta $s(\lambda) = \pm 2$.
    \item Si $B(\mathbf{b}, \lambda) = -B(\mathbf{b}', \lambda)$, entonces $B(\mathbf{b}, \lambda) + B(\mathbf{b}', \lambda) = 0$ y $B(\mathbf{b}, \lambda) - B(\mathbf{b}', \lambda) = \pm 2$. Como $A(\mathbf{a}, \lambda) = \pm 1$, resulta $s(\lambda) = \pm 2$.
\end{itemize}
En todos los casos posibles, $|s(\lambda)| = 2$. Integrando respecto a la medida positiva $\rho(\lambda)$ y aplicando la desigualdad triangular $| \int f(\lambda)\rho(\lambda)d\lambda | \leq \int |f(\lambda)|\rho(\lambda)d\lambda$, obtenemos la **Desigualdad de Bell-CHSH**:
\begin{equation}
\boxed{ |S_{\text{CHSH}}| \equiv \left| E(\mathbf{a}, \mathbf{b}) - E(\mathbf{a}, \mathbf{b}') + E(\mathbf{a}', \mathbf{b}) + E(\mathbf{a}', \mathbf{b}') \right| \leq 2 }
\end{equation}
\end{teorema}

\subsection{Violación Cuántica de la Desigualdad CHSH y la Cota de Tsirelson}

Calculemos ahora la predicción de la Mecánica Cuántica para un par de partículas de espín $1/2$ en el estado singlete $|\Psi^-\rangle = \frac{1}{\sqrt{2}}(|\uparrow\downarrow\rangle - |\downarrow\uparrow\rangle)$. Los operadores de medición de Alice y Bob son $\hat{A}(\mathbf{a}) = \mathbf{a}\cdot\hat{\bm{\sigma}}^A$ y $\hat{B}(\mathbf{b}) = \mathbf{b}\cdot\hat{\bm{\sigma}}^B$.

\begin{ejemplo}[Cálculo del Correlador Cuántico y Violación Máxima a $2\sqrt{2}$]
Evaluemos el valor esperado cuántico $E_{\text{QM}}(\mathbf{a}, \mathbf{b}) = \langle \Psi^- | (\mathbf{a}\cdot\hat{\bm{\sigma}}^A) \otimes (\mathbf{b}\cdot\hat{\bm{\sigma}}^B) | \Psi^- \rangle$. Como $|\Psi^-\rangle$ es invariante bajo rotaciones globales, podemos orientar el eje $\hat{\mathbf{z}}$ paralelo al vector $\mathbf{a}$, de modo que $\mathbf{a}\cdot\hat{\bm{\sigma}}^A = \hat{\sigma}_z^A$ y $\mathbf{b}\cdot\hat{\bm{\sigma}}^B = b_x \hat{\sigma}_x^B + b_y \hat{\sigma}_y^B + b_z \hat{\sigma}_z^B$ con $b_z = \mathbf{a}\cdot\mathbf{b} = \cos\theta_{ab}$. Aplicando los operadores de Pauli sobre $|\Psi^-\rangle$:
\begin{align}
(\hat{\sigma}_z^A \otimes \hat{\sigma}_z^B) |\Psi^-\rangle &= \frac{1}{\sqrt{2}} \left( (+1)(-1)|\uparrow\downarrow\rangle - (-1)(+1)|\downarrow\uparrow\rangle \right) = -|\Psi^-\rangle \\
(\hat{\sigma}_z^A \otimes \hat{\sigma}_x^B) |\Psi^-\rangle &= \frac{1}{\sqrt{2}} \left( |\uparrow\uparrow\rangle + |\downarrow\downarrow\rangle \right) \quad (\text{ortogonal a } |\Psi^-\rangle) \\
(\hat{\sigma}_z^A \otimes \hat{\sigma}_y^B) |\Psi^-\rangle &= \frac{i}{\sqrt{2}} \left( |\uparrow\uparrow\rangle - |\downarrow\downarrow\rangle \right) \quad (\text{ortogonal a } |\Psi^-\rangle)
\end{align}
Por consiguiente, el correlador cuántico de espín es simplemente el producto escalar negativo:
\begin{equation}
E_{\text{QM}}(\mathbf{a}, \mathbf{b}) = -\mathbf{a} \cdot \mathbf{b} = -\cos\theta_{ab}
\end{equation}
Elijamos cuatro vectores coplanares con ángulos sucesivos separados por $45^\circ$ ($\pi/4$ rad):
\begin{equation}
\theta_{\mathbf{a}} = 0^\circ, \qquad \theta_{\mathbf{b}} = 45^\circ, \qquad \theta_{\mathbf{a}'} = 90^\circ, \qquad \theta_{\mathbf{b}'} = 135^\circ
\end{equation}
Los ángulos relativos son $\theta_{ab} = 45^\circ$, $\theta_{ab'} = 135^\circ$, $\theta_{a'b} = 45^\circ$ y $\theta_{a'b'} = 45^\circ$. Sustituyendo en el parámetro CHSH:
\begin{align}
|S_{\text{QM}}| &= \left| -\cos(45^\circ) - \left(-\cos(135^\circ)\right) + \left(-\cos(45^\circ)\right) + \left(-\cos(45^\circ)\right) \right| \\
&= \left| -\frac{\sqrt{2}}{2} - \frac{\sqrt{2}}{2} - \frac{\sqrt{2}}{2} - \frac{\sqrt{2}}{2} \right| = 4 \times \frac{\sqrt{2}}{2} = 2\sqrt{2} \approx 2.8284 > 2
\end{align}
¡La mecánica cuántica viola el límite clásico de Bell ($|S| \leq 2$) alcanzando $2\sqrt{2}$!
\end{ejemplo}

\begin{teorema}[Cota Cuántica de Tsirelson (1980)]
Para cualesquiera operadores hermíticos dicotómicos $\hat{A}_0, \hat{A}_1$ de Alice y $\hat{B}_0, \hat{B}_1$ de Bob tales que $\hat{A}_i^2 = \hat{\mathbb{I}}$, $\hat{B}_j^2 = \hat{\mathbb{I}}$ y $[\hat{A}_i, \hat{B}_j] = 0$, el operador de Bell-CHSH $\hat{\mathcal{B}} = \hat{A}_0 \hat{B}_0 - \hat{A}_0 \hat{B}_1 + \hat{A}_1 \hat{B}_0 + \hat{A}_1 \hat{B}_1$ satisface la identidad algebraica de Landau-Tsirelson:
\begin{equation}
\hat{\mathcal{B}}^2 = 4\hat{\mathbb{I}} + [\hat{A}_0, \hat{A}_1][\hat{B}_0, \hat{B}_1]
\end{equation}
Dado que la norma de operador de cualquier conmutador de operadores unitarios hermíticos cumple $\|[\hat{A}_0, \hat{A}_1]\| \leq 2\|\hat{A}_0\|\|\hat{A}_1\| = 2$, se sigue que $\|\hat{\mathcal{B}}^2\| \leq 4 + 2 \times 2 = 8$, lo que demuestra que $2\sqrt{2}$ es el límite máximo absoluto permitido por la mecánica cuántica:
\begin{equation}
|S_{\text{QM}}| \leq 2\sqrt{2} \quad (\text{Cota de Tsirelson})
\end{equation}
\end{teorema}

\subsection[Implementación óptica y cierre de resquicios]{Implementación Óptica con Fotones Entrelazados y Cierre de \textit{Loopholes}}

En óptica cuántica, la base de espín se codifica en la polarización horizontal $|H\rangle$ y vertical $|V\rangle$ de pares de fotones generados mediante cascadas radiativas atómicas (Clauser y Freedman, 1972; Aspect, Grangier y Roger, 1981-1982) o mediante **Conversión Paramétrica Espontánea Descendente (SPDC)** tipo-II en cristales birrefringentes no lineales como el BBO ($\beta$-$\text{BaB}_2\text{O}_4$, Kwiat \textit{et al.}, 1995; Zeilinger \textit{et al.}).

Dado que una rotación física por un ángulo $\theta$ de un polarizador rota el estado en la esfera de Poincaré un ángulo doble $2\theta$ ($|\theta\rangle = \cos\theta|H\rangle + \sin\theta|V\rangle$, de modo que $|V\rangle$ es ortogonal a $|H\rangle$ a $90^\circ$ en lugar de $180^\circ$), el correlador óptico de polarización para el estado $|\Phi^+\rangle = \frac{1}{\sqrt{2}}(|HH\rangle + |VV\rangle)$ es:
\begin{equation}
E_{\text{óptico}}(\theta_A, \theta_B) = P_{++}(\theta_A, \theta_B) + P_{--}(\theta_A, \theta_B) - P_{+-}(\theta_A, \theta_B) - P_{-+}(\theta_A, \theta_B) = \cos\left[2(\theta_A - \theta_B)\right]
\end{equation}
La violación máxima $|S_{\text{CHSH}}| = 2\sqrt{2}$ se alcanza en óptica orientando los ejes de los polarizadores divisores de haz a los ángulos mitad:
\begin{equation}
\theta_{\mathbf{a}} = 0^\circ, \qquad \theta_{\mathbf{b}} = 22.5^\circ, \qquad \theta_{\mathbf{a}'} = 45^\circ, \qquad \theta_{\mathbf{b}'} = 67.5^\circ
\end{equation}

Históricamente, para descartar definitivamente el realismo local fue necesario cerrar simultáneamente tres resquicios (\textit{loopholes}) experimentales, hazaña lograda en **2015** (Hensen \textit{et al.} en Delft; Giustina \textit{et al.} en Viena; Shalm \textit{et al.} en NIST, motivo del **Premio Nobel de Física 2022** a Alain Aspect, John F. Clauser y Anton Zeilinger):
\begin{enumerate}
    \item \textbf{Loophole de Localidad o Comunicación:} Los cambios de configuración de los analizadores $(\mathbf{a}, \mathbf{a}')$ y $(\mathbf{b}, \mathbf{b}')$ deben realizarse mediante conmutadores electro-ópticos ultrarrápidos mientras los fotones están en vuelo, de modo que la elección de Alice y el registro de detección de Bob estén separados por un intervalo relativista de tipo espacio (iniciado por Aspect en 1982 y perfeccionado por Weihs \textit{et al.} en 1998).
    \item \textbf{Loophole de Detección o Muestreo Justo (\textit{Fair Sampling}):} Si la eficiencia cuántica $\eta$ de los fotodetectores es baja, las variables ocultas podrían instruir a los fotones para no ser detectados cuando el ángulo no les convenga. Para violar la desigualdad de Eberhard sin asumir muestreo justo, la eficiencia total del sistema debe superar el umbral crítico $\eta > \frac{2}{\sqrt{2}+1} \approx 82.8\%$ (para estados máximamente entrelazados) o $\eta > 66.7\%$ (para estados parcialmente entrelazados), logrado mediante detectores superconductores de nanocables de un solo fotón (SNSPD) con $\eta > 90\%$.
    \item \textbf{Loophole de Libertad de Elección:} Uso de generadores cuánticos de números aleatorios ultrarrápidos e incluso fluctuaciones de fotones provenientes de cuásares situados en extremos opuestos del universo a miles de millones de años luz para fijar los ángulos $\theta_A$ y $\theta_B$.
\end{enumerate}

\section[Teorema GHZ y no-localidad]{No-Localidad y Teorema GHZ}

Mientras que el teorema de Bell requiere promedios estadísticos sobre muchas mediciones para demostrar la contradicción con el realismo local ($2\sqrt{2} > 2$), Daniel Greenberger, Michael Horne y Anton Zeilinger demostraron en 1989 que con **tres o más partículas entrelazadas** la contradicción con las variables ocultas locales es del $100\%$ en **una sola medición determinista**.

Consideremos tres qubits (fotones $A, B, C$) preparados en el **estado GHZ tripartito**:
\begin{equation}
|\text{GHZ}\rangle = \frac{1}{\sqrt{2}} \left( |000\rangle + |111\rangle \right)
\end{equation}
Definamos los tres operadores hermíticos conmutativos de tres cuerpos:
\begin{equation}
\hat{\mathcal{O}}_1 = \hat{\sigma}_x^A \otimes \hat{\sigma}_y^B \otimes \hat{\sigma}_y^C, \qquad \hat{\mathcal{O}}_2 = \hat{\sigma}_y^A \otimes \hat{\sigma}_x^B \otimes \hat{\sigma}_y^C, \qquad \hat{\mathcal{O}}_3 = \hat{\sigma}_y^A \otimes \hat{\sigma}_y^B \otimes \hat{\sigma}_x^C
\end{equation}
Dado que $\hat{\sigma}_x|0\rangle = |1\rangle, \hat{\sigma}_x|1\rangle = |0\rangle$ y $\hat{\sigma}_y|0\rangle = i|1\rangle, \hat{\sigma}_y|1\rangle = -i|0\rangle$, es inmediato verificar que $|\text{GHZ}\rangle$ es autoestado simultáneo de $\hat{\mathcal{O}}_1, \hat{\mathcal{O}}_2$ y $\hat{\mathcal{O}}_3$ con autovalor **$-1$**:
\begin{equation}
\hat{\mathcal{O}}_1 |\text{GHZ}\rangle = -|\text{GHZ}\rangle, \qquad \hat{\mathcal{O}}_2 |\text{GHZ}\rangle = -|\text{GHZ}\rangle, \qquad \hat{\mathcal{O}}_3 |\text{GHZ}\rangle = -|\text{GHZ}\rangle
\end{equation}
En cambio, para el operador producto de las tres componentes $x$, $\hat{\mathcal{O}}_4 = \hat{\sigma}_x^A \otimes \hat{\sigma}_x^B \otimes \hat{\sigma}_x^C$, tenemos un autovalor **$+1$**:
\begin{equation}
\hat{\mathcal{O}}_4 |\text{GHZ}\rangle = \left(\hat{\sigma}_x^A \otimes \hat{\sigma}_x^B \otimes \hat{\sigma}_x^C\right) |\text{GHZ}\rangle = +1 |\text{GHZ}\rangle
\end{equation}

Supongamos ahora que existen elementos de realidad preexistentes locales $s_x^i, s_y^i \in \{+1, -1\}$ ($i \in \{A, B, C\}$). Para reproducir las tres primeras predicciones cuánticas con certeza, deben cumplirse las tres ecuaciones:
\[
s_x^A s_y^B s_y^C = -1, \qquad s_y^A s_x^B s_y^C = -1, \qquad s_y^A s_y^B s_x^C = -1
\]
Multiplicando las tres ecuaciones miembro a miembro y recordando que $(s_y^i)^2 = (\pm 1)^2 = +1$:
\begin{equation}
\left( s_x^A s_x^B s_x^C \right) \underbrace{(s_y^A)^2 (s_y^B)^2 (s_y^C)^2}_{=1} = (-1)(-1)(-1) = -1 \implies \left( s_x^A s_x^B s_x^C \right)_{\text{Realismo Local}} = -1
\end{equation}
¡Mientras que el realismo local exige que el producto de las mediciones en $x$ sea siempre **$-1$**, la mecánica cuántica predice y los experimentos confirman que el resultado es siempre **$+1$** ($+1 = -1$, contradicción absoluta)!

% =========================================================================
% BIBLIOGRAFÍA Y LITERATURA COMPLEMENTARIA
% =========================================================================
\begin{thebibliography}{99}

\bibitem{Dirac1930}
P. A. M. Dirac, \textit{The Principles of Quantum Mechanics}, 4th ed., Oxford University Press, Oxford (1958) [1st ed. 1930].

\bibitem{VonNeumann1932}
J. von Neumann, \textit{Mathematische Grundlagen der Quantenmechanik}, Springer, Berlin (1932) [Traducción al español: \textit{Fundamentos matemáticos de la mecánica cuántica}, CSIC, Madrid].

\bibitem{CohenTannoudji}
C. Cohen-Tannoudji, B. Diu, y F. Laloë, \textit{Quantum Mechanics} (Vols. I, II y III), Wiley-VCH (2019).

\bibitem{Sakurai}
J. J. Sakurai y J. Napolitano, \textit{Modern Quantum Mechanics}, 3rd ed., Cambridge University Press (2020).

\bibitem{Shankar}
R. Shankar, \textit{Principles of Quantum Mechanics}, 2nd ed., Springer (1994).

\bibitem{Ballentine}
L. E. Ballentine, \textit{Quantum Mechanics: A Modern Development}, 2nd ed., World Scientific (2014).

\bibitem{Glauber1963}
R. J. Glauber, ``The Quantum Theory of Optical Coherence'', \textit{Physical Review} \textbf{130}, 2529--2539 (1963); ``Coherent and Incoherent States of the Radiation Field'', \textit{Physical Review} \textbf{131}, 2766--2788 (1963).

\bibitem{ScullyZubairy}
M. O. Scully y M. S. Zubairy, \textit{Quantum Optics}, Cambridge University Press (1997).

\bibitem{GerryKnight}
C. C. Gerry y P. L. Knight, \textit{Introductory Quantum Optics}, 2nd ed., Cambridge University Press (2023).

\bibitem{MandelWolf}
L. Mandel y E. Wolf, \textit{Optical Coherence and Quantum Optics}, Cambridge University Press (1995).

\bibitem{WallsMilburn}
D. F. Walls y G. J. Milburn, \textit{Quantum Optics}, 2nd ed., Springer (2008).

\bibitem{Loudon}
R. Loudon, \textit{The Quantum Theory of Light}, 3rd ed., Oxford University Press (2000).

\bibitem{Fox}
M. Fox, \textit{Quantum Optics: An Introduction}, Oxford University Press (2006).

\bibitem{GrynbergAspectFabre}
G. Grynberg, A. Aspect, y C. Fabre, \textit{Introduction to Quantum Optics: From the Semi-classical Approach to Quantized Light}, Cambridge University Press (2010).

\bibitem{HarocheRaimond}
S. Haroche y J.-M. Raimond, \textit{Exploring the Quantum: Atoms, Cavities, and Photons}, Oxford University Press (2006).

\bibitem{JaynesCummings1963}
E. T. Jaynes y F. W. Cummings, ``Comparison of quantum and semiclassical radiation theories with application to the beam maser'', \textit{Proceedings of the IEEE} \textbf{51}, 89--109 (1963).

\bibitem{HOM1987}
C. K. Hong, Z. Y. Ou, y L. Mandel, ``Measurement of subpicosecond time intervals between two photons by interference'', \textit{Physical Review Letters} \textbf{59}, 2044--2046 (1987).

\bibitem{EPR1935}
A. Einstein, B. Podolsky, y N. Rosen, ``Can Quantum-Mechanical Description of Physical Reality Be Considered Complete?'', \textit{Physical Review} \textbf{47}, 777--780 (1935).

\bibitem{Bell1964}
J. S. Bell, ``On the Einstein Podolsky Rosen Paradox'', \textit{Physics Physique Fizika} \textbf{1}, 195--200 (1964); recopilado en J. S. Bell, \textit{Speakable and Unspeakable in Quantum Mechanics}, 2nd ed., Cambridge University Press (2004).

\bibitem{CHSH1969}
J. F. Clauser, M. A. Horne, A. Shimony, y R. A. Holt, ``Proposed Experiment to Test Local Hidden-Variable Theories'', \textit{Physical Review Letters} \textbf{23}, 880--884 (1969).

\bibitem{Aspect1982}
A. Aspect, J. Dalibard, y G. Roger, ``Experimental Test of Bell's Inequalities Using Time-Varying Analyzers'', \textit{Physical Review Letters} \textbf{49}, 1804--1807 (1982).

\bibitem{Brunner2014}
N. Brunner, D. Cavalcanti, S. Pironio, V. Scarani, y S. Wehner, ``Bell nonlocality'', \textit{Reviews of Modern Physics} \textbf{86}, 419--478 (2014).

\bibitem{NielsenChuang}
M. A. Nielsen y I. L. Chuang, \textit{Quantum Computation and Quantum Information}, 10th Anniversary ed., Cambridge University Press (2010).

\bibitem{BreuerPetruccione}
H.-P. Breuer y F. Petruccione, \textit{The Theory of Open Quantum Systems}, Oxford University Press (2002).

\end{thebibliography}

% =========================================================================
\chapter*{Guía Comentada de Literatura Complementaria}
\addcontentsline{toc}{chapter}{Guía Comentada de Literatura Complementaria}

Para profundizar de manera estructurada según el nivel de rigor matemático o interés experimental, se recomienda el siguiente recorrido bibliográfico especializado:

\section*{1. Fundamentos Matemáticos y Espacios de Hilbert Rigurosos}
\begin{itemize}
    \item \textbf{Brian C. Hall, \textit{Quantum Theory for Mathematicians} (Springer, Graduate Texts in Mathematics, 2013):} La referencia definitiva si deseas profundizar en el teorema espectral para operadores autoadjuntos no acotados, dominios densos, el teorema de Stone sobre grupos monoparamétricos unitarios y la estructura rigurosa de la tripleta de Gelfand ($\Phi \subset \mathcal{H} \subset \Phi^\times$).
    \item \textbf{Asher Peres, \textit{Quantum Theory: Concepts and Methods} (Kluwer Academic, 1995):} Una obra maestra conceptual que conecta los postulados operativos de la mecánica cuántica con medidas generalizadas (POVMs), teoremas tipo Kochen-Specker (contextualidad cuántica) y geometría del espacio de estados.
    \item \textbf{Alberto Galindo y Pedro Pascual, \textit{Mecánica Cuántica} (Vols. I y II, Eudema / Springer, 1989):} El tratado más riguroso escrito originalmente en español, con un tratamiento matemático impecable de operadores en espacios de Hilbert, teoría de dispersión y simetrías.
\end{itemize}

\section*{2. Óptica Cuántica Experimental, Cuadraturas y Sistemas Abiertos}
\begin{itemize}
    \item \textbf{Ulrik L. Andersen, Tobias Gehring, Christoph Marquardt y Gerd Leuchs, \textit{30 years of squeezed light generation} (Physica Scripta 91, 053001, 2016):} Excelente artículo de revisión sobre cómo se generan experimentalmente estados comprimidos mediante osciladores paramétricos ópticos (OPO) y detección homodina equilibrada.
    \item \textbf{Howard J. Carmichael, \textit{Statistical Methods in Quantum Optics 1 \& 2} (Springer, 1999, 2008):} Indispensable para dominar las ecuaciones de Fokker-Planck asociadas a las representaciones $P$, $Q$ y Wigner, así como la teoría de trayectorias cuánticas (\textit{quantum jumps} y desenrollado de Monte Carlo de la ecuación de Lindblad).
    \item \textbf{Pierre Meystre y Murray Sargent III, \textit{Elements of Quantum Optics} (4th ed., Springer, 2007):} Muy recomendable para conectar las ecuaciones ópticas de Bloch semiclásicas con la teoría cuántica del láser de Scully-Lamb y el enfriamiento Doppler de átomos por láser.
\end{itemize}

\section*{3. No-Localidad de Bell, Entrelazamiento e Información Cuántica}
\begin{itemize}
    \item \textbf{Ryszard Horodecki, Paweł Horodecki, Michał Horodecki y Karol Horodecki, \textit{Quantum entanglement} (Reviews of Modern Physics 81, 865--942, 2009):} El artículo de revisión enciclopédico estándar sobre teoría del entrelazamiento, testigos de entrelazamiento (\textit{entanglement witnesses}), destilación, entrelazamiento ligado (\textit{bound entanglement}) y medidas monotónicas bajo operaciones locales y comunicación clásica (LOCC).
    \item \textbf{Valerio Scarani, \textit{Bell Nonlocality} (Oxford University Press, 2019):} Un libro moderno dedicado exclusivamente a los politopos de correlación local vs. cuántico vs. no-señalización (cajas de Popescu-Rohrlich), criptografía cuántica independiente del dispositivo (\textit{Device-Independent QKD}) y certificación de aleatoriedad cuántica.
    \item \textbf{John Preskill, \textit{Lecture Notes for Ph219/CS219: Quantum Information} (Caltech):} Apuntes legendarios de acceso libre que desarrollan con extraordinaria claridad los canales cuánticos (representación de suma de operadores de Kraus), el teorema de Stinespring, la corrección de errores cuánticos y la tolerancia a fallos.
\end{itemize}

\end{document}